\chapter{Cartier b-divisors and homogeneous realization}
\index{b-divisor@$b$-divisor!Cartier}\index{homogeneous realization}
\label{chap:bj-cartier-realization}

This chapter constructs the first independent ingredient in the envelope
proof.  An anti-effective real Cartier b-divisor with nef total class will
be realized as the divisorial trace of a homogeneous psh function.  The
construction uses ordered rational approximation, semiample coefficient
functions, the decreasing-limit axiom, comparison, and homogeneous Siu
division.  It does not use the envelope property.

Throughout, $X$ is a normal integral projective variety over the
algebraically closed trivially valued field $k$.  A model over $X$ is a
normal projective birational morphism.  Inequalities between divisors with
one fixed finite support are coefficientwise.

\section{Orientation}
\label{sec:ch38-orientation}

\subsection{What is built here, and for whom}
\label{sec:ch38-what}

\Cref{chap:bj-homogeneity-bdivisors} established a dictionary between three
descriptions at the Cartier level: a homogeneous PL function on
$X^{\val}$, its divisorial trace, and a Cartier b-divisor over $X$.
The coefficient dictionary is a bijection
(\Cref{prop:cartier-homogeneous-pl}), and the positivity condition
characterizing the Fubini--Study cone inside it is
\emph{semiampleness} of the total class $\pi^*L+B_Y$
(\Cref{thm:ch35-semiample-criterion}).  What is not yet available in
arbitrary characteristic is the corresponding statement with
\emph{nefness} in place of semiampleness.  That statement,
\Cref{thm:ch35-cartier-psh-criterion}, was proved in
\Cref{chap:bj-homogeneity-bdivisors} only under the hypothesis that ample
classes on $X$ have the envelope property, and its proof passes through
an increasing limit.  Since
\Crefrange{chap:bj-cartier-realization}{chap:bj-subintegral-normalization}
exist in order to \emph{prove} the envelope property, that route is closed
to us.

The present chapter supplies the substitute.  Its main theorem,
\Cref{thm:bj-real-cartier-realization}, says:

\begin{quote}
if $B\leq0$ is a \emph{real} Cartier b-divisor over a normal integral
projective $X$, determined on a model $\pi:Y\to X$, and if the total
class $\pi^*\theta+B_Y$ is nef, then the coefficient function
$\psi_B$ on $X^{\divis}$ is the divisorial trace of a unique
homogeneous $\theta$-psh function $\chi_B$ on $X^{\An}$.
\end{quote}

The gain over \Cref{thm:ch35-cartier-psh-criterion} is threefold.  The
coefficients are real, not rational; no smoothness or characteristic
hypothesis is imposed; and, decisively, no envelope property is used.  The
price is that the conclusion is stated for \emph{Cartier} b-divisors only.
An arbitrary Weil b-divisor is not realized here, and
\Cref{rem:ch38-cartier-not-weil} explains why the Cartier case is
nevertheless exactly what the later chapters consume.

Three later interfaces motivate the detailed form of the result.
\Cref{chap:bj-reflection} uses the realization theorem, its order
preservation, and the fixed-pole data of \Cref{lem:ch38-activation}.
\Crefrange{chap:bj-smooth-characteristic}{chap:bj-subintegral-normalization}
use the decreasing model approximation of
\Cref{cor:ch38-gf-regularization-bdivisor}.  Finally,
\Cref{chap:bj-envelope-consequences} uses
\Cref{cor:ch38-cartier-psh-criterion-unconditional}, which upgrades the
Cartier psh criterion to real coefficients without an envelope, smoothness,
or characteristic hypothesis.

\subsection{The firewall}
\label{sec:ch38-firewall}

\Cref{chap:bj-envelope-history} warned, in the discussion following
\Cref{def:bj-obstacle-envelope}, that before the envelope theorem the
operator $P_\theta$ is a map from obstacles to functions and not to
$\PSH_X(\theta)$.  The same warning applies here in a stronger form,
because this chapter is one of the inputs of
\Cref{thm:bj-final-envelope} and \Cref{chap:bj-envelope-statement}
cites it in the proof of that theorem.  We therefore state the firewall
explicitly and once.

\begin{warning}{The envelope theorem is not available in this chapter}
\label{warn:ch38-firewall}
Nothing in \Cref{chap:bj-cartier-realization} may use, directly or
transitively:
\begin{enumerate}
\item \Cref{thm:bj-final-envelope}, or the envelope property
 $\operatorname{EP}$ of \Cref{def:bj-envelope-property} for any class on
 any variety;
\item closure of $\PSH_X(\theta)$ under suprema, regularized or not, and
 in particular any appeal to an increasing net of psh functions having a
 psh limit;
\item the obstacle envelope $P_\theta$ regarded as an operator with
 values in $\PSH_X(\theta)$;
\item \Cref{thm:ch35-cartier-psh-criterion} in its nontrivial direction
 $\Leftarrow$, whose proof in
 \Cref{chap:bj-homogeneity-bdivisors} is an increasing limit governed by
 $\operatorname{EP}$ for ample classes;
\item \Cref{thm:bj-nef-bdivisor-correspondence} and
 \Cref{cor:ch35-monotone-approximation}, which rest on the previous item
 and on smoothness in characteristic zero;
\item \Cref{thm:ch34-psh-compactness} in the homogeneous form used in the
 proof of \Cref{thm:bj-nef-bdivisor-correspondence}.
\end{enumerate}
The only limit operation performed below is a decreasing one.  Every
appearance of the word ``limit'' in this chapter refers to a decreasing
net or to a convergent net in a finite-dimensional real vector space.
\end{warning}

It is worth saying what \emph{is} available, since the list is short and
each item is used.  From \Cref{chap:bj-fs-pl-psh}: the definition
\Cref{def:bj-psh} of $\PSH_X(\theta)$, the closure calculus
\Cref{thm:bj-basic-closure}, the existence statement
\Cref{lem:ch34-psh-nonempty}, the hierarchy \eqref{eq:ch34-hierarchy}, and
the Fubini--Study regularization \Cref{thm:ch34-fs-regularization}.  From
\Cref{chap:bj-homogeneity-bdivisors}: the evaluation pairing
\Cref{prop:ch35-weil-coefficient-bijection}, its Cartier refinement
\Cref{prop:ch35-cartier-continuous-extension}, the identification
\Cref{prop:cartier-homogeneous-pl} of Cartier b-divisors with homogeneous
PL functions, the semiample criterion
\Cref{thm:ch35-semiample-criterion}, the normalization
\Cref{prop:ch35-homogeneous-normalization}, the comparison theorem
\Cref{thm:ch35-divisorial-comparison} with its corollary
\Cref{cor:ch35-divisorial-determination}, the homogenization corollary
\Cref{cor:ch35-homogenization-psh}, and the homogeneous Siu decomposition
\Cref{thm:ch35-homogeneous-siu} together with its general form
\Cref{thm:ch35-homogeneous-siu-general}.  Following the usage of
\Crefrange{chap:bj-reflection}{chap:bj-envelope-consequences}, we
abbreviate as \emph{PSH2} the decreasing-limit axiom, that is
\Cref{thm:bj-basic-closure}\,(5): if $\psi_j\in\PSH_X(\theta_j)$ with
$\theta_j\to\theta$, the net $(\psi_j)$ is decreasing and the limit is
generically finite, then the limit is $\theta$-psh.  Part~(4), the
fixed-class case, is its specialization.  It is the psh limit operation
used below; the nontriviality condition on the limit is
checked explicitly at every application.  We nevertheless apply it in the
fixed-class form throughout, for the reason given in
\Cref{warn:ch38-drift-not-absorbed}.  These available inputs are unconditional.
The Cartier psh criterion has an explicit envelope hypothesis, and the
numerical b-class correspondence also uses it in its smooth
characteristic-zero setting. Both are excluded by the firewall above.

\begin{remark}[Why the firewall is not vacuous]
\label{rem:ch38-firewall-not-vacuous}
A reader who has met the theory in its finished form will find the
statement of \Cref{thm:bj-real-cartier-realization} indistinguishable from
the implication $\Leftarrow$ of \Cref{thm:ch35-cartier-psh-criterion}
with real coefficients, and may wonder what there is to prove.  The answer
is contained in the proof of \Cref{thm:ch35-cartier-psh-criterion}: there
one perturbs $B$ by $\varepsilon H$ with $H$ relatively ample and
\emph{negative}, obtains $\psi_B+\varepsilon\psi_H\in\cH_{\mathrm{hom}}(L)$
for each $\varepsilon>0$, and lets $\varepsilon\downarrow0$.  Because
$\psi_H\leq0$, the resulting family \emph{increases} to $\psi_B$, and
an increasing limit of psh functions is psh only by the envelope property.
The whole point of the construction below is to arrange an approximating
family that decreases.  This is why the activated divisor $B'$ is approached from
above by $-R_n$ with $R_n$ increasing, and why the auxiliary positivity
is inserted as an \emph{ample buffer} whose effect vanishes in the limit
rather than as a negative perturbation of the b-divisor.
\end{remark}

\subsection{The five tasks, and where each is discharged}
\label{sec:ch38-payoff}

\begin{enumerate}
\item\label{itm:ch38-task:evaluation}
 \emph{Evaluation of a Cartier b-divisor on valuations.}
 \Cref{sec:ch38-evaluation} defines
 $\psi_B$ on $X^{\divis}$, proves that the definition is independent
 of every choice
 (\Cref{prop:ch38-evaluation-welldefined}), that it is functorial for
 birational pullback (\Cref{prop:ch38-evaluation-functorial}), and that it
 extends continuously to the full-support locus but \emph{not} by a
 formula to all of $X^{\An}$
 (\Cref{prop:ch38-evaluation-domains},
 \Cref{ex:ch38-evaluation-no-lift}).
\item\label{itm:ch38-task:fs}
 \emph{Model and homogeneous Fubini--Study functions.}
 \Cref{sec:ch38-model-fs} isolates the exact dictionary between rational
 antieffective Cartier b-divisors with semiample total class and
 homogeneous Fubini--Study functions
 (\Cref{prop:ch38-model-fs-dictionary}), separates the two gradings that
 occur in the underlying Rees and base-ideal constructions
 (\Cref{prop:ch38-two-gradings}), and computes the correspondence in both
 directions on an explicit surface
 (\Cref{ex:ch38-fs-from-bdivisor-and-back}).
\item\label{itm:ch38-task:realization}
 \emph{Realization and approximation.}
 \Cref{sec:ch38-realization} proves
 \Cref{thm:bj-real-cartier-realization}, and
 \Cref{sec:ch38-consequences} packages the approximation statements
 \Cref{cor:ch38-gf-regularization-bdivisor} and
 \Cref{cor:ch38-realization-pullback} that later chapters consume.
\item\label{itm:ch38-task:classes}
 \emph{Passage between numerical classes on birational models.}
 \Cref{sec:ch38-classes} treats $\pi^*$, $\pi_*$, and the total class
 $P=\pi^*\theta+B_Y$; it proves the drift ledger
 \Cref{prop:ch38-drift-ledger}, which tracks the $O(m^{-1})$ error in
 numerical class through the construction and shows that what makes the
 error harmless is not its size but its \emph{sign}
 (\Cref{warn:ch38-drift-not-absorbed}).
\item\label{itm:ch38-task:extended}
 \emph{Extended-real safeguards.}  \Cref{sec:ch38-safeguards} lists every
 operation in the chapter that could produce the forbidden sum
 $(+\infty)+(-\infty)$, proves that none of them does
 (\Cref{prop:ch38-safeguards}), and exhibits a naive realization that
 fails at a proper-support point
 (\Cref{ex:ch38-naive-realization-fails}).
\end{enumerate}

Two negative statements accompany the list.
\Cref{ex:ch38-usc-not-psh} shows that the coefficient function of an
antieffective Cartier b-divisor with \emph{non-nef} total class is a
perfectly good upper-semicontinuous homogeneous function which is not psh
in the specified class of degree smaller than its pole coefficient; so \Cref{thm:bj-real-cartier-realization} is not a
regularity statement about extending functions, and the nefness hypothesis
cannot be dropped.  And \Cref{ex:ch38-strict-room-fails} shows that the
Siu division at the end of the proof genuinely occurs at the boundary of
the strict regime of \Cref{thm:ch35-homogeneous-siu}, so that the general
form \Cref{thm:ch35-homogeneous-siu-general} is the correct citation.

\subsection{Standing hypotheses and notation}
\label{sec:ch38-hypotheses}

Throughout, $k$ is algebraically closed and carries the trivial absolute
value, and
\[
 X\ \text{is a normal integral projective $k$-variety},\qquad
 n:=\dim X .
\tag{38.1}\label{eq:ch38-standing}
\]
Normality is in force from the first line, because the model category of
\Cref{sec:appF-models} and the identification
$\pi_*\cO_Y=\cO_X$ are used everywhere.  Numerical classes are real,
$N^1(X):=N^1(X)_{\mathbf R}$, and $\Pic(X)_{\mathbf Q}
=\Pic(X)\otimes_{\mathbf Z}\mathbf Q$.  A \emph{model} of $X$ is a
normal projective variety $Y$ together with a projective birational
morphism $\pi:Y\to X$; models form a cofiltered category under
domination, by \Cref{lem:appF-common-domination}, and no resolution of
singularities is used.

We write $X^{\An}$ for the Berkovich analytification,
$X^{\val}$ for the full-support locus
\eqref{eq:triv-full-support-locus}, $X^{\divis}\subset X^{\val}$ for the
divisorial locus, and $v_{\triv}$ for the trivial valuation of
$k(X)$, which lies in $X^{\divis}$ by convention.  A point of
$X^{\An}\setminus X^{\val}$ has \emph{proper support}
(\Cref{def:ch32-proper-support}).  For an $\mathbf R$-Cartier divisor
$R$ on a model $Y$ we write $\overline R$ for the Cartier b-divisor
determined by the pair $(Y,R)$; when $Y=X$ this is the notation of
\Cref{sec:ch35-bdivisors}, and the general case is the same construction
after \Cref{def:appF-cartier-bdivisor}.  Ordering of b-divisors is
coefficientwise, and for a Cartier b-divisor determined on $Y$ it
suffices to test $B_Y\leq0$.

Inequalities between extended-real functions are always written in the
relative form $\phi\leq\psi+C$ and never as $\phi-\psi\leq C$,
following ledger entry~(4) of \Cref{sec:extended-arithmetic} and
\Cref{rem:ledger-uses}.  The only extended value that occurs is
$-\infty$; $+\infty$ occurs only as the value $v(f)=+\infty$ of a
semivaluation on a section vanishing along its support, and the
translation \eqref{eq:bj-divisor-weight-sign} between the two sign
conventions is fixed in \Cref{sec:ch38-signs} once and for all.

\subsection{Dependencies}
\label{sec:ch38-dependencies}

The chapter uses \Cref{chap:trivial-analytification} for the
analytification, the scaling action, divisorial and quasimonomial points,
and the density theorem \Cref{thm:triv-divisorial-density};
\Cref{chap:bj-fs-pl-psh} for Fubini--Study, PL, and psh functions;
\Cref{chap:bj-homogeneity-bdivisors} for everything homogeneous;
\Cref{app:birational-geometry} for normalized blowups, relative
ampleness, and the domination theorem
\Cref{thm:appB-admissible-domination};
\Cref{app:intersection-bdivisors} for the b-divisor formalism and the
projection formula; and \Cref{sec:extended-arithmetic} for the
extended-real ledger. It uses no later envelope theorem or
reflection result as an input.

\section{Evaluation of a Cartier b-divisor on valuations}
\label{sec:ch38-evaluation}

\subsection{Cartier b-divisors and determinations}
\label{sec:ch38-cartier-bdivisors}

\begin{definition}[Real Cartier b-divisor]
\label{def:bj-real-cartier-bdivisor}
A \emph{real Cartier b-divisor} $B$ over $X$ is represented by
\[
                     (\pi:Y\longrightarrow X,B_Y),
\]
where $\pi$ is a model and $B_Y$ is an $\mathbf R$-Cartier divisor
on $Y$.  Two pairs are equivalent if their pullbacks agree on a common
dominating model.  The pair is a \emph{determination} of $B$.
\end{definition}

If $B_Y\le0$ on one determination, then $B\le0$.  Pullback preserves
this inequality on every higher model.

\Cref{def:bj-real-cartier-bdivisor} agrees with
\Cref{def:appF-cartier-bdivisor}, and the two clauses of the latter --- a
compatible family in the sense of \eqref{eq:appF-weil-compatibility}, which
is a pullback of a single divisor on all higher models --- are recovered
below.  We spell the elementary points out, because the whole chapter rests
on them and because the difference between a Cartier and a Weil b-divisor
is exactly the difference between a direct and an inverse limit in
\eqref{eq:appF-direct-inverse}.

\begin{lemma}
\label{lem:ch38-determination-equivalence}
Let $\mathcal P$ be the class of pairs $(\pi:Y\to X,\,B_Y)$ with $Y$
a model and $B_Y\in\Div(Y)_{\mathbf R}$ Cartier, and let
$(Y,B_Y)\sim(Y',B_{Y'})$ mean that there is a model $Z$ dominating both,
say $\mu:Z\to Y$ and $\mu':Z\to Y'$, with
$\mu^*B_Y=\mu'^*B_{Y'}$.  Then:
\begin{enumerate}
\item\label{itm:ch38-de:equiv} $\sim$ is an equivalence relation on
 $\mathcal P$, and if $\mu^*B_Y=\mu'^*B_{Y'}$ on one common domination
 then the same holds on every common domination;
\item\label{itm:ch38-de:family} each equivalence class determines a unique
 Weil b-divisor $B=(B_{Y''})_{Y''}$ with
 $B_{Y''}=\nu^*B_Y$ whenever $\nu:Y''\to Y$ is a domination of a
 determination, and $B_{Y''}=\nu'_*(\nu^*B_Y)$ in general, where
 $Z\xrightarrow{\nu}Y$ and $Z\xrightarrow{\nu'}Y''$ is any common
 domination;
\item\label{itm:ch38-de:order} $B\leq0$ if and only if $B_Y\leq0$ for
 one, equivalently every, determination.
\end{enumerate}
\end{lemma}

\begin{proof}
\emph{\cref{itm:ch38-de:equiv}.}  Reflexivity and symmetry are immediate.
For transitivity, let $(Y_1,B_1)\sim(Y_2,B_2)$ via $Z_{12}$ and
$(Y_2,B_2)\sim(Y_3,B_3)$ via $Z_{23}$.  By
\Cref{lem:appF-common-domination} choose a model $W$ dominating
$Z_{12}$ and $Z_{23}$; pulling back the two equalities to $W$ and
using functoriality of pullback for Cartier divisors along the composites
gives $\alpha_1^*B_1=\alpha_2^*B_2=\alpha_3^*B_3$ on $W$, where
$\alpha_i:W\to Y_i$.  Hence $(Y_1,B_1)\sim(Y_3,B_3)$.

For the last clause of \cref{itm:ch38-de:equiv}, suppose
$\mu^*B_Y=\mu'^*B_{Y'}$ on $Z$, and let $\widetilde Z$ be another
common domination, with maps $\widetilde\mu,\widetilde\mu'$.  Choose
$W$ dominating $Z$ and $\widetilde Z$, with
$\beta:W\to Z$, $\widetilde\beta:W\to\widetilde Z$.  Then
$\beta^*\mu^*B_Y=\beta^*\mu'^*B_{Y'}$, that is,
$\widetilde\beta^*\widetilde\mu^*B_Y=\widetilde\beta^*\widetilde\mu'^*
B_{Y'}$, because the two composites $W\to Y$ agree, both being the
canonical birational morphism over $X$, and similarly for $Y'$.  As
$\widetilde\beta$ is birational, $\widetilde\beta^*$ is injective on
$\mathbf R$-Cartier divisors: a Cartier divisor is determined by its
local equations, and $\widetilde\beta^\#$ is injective on function
fields, both being $k(X)$.  Hence
$\widetilde\mu^*B_Y=\widetilde\mu'^*B_{Y'}$.

\emph{\cref{itm:ch38-de:family}.}  Fix a determination $(Y,B_Y)$.  For an
arbitrary model $Y''$, choose a common domination
$Z\xrightarrow{\nu}Y$, $Z\xrightarrow{\nu'}Y''$ and set
$B_{Y''}:=\nu'_*\nu^*B_Y$.  Independence of the choice of $Z$ follows
from \cref{itm:ch38-de:equiv} together with the projection formula
\eqref{eq:appF-projection}: if $W$ dominates two such $Z$, pushing
forward from $W$ factors through either.  Compatibility
\eqref{eq:appF-weil-compatibility} under a further domination
$\lambda:Y_1''\to Y''$ is
$\lambda_*B_{Y_1''}=\lambda_*\lambda'_*\nu_1^*B_Y=\nu'_*\nu^*B_Y=B_{Y''}$,
again by functoriality of pushforward.  If $Y''\succeq Y$, one may take
$Z=Y''$ and $\nu'=\operatorname{id}$, giving
$B_{Y''}=\nu^*B_Y$.

\emph{\cref{itm:ch38-de:order}.}  If $B_Y\leq0$, then
$\nu^*B_Y\leq0$ for every domination $\nu$, because the pullback of an
effective Cartier divisor is effective; and pushforward preserves
effectivity, so $B_{Y''}\leq0$ for every model, that is, $B\leq0$.  The
converse is trivial.
\end{proof}

\begin{remark}[Cartier is a direct limit, Weil an inverse limit]
\label{rem:ch38-direct-inverse}
\Cref{lem:ch38-determination-equivalence}\,\cref{itm:ch38-de:family} is the
statement that the natural map
$\varinjlim_Y\operatorname{CDiv}_{\mathbf R}(Y)
\to\varprojlim_Y\Div_{\mathbf R}(Y)$ of
\eqref{eq:appF-direct-inverse} is well defined and injective.  It is not
surjective: \Cref{ex:appF-non-cartier-bdivisor} produces a Weil b-divisor
supported on an infinite chain of exceptional primes, which no single
model determines.  Everything in this chapter takes place in the image, and
\Cref{rem:ch38-cartier-not-weil} explains why that is the right domain for
the envelope proof.
\end{remark}

\subsection{The evaluation pairing on divisorial points}
\label{sec:ch38-evaluation-divisorial}

For a divisorial valuation $v=c\ord_E$, choose a model
$\rho:Z\to Y$ on which $E$ appears and define
\[
                       B(v):=c\ord_E(\rho^*B_Y).
\tag{38.2}\label{eq:bj-bdivisor-evaluation}
\]
The value is independent of the chosen model.  We denote this coefficient
function on $X^{\divis}$ by $\psi_B$.

We now prove the independence, and record the properties of the pairing
that are used later.  All of them are elementary, and all of them are used.

\begin{proposition}
\label{prop:ch38-evaluation-welldefined}
Let $B$ be a real Cartier b-divisor over $X$ with determination
$(\pi:Y\to X,B_Y)$.
\begin{enumerate}
\item\label{itm:ch38-ev:welldef}
 The right-hand side of \eqref{eq:bj-bdivisor-evaluation} does not depend
 on the model $Z$ carrying $E$, nor on the determination of $B$;
 thus $\psi_B:X^{\divis}\to\mathbf R$ is well defined, and it coincides
 with the restriction to $X^{\divis}$ of the coefficient function
 \eqref{eq:bdivisor-coefficient-function} of the underlying Weil
 b-divisor.
\item\label{itm:ch38-ev:linear}
 $B\mapsto\psi_B$ is $\mathbf R$-linear and order preserving:
 $\psi_{aB+a'B'}=a\psi_B+a'\psi_{B'}$ for $a,a'\in\mathbf R$, and
 $B\leq B'$ if and only if $\psi_B\leq\psi_{B'}$ on $X^{\divis}$.
\item\label{itm:ch38-ev:homog}
 $\psi_B(sv)=s\psi_B(v)$ for $s\in\mathbf Q_{>0}$ and
 $v\in X^{\divis}$, and $\psi_B(v_{\triv})=0$. Its continuous
 extension to $X^{\val}$ is homogeneous for every real $s>0$.
\item\label{itm:ch38-ev:finite}
 $\psi_B$ is finite at every point of $X^{\divis}$; if $B\leq0$
 then $\psi_B\leq0$.
\item\label{itm:ch38-ev:pl}
 $\psi_B$ lies in $\PL_{\mathrm{hom}}(X)$ if $B$ is rational, and in
 the real span $\PL_{\mathrm{hom},\mathbf R}(X)$ in general; if
 $B\leq0$ is rational and $B_Y$ is relatively semiample, then
 $\psi_B\in\PL^+_{\mathrm{hom}}(X)$.
\end{enumerate}
\end{proposition}

\begin{proof}
\emph{\cref{itm:ch38-ev:welldef}.}  We use throughout the following
description of the evaluation.  If $C$ is an $\mathbf R$-Cartier
divisor on a normal variety $W$ and $v$ is a valuation of $k(W)$
admitting a centre $c_W(v)\in W$, set
$v(C):=v(f)$, where $f\in k(W)^\times_{\mathbf R}$ is a local equation
of $C$ at $c_W(v)$; this is well defined because two local equations
differ by a unit of $\cO_{W,c_W(v)}$, on which $v$ vanishes.  If
$\mu:W'\to W$ is a projective birational morphism of normal varieties and
$v$ has a centre on both, then
\begin{equation*}
 v(\mu^*C)=v(C),
\end{equation*}
because a local equation of $\mu^*C$ at $c_{W'}(v)$ is the pullback of
a local equation of $C$ at $c_W(v)=\mu(c_{W'}(v))$, and the pullback of
a rational function along a birational morphism is the same element of the
common function field.

Now let $E$ be a prime divisor over $X$, let $v:=c\ord_E$ with
$c>0$, and let $Z$ and $Z'$ be two models dominating the
determination $Y$ and carrying $E$, with structure morphisms
$\rho:Z\to Y$ and $\rho':Z'\to Y$.  Then
$\ord_E(\rho^*B_Y)=\ord_E(B_Y)=\ord_E(\rho'^*B_Y)$ by the displayed
identity, so \eqref{eq:bj-bdivisor-evaluation} does not depend on the model
carrying $E$.  If $(Y',B_{Y'})$ is another determination, choose by
\Cref{lem:appF-common-domination} a model $W$ dominating $Y$ and
$Y'$; there the pullbacks of $B_Y$ and $B_{Y'}$ agree, by
\Cref{lem:ch38-determination-equivalence}\,\cref{itm:ch38-de:equiv}, and the
displayed identity again gives $\ord_E(B_Y)=\ord_E(B_{Y'})$.  Finally, by
\Cref{lem:ch38-determination-equivalence}\,\cref{itm:ch38-de:family} the
coefficient of $E$ in the trace $B_Z$ of the underlying Weil b-divisor
on a model $Z$ carrying $E$ is exactly
$\ord_E(\rho^*B_Y)$; so \eqref{eq:bj-bdivisor-evaluation} is
\eqref{eq:bdivisor-coefficient-function}, and
\Cref{prop:ch35-weil-coefficient-bijection} applies.

\emph{\cref{itm:ch38-ev:linear}.}  Two Cartier b-divisors admit a common
determination, by \Cref{lem:appF-common-domination}, and on a fixed model
both pullback and the coefficient $\ord_E$ are $\mathbf R$-linear in
the divisor.  The order statement is
\Cref{prop:ch35-weil-coefficient-bijection}, whose proof identifies
$B\leq B'$ with the coefficientwise inequality over all primes over
$X$, and homogeneity turns the latter into
$\psi_B\leq\psi_{B'}$ on $X^{\divis}$.

\emph{\cref{itm:ch38-ev:homog}.}  By construction
$\psi_B(s\cdot c\ord_E)=sc\,\ord_E(\rho^*B_Y)=s\,\psi_B(c\ord_E)$ for
$s\in\mathbf Q_{>0}$, which preserves $X^{\divis}$; the same
local-equation calculation on $X^{\val}$ holds for all real $s>0$.
At the trivial valuation every nonzero local equation has value zero,
so $\psi_B(v_{\triv})=0$.

\emph{\cref{itm:ch38-ev:finite}.}  Finiteness is clear from
\eqref{eq:bj-bdivisor-evaluation}: $\ord_E(\rho^*B_Y)$ is a real
coefficient of a divisor.  If $B\leq0$ then every coefficient is
$\leq0$, by
\Cref{lem:ch38-determination-equivalence}\,\cref{itm:ch38-de:order}.

\emph{\cref{itm:ch38-ev:pl}.}  This is
\Cref{prop:cartier-homogeneous-pl}; the relatively semiample antieffective
case is Step~1 of its proof, which exhibits
$\psi_B=m^{-1}\log|\mathfrak b|$ for a nonzero coherent ideal
$\mathfrak b\subseteq\cO_X$.
\end{proof}

\begin{proposition}
\label{prop:ch38-evaluation-functorial}
Let $\beta:X'\to X$ be a projective birational morphism of normal
integral projective varieties, and identify $k(X')=k(X)$, hence
$X'^{\divis}=X^{\divis}$ and $X'^{\val}=X^{\val}$.  Then:
\begin{enumerate}
\item\label{itm:ch38-fun:cofinal} models over $X'$ are cofinal among
 models over $X$, and the induced map
 $\operatorname{CDiv}_{\mathbf R}(\mathfrak X)
 \to\operatorname{CDiv}_{\mathbf R}(\mathfrak X')$,
 $B\mapsto B'$, sending a determination $(Y,B_Y)$ with $Y\succeq X'$
 to itself, is an isomorphism;
\item\label{itm:ch38-fun:psi} $\psi_{B'}=\psi_B$ as functions on
 $X^{\divis}=X'^{\divis}$;
\item\label{itm:ch38-fun:class} if $(Y,B_Y)$ is a determination with
 $\pi:Y\to X$ factoring as $\pi=\beta\pi'$, then for every
 $\theta\in N^1(X)$,
\[
 \pi'^*(\beta^*\theta)+B_Y=\pi^*\theta+B_Y,
\tag{38.3}\label{eq:ch38-functorial-total}
\]
 so the total class of $B$ in the class $\theta$ over $X$ and the
 total class of $B'$ in the class $\beta^*\theta$ over $X'$ are the
 same class on $Y$; in particular one is nef if and only if the other
 is.
\end{enumerate}
\end{proposition}

\begin{proof}
\emph{\cref{itm:ch38-fun:cofinal}.}  A model $\pi':Y\to X'$ composed with
$\beta$ is a model of $X$, since a composite of projective birational
morphisms of normal projective varieties is one.  Conversely, given a model
$\pi:Y\to X$, \Cref{lem:appF-common-domination} applied to $Y$ and
$X'$ produces a model $Z$ dominating both, and $Z\to X'$ is a model of
$X'$ dominating $Y$.  Hence models over $X'$ are cofinal.  A direct
limit over a cofinal subcategory is the same direct limit, which is
\eqref{eq:appF-direct-inverse}; explicitly, a Cartier b-divisor over $X$
is determined on some $Y$, which after passing to a dominating model may
be taken over $X'$, and two determinations over $X'$ are equivalent over
$X'$ if and only if they are equivalent over $X$, since the common
dominations may be taken over $X'$.

\emph{\cref{itm:ch38-fun:psi}.}  Both sides are computed by
\eqref{eq:bj-bdivisor-evaluation} on a model over $X'$ carrying $E$,
and the formula does not refer to the base.

\emph{\cref{itm:ch38-fun:class}.}  Immediate from
$\pi'^*\beta^*=(\beta\pi')^*=\pi^*$ on $N^1$.  Nefness is a property of
the class on $Y$.
\end{proof}

\begin{remark}[Why \Cref{prop:ch38-evaluation-functorial} is stated now]
\label{rem:ch38-functorial-use}
\Cref{prop:ch38-evaluation-functorial} is the reason
\Cref{chap:bj-reflection} can manufacture a b-divisor upstairs and realize
it downstairs.  In the notation of that chapter, a homogeneous psh function
on $Y$ is approximated by homogeneous GF stages whose Cartier
coefficient b-divisors, by
\cref{itm:ch38-fun:cofinal}, are literally the same objects over $X$; the
only thing that changes is the class in which one measures positivity, and
\eqref{eq:ch38-functorial-total} says that the total class does not change
either.  All the work in \Cref{chap:bj-reflection} therefore goes into
\emph{correcting} the class $\beta^*\theta$ back to $\theta$, not into
transporting the divisor.
\end{remark}

\subsection{The three domains of definition}
\label{sec:ch38-domains}

The coefficient function of a Cartier b-divisor lives naturally on three
nested spaces, and the difference between them is the source of every
extended-real difficulty in this chapter.

\begin{proposition}
\label{prop:ch38-evaluation-domains}
Let $B$ be a real Cartier b-divisor with determination $(Y,B_Y)$.
\begin{enumerate}
\item\label{itm:ch38-dom:divis} On $X^{\divis}$ the function
 $\psi_B$ is defined by \eqref{eq:bj-bdivisor-evaluation}, is real
 valued, homogeneous, and determines $B$.
\item\label{itm:ch38-dom:val} On $X^{\val}$ the formula
 $\psi_B(v)=v(B_Y)$, computed through the canonical homeomorphism
 $Y^{\val}\simeq X^{\val}$, extends $\psi_B$ to a real-valued
 continuous function, and the extension is unique.
\item\label{itm:ch38-dom:an} If $B\leq0$ is rational and $B_Y$ is
 relatively semiample, then $\psi_B$ extends further to a
 continuous function
 $X^{\An}\to[-\infty,0]$, namely $m^{-1}\log|\mathfrak b|$ for the
 ideal $\mathfrak b=\pi_*\cO_Y(mB_Y)\subseteq\cO_X$ of Step~1 of the
 proof of \Cref{prop:cartier-homogeneous-pl}; this \emph{continuous}
 extension is unique, but the upper-semicontinuous extension is not.
\item\label{itm:ch38-dom:general} For a general real Cartier b-divisor
 there is no extension of $\psi_B$ to $X^{\An}$ given by a formula in
 $B_Y$: the expression $v(B_Y)$ is not defined at a point of proper
 support, and even when it is defined after a choice of lift it depends on
 that choice.
\end{enumerate}
\end{proposition}

\begin{proof}
\emph{\cref{itm:ch38-dom:divis}} is
\Cref{prop:ch38-evaluation-welldefined} together with
\Cref{prop:ch35-weil-coefficient-bijection}.

\emph{\cref{itm:ch38-dom:val}} is
\Cref{prop:ch35-cartier-continuous-extension}.  We recall the mechanism,
because \cref{itm:ch38-dom:general} is its exact failure.  A birational
morphism $\pi:Y\to X$ of normal projective varieties induces a bijection
$Y^{\val}\to X^{\val}$: both sides are the set of real valuations of the
common function field $k(X)=k(Y)$ admitting a centre, and properness
supplies the centre on either side.  Writing $B_Y$ locally as a
difference of two effective Cartier divisors, $v\mapsto v(B_Y)$ is a
difference of two ideal evaluations, both finite on $X^{\val}$ because a
nonzero ideal is not contained in the ideal of an irreducible component;
hence the difference is defined and continuous.  Uniqueness of the
extension follows from density of $X^{\divis}$ in $X^{\val}$.

\emph{\cref{itm:ch38-dom:an}.}  With $m\geq1$ chosen so that $mB_Y$ is
integral Cartier and $\cO_Y(mB_Y)$ is globally generated over $X$,
Step~1 of the proof of \Cref{prop:cartier-homogeneous-pl} produces a
nonzero coherent ideal $\mathfrak b=\pi_*\cO_Y(mB_Y)\subseteq\cO_X$ with
$\mathfrak b\cdot\cO_Y=\cO_Y(mB_Y)$ and
$\psi_B=m^{-1}\log|\mathfrak b|$ on $X^{\val}$.  The right-hand side is
defined on all of $X^{\An}$ with values in $[-\infty,0]$; it is upper
semicontinuous by \Cref{prop:ch35-hom-pl-semiring}, and in fact
\emph{continuous} into $[-\infty,0]$, because
$v\mapsto v(\mathfrak b)$ is continuous into $[0,+\infty]$ by
\Cref{thm:ch32-ideals-generate}\,\cref{itm:ch32-ig:continuous}.  Since
$X^{\val}$ is dense (\Cref{thm:triv-divisorial-density}) and continuous
functions agreeing on a dense subset agree, this is the unique continuous
extension.  Upper semicontinuity alone does \emph{not} force uniqueness:
with $X=\mathbf P^1$, $B=-\overline p$ and $m=1$, the function equal
to $\log|s_p|$ everywhere except at $v_{p,\triv}$, where it is given
the value $-5$, is still upper semicontinuous --- every neighbourhood
$\{v(\mathfrak m_p)>t\}$ of $v_{p,\triv}$ carries values $\leq-t$ ---
and it agrees with $\psi_B$ on $X^{\val}$.  This is why
\cref{itm:ch38-dom:an} is stated for the continuous extension, and why
\Cref{warn:ch38-proper-support-value} appeals to plurisubharmonicity and
not merely to upper semicontinuity.

\emph{\cref{itm:ch38-dom:general}} is
\Cref{ex:ch38-evaluation-no-lift} and
\Cref{ex:ch38-fractional-undefined} below.
\end{proof}

\begin{warning}{Three extensions, three hypotheses}
\label{warn:ch38-three-extensions}
The three extensions in \Cref{prop:ch38-evaluation-domains} are genuinely
different statements and are proved by different means.  The extension to
$X^{\val}$ is formal and uses only that $B_Y$ is Cartier.  The
extension to $X^{\An}$ in \cref{itm:ch38-dom:an} uses antieffectivity
\emph{and} relative semiampleness, and its value at a proper-support point
is read off from an ideal, not from a divisor.  The extension of
$\psi_B$ to $X^{\An}$ for a general $B$ with nef total class is not
available through a general divisor-evaluation formula at
proper-support points; it is the content of
\Cref{thm:bj-real-cartier-realization}, which produces it as the unique
homogeneous psh function with the given divisorial trace.  Confusing the
second and the third is the standard way of writing a proof of the
realization theorem that is not a proof.
\end{warning}

\subsection{The sign convention}
\label{sec:ch38-signs}

For an effective rational Cartier divisor $F$, the sign convention is
\[
                         \log|s_F|(v)=-v(F).
\tag{38.4}\label{eq:bj-divisor-weight-sign}
\]
Thus the coefficient function of $-\overline F$ is
$\log|s_F|$.

This is \eqref{eq:section-bdivisor-sign}, and the three unit tests
\eqref{eq:ch35-sign-tests} fix it beyond doubt.  We add the two tests that
are specific to the present chapter, and which the reader should apply
whenever a sign is in doubt below.

\begin{lemma}
\label{lem:ch38-sign-tests}
Let $F\geq0$ be an effective $\mathbf Q$-Cartier divisor on $X$ and
let $\pi:Y\to X$ be a model.
\begin{enumerate}
\item\label{itm:ch38-sign:pole}
 $\psi_{-\overline F}=\log|s_F|\leq0$, with value $-v(F)$ at
 $v\in X^{\val}$ and value $-\infty$ at a semivaluation whose support
 is contained in $\Supp F$.
\item\label{itm:ch38-sign:exceptional}
 If $G\geq0$ is an effective Cartier divisor on $Y$, then
 $\psi_{-\overline G}\leq0$, and
 $\psi_{-\overline G}(\ord_E)=-\ord_E(G)$ for every prime divisor $E$
 over $X$; in particular $\psi_{-\overline G}$ vanishes on the primes
 of $X$ itself when $G$ is $\pi$-exceptional.
\item\label{itm:ch38-sign:total}
 With $B\leq0$ determined on $Y$ and $D:=-B_Y\geq0$, the total class
 of \Cref{def:bj-nef-total} is $P=\pi^*\theta-D$.
 If $D'\geq D$, then $\psi_{-\overline{D'}}\leq\psi_{-\overline D}$,
 and the total classes differ by $-[D'-D]$. This coefficient order
 does not give an order of nef classes: an effective divisor need
 not be nef. The coefficient inequality agrees with
 $\mathfrak b\subseteq\mathfrak b'\Rightarrow\log|\mathfrak b|\leq
 \log|\mathfrak b'|$.
\end{enumerate}
\end{lemma}

\begin{proof}
\emph{\cref{itm:ch38-sign:pole}.}  By
\eqref{eq:bj-bdivisor-evaluation} and
\Cref{prop:ch38-evaluation-welldefined}\,\cref{itm:ch38-ev:linear},
$\psi_{-\overline F}(v)=-v(F)$ on $X^{\divis}$, which is
\eqref{eq:bj-divisor-weight-sign}.  The values on $X^{\val}$ and at
proper-support points are \eqref{eq:divisor-section-pole}, applied to the
canonical section $s_F$ of $\cO_X(F)$.

\emph{\cref{itm:ch38-sign:exceptional}.}  Apply
\eqref{eq:bj-bdivisor-evaluation} on a model dominating $Y$.  If $G$ is
$\pi$-exceptional and $E\subset X$ is a prime divisor, then the strict
transform $\widetilde E\subset Y$ is not a component of $G$, so
$\ord_E(G)=\operatorname{coeff}_{\widetilde E}(G)=0$.

\emph{\cref{itm:ch38-sign:total}.}  The identity
$P=\pi^*\theta+B_Y=\pi^*\theta-D$ is
\eqref{eq:bj-nef-total} with $D=-B_Y$.  The coefficient inequality follows from
\cref{itm:ch38-ev:linear} of
\Cref{prop:ch38-evaluation-welldefined}; the total-class difference
is an identity. For example, on the blowup of a surface at a point,
the effective exceptional divisor has negative intersection with itself,
so its class is not nef.
\end{proof}

\begin{remark}[Trace and analytic extension]
\label{rem:ch38-trace-analytic-extension}
For $v\in X^{\divis}$, the number $v(F)$ is finite.  At a
semivaluation supported in $F$, however, $\log|s_F|=-\infty$.  The
function $\psi_B$ is initially only a divisorial trace.  The theorem
below constructs its psh extension and never defines that extension by
subtracting divisor functions at a proper-support point.
\end{remark}

\subsection{Nef total classes}
\label{sec:ch38-nef-total}

\begin{definition}[Nef total class]
\label{def:bj-nef-total}
Let $\theta\in N^1(X)_{\mathbf R}$, and let $B\le0$ be determined on
$\pi:Y\to X$.  The b-divisor is \emph{nef-total in the class $\theta$}
if
\[
                         P:=\pi^*\theta+B_Y
\tag{38.5}\label{eq:bj-nef-total}
\]
is nef on $Y$.
\end{definition}

Neither $\theta$ nor $B_Y$ is required to be nef separately.

\begin{lemma}
\label{lem:ch38-nef-total-intrinsic}
Let $B\leq0$ be a real Cartier b-divisor with determination
$(\pi:Y\to X,B_Y)$ and let $\theta\in N^1(X)$.
\begin{enumerate}
\item\label{itm:ch38-nt:indep} Nef-totality in the class $\theta$ does
 not depend on the determination: if $\mu:Y'\to Y$ is a domination, then
 $\pi^*\theta+B_Y$ is nef if and only if
 $(\pi\mu)^*\theta+B_{Y'}$ is.
\item\label{itm:ch38-nt:exceptional} If $B$ is nef-total in the class
 $\theta$ and its Weil trace $B_X$ is $\mathbf R$-Cartier on $X$,
 then the exceptional part $B-\overline{B_X}$ satisfies
 $B-\overline{B_X}\leq0$, and it is nef-total in the class
 $\theta+[B_X]$.
\item\label{itm:ch38-nt:trace} If $B$ is nef-total in the class
 $\theta$ and $B_X$ is $\mathbf R$-Cartier, put
 $\eta:=\theta+[B_X]$. Then $\eta\cdot\Gamma\geq0$ for every
 integral curve $\Gamma\subset X$ not contained in the image of
 the non-isomorphism locus of $\pi$. If $\dim X\leq2$, this covers
 every curve, so $\eta$ is nef and constants are $\eta$-psh.
 Nefness of $\eta$ does not follow in arbitrary dimension.
\end{enumerate}
\end{lemma}

\begin{proof}
\emph{\cref{itm:ch38-nt:indep}.}  $B_{Y'}=\mu^*B_Y$, so
$(\pi\mu)^*\theta+B_{Y'}=\mu^*(\pi^*\theta+B_Y)$, and nefness of a class
on $Y$ is equivalent to nefness of its pullback to a dominating model.
This is \Cref{lem:appF-nef-determination}; the nontrivial half is that
$\mu^*\alpha$ nef implies $\alpha$ nef, which follows from the
projection formula \eqref{eq:appF-projection}.  Indeed, let
$\Gamma\subset Y$ be an integral curve.  Since $\mu$ is proper and
surjective, $\mu^{-1}(\Gamma)$ has an irreducible component $T$
dominating $\Gamma$; cutting $T$ with $\dim T-1$ general very ample
divisors on $Y'$ yields an integral curve $C\subset Y'$ still
dominating $\Gamma$, so that $\mu_*C=d\,\Gamma$ with
$d=\deg(C\to\Gamma)>0$.  Then
$d(\alpha\cdot\Gamma)=\alpha\cdot\mu_*C=\mu^*\alpha\cdot C\geq0$.

\emph{\cref{itm:ch38-nt:exceptional}.}  Put $C:=B-\overline{B_X}$, a
Cartier b-divisor determined on $Y$ by $C_Y=B_Y-\pi^*B_X$, and
$\eta:=\theta+[B_X]$.  Then
\[
 \pi^*\eta+C_Y=\pi^*\theta+\pi^*[B_X]+B_Y-\pi^*B_X=\pi^*\theta+B_Y=P,
\]
which is nef; so $C$ is nef-total in the class $\eta$.  Moreover
$C_Y$ is $\pi$-exceptional, since
$\pi_*C_Y=B_X-B_X=0$, and it is $\pi$-nef, since
$C_Y=P-\pi^*\eta$ with $P$ nef, hence $\pi$-nef, and $\pi^*\eta$
numerically trivial on $\pi$-contracted curves.  The negativity lemma
\cite[Lemma~3.39]{KM08}, applied to the $\pi$-exceptional divisor
$-C_Y$, whose negative $C_Y$ is $\pi$-nef, gives: $-C_Y$ is
effective if and only if $\pi_*(-C_Y)$ is, and the latter is $0$.
Hence $C_Y\leq0$, that is, $C\leq0$.

\emph{\cref{itm:ch38-nt:trace}.} Keep the notation of
\cref{itm:ch38-nt:exceptional}. For a curve $\Gamma$ outside the
stated exceptional image, its strict transform $\widetilde\Gamma$
is a curve not contained in $\Supp(-C_Y)$, and
$\pi_*\widetilde\Gamma=\Gamma$. Thus
\[
 \eta\cdot\Gamma
 =P\cdot\widetilde\Gamma+(-C_Y)\cdot\widetilde\Gamma\geq0.
\]
The first term is nonnegative by nefness; the second is nonnegative
because $-C_Y$ is an effective $\mathbf R$-Cartier divisor and the
curve is not contained in its support. A proper birational morphism
onto a normal variety is an isomorphism in codimension one.
For a surface its non-isomorphism image is therefore finite, and
no integral curve is contained in it. In higher dimension curves
can be contained in that image, and this argument does not apply.
\end{proof}

\begin{warning}{The hypothesis does not make $\theta$ nef}
\label{warn:ch38-theta-not-nef}
\Cref{lem:ch38-nef-total-intrinsic}\,\cref{itm:ch38-nt:trace} gives
nefness of $\theta+[B_X]$ on a surface when $B_X$ is
$\mathbf R$-Cartier. Even there it does not imply that $\theta$
is nef.  Let $X$ be the blowup of
$\mathbf P^2$ at a point, with exceptional curve $E$, let
$\theta:=[E]$, and let $B:=-\overline E$, determined on $X$ itself.
Then $B\leq0$, the total class is
$\theta+[B_X]=[E]-[E]=0$, which is nef, and the hypotheses of
\Cref{thm:bj-real-cartier-realization} are satisfied; but
$\theta\cdot E=E^2=-1<0$, so $\theta$ is not nef.  The corresponding
coefficient function is $\psi_B=\log|s_E|$, which is a homogeneous
Fubini--Study function for the line bundle $\cO_X(E)$, by
\Cref{prop:hom-fs-ideal}, since
$\cO_X(E)\otimes\cO_X(-E)\simeq\cO_X$ is globally generated.

There is also a higher-dimensional obstruction to nefness of the
trace class itself. Let $X$ be the blowup of $\mathbf P^3$ at two
distinct points $p_1,p_2$, with exceptional divisors $E_1,E_2$,
and let $C$ be the strict transform of their joining line.
Write $H$ for the pullback hyperplane class and
$\theta=H-[E_1]-[E_2]$, so $\theta\cdot C=1-1-1=-1$.
The pencil of planes containing that line has base ideal
$\mathcal I_C$ after removing $E_1+E_2$: locally at either point
the two linear equations pull back with one common exceptional
factor, and their residual equations define $C$.
On the blowup $\pi:Y\to X$ of $C$, with exceptional divisor $T$,
the line bundle of class $\pi^*\theta-[T]$ is generated by this
pencil, hence nef. The b-divisor determined by $-T$ has trace zero
on $X$ and is nef-total in class $\theta$, although
$\theta+[B_X]=\theta$ is not nef. Thus the surface trace conclusion
cannot be used for curves lying in a higher-dimensional exceptional
centre.

Three consequences follow. First, the hypothesis
of \Cref{thm:bj-real-cartier-realization} is a positivity statement about
the \emph{pair} $(\theta,B)$, and the positivity may be entirely consumed
by the negative part of $B$.  Second, no step of the proof below may
quote nefness of $\theta$ as if it were part of the hypotheses; and none
does.  All the positivity used in the construction lives in the
\emph{ample} class $\omega=\theta+[F]$ manufactured by the activation
step of \Cref{sec:ch38-activation}, and $\theta$ itself appears only as
the target class of the final Siu division.  Third,
\Cref{lem:ch34-psh-nonempty} makes the distinction precise: a bounded
$\theta$-psh function forces $\theta$ to be nef, whereas an unbounded
psh function may live in a non-nef pseudoeffective class.  Here
$\log|s_E|\in\PSH_{\mathrm{hom},X}([E])$ is unbounded, so there is no
contradiction with that lemma.  Accordingly, no argument in
\Crefrange{chap:bj-cartier-realization}{chap:bj-envelope-consequences} may
silently impose nefness on a target class; whenever constants are used as
psh functions, nefness is stated separately.
\end{warning}

\begin{example}[The hypothesis is not vacuous and not automatic]
\label{ex:ch38-usc-not-psh}
Let $X=\mathbf P^1_k$, $p\in X(k)$, and let
$B:=-c\,\overline p$ for a real number $c>0$; this is an antieffective
real Cartier b-divisor determined on $X$ itself, and
$\psi_B(s\,\ord_p)=-cs$, $\psi_B(s\,\ord_q)=0$ for $q\neq p$.  Put
$\theta=c_1(\cO(d))$, so $\deg\theta=d$.  The total class is the
number $d-c$, and it is nef exactly when $c\leq d$.

For $c\leq d$, \Cref{prop:ch35-curve-classification} shows that
$\psi_B$ is the divisorial trace of the homogeneous $\theta$-psh
function $c\log|s_p|$, which equals $-cs$ on the ray through $p$,
$0$ on every other ray, and $-\infty$ at $v_{p,\triv}$.  This is
\Cref{thm:bj-real-cartier-realization} in dimension one.

For $c>d$ the function $\psi_B$ is still homogeneous, still the
coefficient function of an antieffective real Cartier b-divisor, and still
extends to an upper-semicontinuous function on $X^{\An}$ --- namely
$c\log|s_p|$ again.  But it is $\theta$-psh for no class $\theta$ of
degree $d<c$, by \Cref{prop:ch35-curve-classification}: the slope
inequality $\sum_rc_r\leq\deg\theta$ fails.  Thus
\Cref{thm:bj-real-cartier-realization} is not an assertion that
homogeneous PL functions extend, nor that the extension is automatically
psh; it is an assertion that the numerical hypothesis \eqref{eq:bj-nef-total}
is exactly what makes the extension psh.  The example says nothing, however, about weakening the
hypothesis to pseudoeffectivity of $P$: on a curve the pseudoeffective
and nef cones coincide, so the two hypotheses are the same condition here
and no one-dimensional example can separate them.  That question first
arises in dimension $\geq2$, where $\operatorname{Psef}(X)$ may strictly
contain $\Nef(X)$ --- as it does on a blowup of $\mathbf P^2$, though not
on $\mathbf P^2$ itself, where both cones are $\mathbf R_{\geq0}H$ --- and
it is not addressed here.
\end{example}

\subsection{A Cartier b-divisor on a blown-up surface, evaluated}
\label{sec:ch38-surface-example}

The next example is carried to explicit numbers and is used four times
below: in \Cref{ex:ch38-fs-from-bdivisor-and-back} to compute the
Fubini--Study dictionary in both directions, in
\Cref{ex:ch38-activation-surface} to exhibit the fixed-pole data, in
\Cref{ex:ch38-buffer-surface} to exhibit the ample buffer, and in
\Cref{ex:ch38-drift-irrational} to exhibit an unremovable class drift.

\begin{example}[Evaluation at divisorial and quasimonomial points]
\label{ex:ch38-blowup-surface-evaluation}
\index{b-divisor@$b$-divisor!evaluated on a surface}
Let $X=\mathbf P^2_k$ with homogeneous coordinates $x,y,z$, let
$p=[0:0:1]$, and work on the affine chart $z=1$ with coordinates
$(x,y)$.  Let $\pi_1:Y_1\to X$ be the blowup of $p$, with
exceptional prime $E_1$; let $q\in E_1$ be the point corresponding to
the tangent direction $y=0$, let $\pi_2:Y_2\to Y_1$ be the blowup of
$q$, with exceptional prime $E_2$, and let $E_1'$ be the strict
transform of $E_1$ on $Y_2$.  Write $\pi=\pi_1\pi_2:Y_2\to X$.  This
is the tower of \Cref{ex:ch32-blowup-surface} over the affine chart, and
we use its coordinates: on $Y_2$ there is a chart $(z_1,z_2)$ with
\[
 x=z_1z_2,\qquad y=z_1z_2^2,
 \qquad
 E_1'=\{z_1=0\},\quad E_2=\{z_2=0\},
\]
which is \eqref{eq:ch32-blowup-coordinates}.  For
$\alpha=(\alpha_1,\alpha_2)\in\mathbf R^2_{\geq0}\setminus\{0\}$ let
$v_\alpha$ be the quasimonomial valuation at the zero-dimensional
stratum $E_1'\cap E_2$ with $v_\alpha(z_i)=\alpha_i$, as in
\eqref{eq:quasimonomial-valuation}.  Then
$v_{(1,0)}=\ord_{E_1}$ and $v_{(0,1)}=\ord_{E_2}$, and
\[
 v_\alpha(x)=\alpha_1+\alpha_2,
 \qquad
 v_\alpha(y)=\alpha_1+2\alpha_2 .
\tag{38.6}\label{eq:ch38-surface-coordinates}
\]

\emph{The b-divisor.}  Let $B$ be the Cartier b-divisor determined on
$Y_2$ by
\[
 B_{Y_2}:=-2E_1'-3E_2 .
\tag{38.7}\label{eq:ch38-surface-bdivisor}
\]
It is antieffective, and it is genuinely determined on $Y_2$ and on no
lower model: its trace on $Y_1$ is
$\pi_{2*}(-2E_1'-3E_2)=-2E_1$, which is nonzero, and
$\pi_2^*(-2E_1)=-2E_1'-2E_2\neq B_{Y_2}$; its trace on $X$ is
$\pi_*B_{Y_2}=0$, so $B$ is exceptional in the sense of
\Cref{def:ch35-exceptional-bdivisor}.

\emph{Evaluation at divisorial points.}  By
\eqref{eq:bj-bdivisor-evaluation},
\[
 \psi_B(\ord_{E_1})=-2,
 \qquad
 \psi_B(\ord_{E_2})=-3,
 \qquad
 \psi_B(\ord_C)=0
\tag{38.8}\label{eq:ch38-surface-divisorial-values}
\]
for every prime divisor $C\subset X$, the last because the strict
transform of $C$ on $Y_2$ is not a component of $B_{Y_2}$.  On the continuous $X^{\val}$ extension, for a
positive real $s$, homogeneity gives
$\psi_B(s\ord_{E_1})=-2s$ and $\psi_B(s\ord_{E_2})=-3s$.

\emph{Evaluation at quasimonomial points.}  In the chart $(z_1,z_2)$ the
divisor $B_{Y_2}$ has local equation $z_1^{-2}z_2^{-3}$, so
\[
 \psi_B(v_\alpha)=v_\alpha(B_{Y_2})=-2\alpha_1-3\alpha_2 .
\tag{38.9}\label{eq:ch38-surface-quasimonomial-value}
\]
This is consistent with \eqref{eq:ch38-surface-divisorial-values} at
$\alpha=(1,0)$ and $\alpha=(0,1)$, and it is linear in $\alpha$,
as it must be: on the two-dimensional cone of quasimonomial valuations at
a fixed zero-dimensional stratum, the evaluation of a Cartier divisor
determined on that model is the linear function read off from the local
equation.  At the irrational point
$\alpha=(1,\sqrt2)$ of
\Cref{ex:ch32-blowup-surface}\,\cref{itm:ch32-bs:irrational} the value is
$-2-3\sqrt2$, an irrational number attached to a rational b-divisor;
irrationality of a value is a property of the point, not of the divisor.

\emph{The ideal.}  By
\Cref{prop:ch38-evaluation-domains}\,\cref{itm:ch38-dom:an}, with
$m=1$, the ideal $\mathfrak b:=\pi_*\cO_{Y_2}(B_{Y_2})$ consists of
those $f\in\cO_X$ with $\ord_{E_1}(f)\geq2$ and
$\ord_{E_2}(f)\geq3$.  Since $\ord_{E_1}(x^iy^j)=i+j$ and
$\ord_{E_2}(x^iy^j)=i+2j$, the monomials in $\mathfrak b$ are those
with $i+j\geq2$ and $i+2j\geq3$, and the minimal ones are
$x^3$, $xy$, $y^2$.  Hence
\[
 \mathfrak b=(x^3,\;xy,\;y^2)
 \qquad\text{near }p,\qquad
 \mathfrak b=\cO_X\ \text{elsewhere}.
\tag{38.10}\label{eq:ch38-surface-ideal}
\]
The check that $\psi_B=\log|\mathfrak b|$ on the quasimonomial cone is
immediate from \eqref{eq:ch38-surface-coordinates}:
\[
 v_\alpha(x^3)=3\alpha_1+3\alpha_2,
 \quad
 v_\alpha(xy)=2\alpha_1+3\alpha_2,
 \quad
 v_\alpha(y^2)=2\alpha_1+4\alpha_2,
\]
whose minimum is $2\alpha_1+3\alpha_2$, because
$3\alpha_1+3\alpha_2-(2\alpha_1+3\alpha_2)=\alpha_1\geq0$ and
$2\alpha_1+4\alpha_2-(2\alpha_1+3\alpha_2)=\alpha_2\geq0$.  Thus
$\log|\mathfrak b|(v_\alpha)=-2\alpha_1-3\alpha_2=\psi_B(v_\alpha)$,
in agreement with \eqref{eq:ch38-surface-quasimonomial-value}.  The
colength of $\mathfrak b$ at $p$ is $4$, the monomials outside
$\mathfrak b$ being $1,x,y,x^2$.

\emph{Total classes.}  On $Y_2$ write $H:=\pi^*c_1(\cO(1))$,
$e_1:=[\pi_2^*E_1]=[E_1'+E_2]$ and $e_2:=[E_2]$.  The intersection
table is
\[
 H^2=1,\quad H\cdot e_1=H\cdot e_2=0,\quad
 e_1^2=e_2^2=-1,\quad e_1\cdot e_2=0,
\tag{38.11}\label{eq:ch38-surface-intersection}
\]
where $e_1\cdot e_2=\pi_2^*E_1\cdot E_2=E_1\cdot\pi_{2*}E_2=0$ by the
projection formula \eqref{eq:appF-projection}, and the remaining entries
are \eqref{eq:appF-blowup-table} applied twice.  Since
$[E_1']=e_1-e_2$,
\[
 [B_{Y_2}]=-2(e_1-e_2)-3e_2=-2e_1-e_2 .
\]
Taking $\theta=c_1(\cO(d))$, the total class of
\Cref{def:bj-nef-total} is
\[
 P=dH-2e_1-e_2 .
\tag{38.12}\label{eq:ch38-surface-total}
\]
The extremal curves of $Y_2$ are $E_2$, the $(-2)$-curve $E_1'$,
and the strict transform $\Lambda$ of the line through $p$ in the
direction $q$, with classes $e_2$, $e_1-e_2$ and
$H-e_1-e_2$.  Using \eqref{eq:ch38-surface-intersection},
\[
 P\cdot E_2=1,\qquad
 P\cdot E_1'=1,\qquad
 P\cdot\Lambda=d-3 .
\tag{38.13}\label{eq:ch38-surface-nef-test}
\]
Hence $P$ is nef if and only if $d\geq3$, and ample if and only if
$d\geq4$.  For $d=3$ the class $P$ is nef and not ample, which is
exactly the situation \Cref{thm:bj-real-cartier-realization} is designed
for and the situation in which the ample buffer of
\Cref{sec:ch38-buffer} is needed.  For $d=3$ the total class is
moreover semiample, as \Cref{ex:ch38-fs-from-bdivisor-and-back} verifies
by exhibiting the linear system.
\end{example}

\begin{example}[A lift is needed, and lifts disagree]
\label{ex:ch38-evaluation-no-lift}
Continue with $X=\mathbf P^2$ and the blowup $\pi_1:Y_1\to X$ of the
point $p$, and let $v\in X^{\An}$ be a point of proper support, say
$v=v_{p,\triv}$, the trivial semivaluation of the residue field at
$p$.  The expression $v(B_{Y_1})$ has no meaning as it stands: $v$
is a semivaluation on $\cO_X$, and to evaluate it on a local equation of
a divisor on $Y_1$ one must first lift $v$ to a point of
$Y_1^{\An}$.  The fibre of $\pi_1^{\An}:Y_1^{\An}\to X^{\An}$ over
$v_{p,\triv}$ consists of the semivaluations $w$ on $\cO_{Y_1}$ with
$w|_{\cO_X}=v_{p,\triv}$; since
$\mathfrak m_p\cdot\cO_{Y_1}=\cO_{Y_1}(-E_1)$, such a $w$ satisfies
$w(\mathcal I_{E_1})=+\infty$, so its support contains $E_1$, and the
fibre is the set of semivaluations of $\cO_{Y_1}$ with support
$\mathcal I_{E_1}$ or the maximal ideal of a closed point of $E_1$.
The fibre is therefore infinite: it contains
$v_{E_1,\triv}$, every $v_{\eta,\triv}$ with $\eta\in E_1$ a closed
point, and all the rays above them.

Now take the Cartier b-divisor $C$ determined on $Y_1$ by
$C_{Y_1}:=\pi_1^*L-E_1$, where $L\subset X$ is the line $y=0$, which
passes through $p$.  Let $\widetilde L\subset Y_1$ be its strict
transform, so $C_{Y_1}=\widetilde L$.  Evaluating on lifts of
$v_{p,\triv}$:
\[
 v_{E_1,\triv}(C_{Y_1})=v_{E_1,\triv}(\widetilde L)=0,
 \qquad
 v_{\eta_0,\triv}(C_{Y_1})=+\infty
 \quad\text{for }\eta_0:=E_1\cap\widetilde L,
\]
because $\widetilde L$ does not contain $E_1$ but does contain
$\eta_0$.  The two lifts of the same point $v_{p,\triv}$ give the two
values $0$ and $+\infty$.  Hence there is no function
``$v\mapsto v(C_{Y_1})$'' on $X^{\An}$, and
\Cref{prop:ch38-evaluation-domains}\,\cref{itm:ch38-dom:general} is proved.
Note that $C$ is not antieffective; \Cref{ex:ch38-fractional-undefined}
shows that antieffectivity alone does not repair the situation either,
once differences are formed.
\end{example}

\begin{example}[A fractional difference undefined at a proper-support
point]
\label{ex:ch38-fractional-undefined}
Let $X=\mathbf P^2$, let $L_1$ and $L_2$ be two distinct lines, and
let $C:=\overline{L_1}-\overline{L_2}$, a real (indeed integral) Cartier
b-divisor determined on $X$ itself.  Its coefficient function on
$X^{\val}$ is $\psi_C(v)=v(L_1)-v(L_2)$, a difference of two finite
numbers, hence continuous and real valued, in accordance with
\Cref{prop:ch38-evaluation-domains}\,\cref{itm:ch38-dom:val}.  At the
proper-support point $v=v_{\xi,\triv}$ with
$\xi\in L_1\cap L_2$, however, both terms equal $+\infty$ and the
difference is undefined by ledger entry~(4) of
\Cref{sec:extended-arithmetic}.  This is the phenomenon recorded in the
warning after \Cref{prop:ch35-hom-pl-semiring}: a homogeneous PL function
is a difference of two ideal functions and therefore lives on
$X^{\val}$, not on $X^{\An}$.  It is the reason
\Cref{thm:bj-real-cartier-realization} is stated with $B\leq0$ and
produces its extension by a decreasing limit rather than by a formula, and
it is the mechanism behind
\Cref{ex:ch38-naive-realization-fails}.
\end{example}

\section{Model and homogeneous Fubini--Study functions}
\label{sec:ch38-model-fs}

\subsection{Rational coefficient functions}
\label{sec:ch38-rational-coefficients}

Let $L\in\Pic(X)_{\mathbf Q}$, and let $C\le0$ be a rational Cartier
b-divisor determined on $\pi:Y\to X$.  If the actual rational line bundle
$\pi^*L+C_Y$ is absolutely semiample, the Cartier b-divisor description
of homogeneous PL functions gives
\[
 \psi_C\in\cH_{\mathrm{hom}}(L)
       =\cH^{\mathrm{gf}}_0(L)
       \subset\PSH_{\mathrm{hom},X}(c_1(L)).
\tag{38.14}\label{eq:bj-rational-coefficient}
\]
The construction is ordered: $C_1\le C_2$ implies
$\psi_{C_1}\le\psi_{C_2}$.  Every homogeneous coefficient function
vanishes at the trivial valuation.

The first membership in \eqref{eq:bj-rational-coefficient} is
\Cref{thm:ch35-semiample-criterion}, the equality is
\Cref{def:homogeneous-fs-psh}, and the inclusion is the same definition
together with \eqref{eq:ch34-hierarchy}.  In every application below the
class $c_1(L)$ is \emph{ample}, so that the last inclusion is used only
in the range in which \Cref{lem:ch34-psh-nonempty} guarantees that
$\PSH_X(c_1(L))$ is nonempty; \Cref{warn:ch38-theta-not-nef} explains why
this restriction is worth observing.  The ordering statement is
\Cref{prop:ch38-evaluation-welldefined}\,\cref{itm:ch38-ev:linear}, and the
value at $v_{\triv}$ is
\Cref{prop:ch38-evaluation-welldefined}\,\cref{itm:ch38-ev:homog}.

\begin{example}[A fixed divisor]
\label{ex:bj-fixed-divisor-coefficient}
Suppose $F\ge0$ is rational Cartier and $L-F$ is semiample.  Taking
$C=-\overline F$ in \eqref{eq:bj-rational-coefficient} gives
$\psi_C=\log|s_F|$.  It is nonpositive, vanishes at the trivial
valuation, and equals $-\infty$ at a semivaluation supported in $F$.
Thus a generically finite function need not be finite-valued everywhere.
\end{example}

The example is the smallest instance of the phenomenon that governs the
whole chapter, and it is worth naming its three features separately.  The
b-divisor $-\overline F$ is determined on $X$ itself, so no birational
model intervenes.  Its coefficient function is finite on $X^{\val}$ and
takes the value $-\infty$ exactly on the analytic subset
$\Supp(F)^{\An}$ of proper-support points, by
\eqref{eq:divisor-section-pole}.  And, crucially, the value $-\infty$ is
attained at the pole, where the finite approximants tend to $-\infty$:
$\log|s_F|$ is a decreasing limit of the
finite Fubini--Study functions $\max\{\log|s_F|,-N\}$, and the pole
appears only in the limit, as in \Cref{ex:ch34-decreasing-to-pole}.  Every
occurrence of $-\infty$ in this chapter arises in this way.

\subsection{The dictionary, in both directions}
\label{sec:ch38-dictionary-fs}

We isolate the statement that the realization theorem applies at every
rational stage.  Nothing in it is new; the point is to have the two
directions and their hypotheses displayed together, since the proof of
\Cref{thm:bj-real-cartier-realization} alternates between them.

\begin{proposition}
\label{prop:ch38-model-fs-dictionary}
Let $X$ be normal integral projective and $L\in\Pic(X)_{\mathbf Q}$.
\begin{enumerate}
\item\label{itm:ch38-dict:forward}
 Let $C$ be a rational Cartier b-divisor with determination
 $\pi:Y\to X$ such that $C\leq0$ and $\pi^*L+C_Y$ is semiample.
 Choose $m\geq1$ with $mC_Y$ integral Cartier, $mL$ a line bundle
 and $m(\pi^*L+C_Y)$ globally generated.  Then
 $\mathfrak b_m\subseteq\pi_*\cO_Y(mC_Y)\subseteq\cO_X$, defined
 by the coefficients of the global sections of
 $m(\pi^*L+C_Y)$, is a nonzero coherent ideal,
 $mL\otimes\mathfrak b_m$ is globally generated,
 $\mathfrak b_m\cdot\cO_Y=\cO_Y(mC_Y)$, and
\[
 \psi_C=m^{-1}\log|\mathfrak b_m|\in\cH_{\mathrm{hom}}(L).
\tag{38.15}\label{eq:ch38-forward-dictionary}
\]
\item\label{itm:ch38-dict:backward}
 Conversely, let $\psi\in\cH_{\mathrm{hom}}(L)$, and write
 $\psi=m^{-1}\log|\mathfrak b|$ with $0\neq\mathfrak b\subseteq\cO_X$
 and $mL\otimes\mathfrak b$ globally generated, as in
 \Cref{prop:hom-fs-ideal}.  Let $\rho:Z\to X$ be the normalized blowup
 of $\mathfrak b$, with $\mathfrak b\cdot\cO_Z=\cO_Z(-E)$.  Then the
 Cartier b-divisor $C$ determined on $Z$ by
 $C_Z:=-m^{-1}E$ satisfies $C\leq0$, $\rho^*L+C_Z$ is semiample,
 and $\psi_C=\psi$.
\item\label{itm:ch38-dict:unique}
 The Cartier b-divisor in \cref{itm:ch38-dict:backward} is unique: if
 $\psi_{C}=\psi_{C'}$ for two Cartier b-divisors, then $C=C'$.
\item\label{itm:ch38-dict:order}
 Both directions are order preserving, and
 $\psi_C\leq0$ with $\psi_C(v_{\triv})=0$.
\end{enumerate}
\end{proposition}

\begin{proof}
\emph{\cref{itm:ch38-dict:forward}.} This is the implication
$\Leftarrow$ of \Cref{thm:ch35-semiample-criterion}. Normality gives
$\pi_*\cO_Y=\cO_X$, so
$\pi_*\cO_Y(mC_Y)\subseteq\cO_X$. The projection formula identifies
the global sections upstairs with
$H^0(X,mL\otimes\pi_*\cO_Y(mC_Y))$. Define
$mL\otimes\mathfrak b_m$ as the image of their evaluation map.
It is globally generated by construction, and is nonzero because
the globally generated line bundle upstairs has a nonzero section.
Pulling these sections to $Y$ gives
$\mathfrak b_m\cdot\cO_Y=\cO_Y(mC_Y)$. Evaluating on $X^{\val}$ gives
$\psi_C=m^{-1}\log|\mathfrak b_m|$.
Global generation upstairs does not assert that the whole
pushforward sheaf is globally generated downstairs; the section
image is the ideal used here.

\emph{\cref{itm:ch38-dict:backward}.}  This is Step~2 of the proof of
\Cref{prop:cartier-homogeneous-pl} together with the implication
$\Rightarrow$ of \Cref{thm:ch35-semiample-criterion}.  On the normalized
blowup, $-E$ is $\rho$-ample, and global generation of
$mL\otimes\mathfrak b$ says that
$\rho^*(mL)\otimes\cO_Z(-E)=m(\rho^*L+C_Z)$ is globally generated,
whence $\rho^*L+C_Z$ is semiample.  The evaluation
$\psi_C(v)=-m^{-1}v(E)=-m^{-1}v(\mathfrak b)=\psi(v)$ holds on
$X^{\val}$, hence on $X^{\divis}$.

\emph{\cref{itm:ch38-dict:unique}} is
\Cref{prop:ch35-weil-coefficient-bijection}: the coefficient function
determines all coefficients, and a Cartier b-divisor is a Weil b-divisor.

\emph{\cref{itm:ch38-dict:order}} is
\Cref{prop:ch38-evaluation-welldefined}\,\cref{itm:ch38-ev:linear} and
\cref{itm:ch38-ev:homog}, together with
\Cref{prop:ch35-hom-pl-semiring}, which gives
$\log|\mathfrak b|\leq0$ with value $0$ at $v_{\triv}$.
\end{proof}

\begin{remark}[Semiample, not nef, and not ample]
\label{rem:ch38-semiample-not-nef}
\Cref{prop:ch38-model-fs-dictionary} is a statement about linear series.
It is the reason the construction below inserts an ample buffer: the
hypothesis of \Cref{thm:bj-real-cartier-realization} is that the total
class $P$ is \emph{nef}, and a nef class need not be semiample or supply the globally
generated multiple required by this dictionary.  What the buffer does is replace $P$ by a family of
\emph{ample} total classes, for which semiampleness is automatic on the
rational members, at the cost of a controlled error in the numerical class
and in the divisor.  The two errors are accounted for in
\Cref{sec:ch38-classes} and \Cref{sec:ch38-approximation} respectively.
The corresponding warning in
\Cref{chap:bj-homogeneity-bdivisors}, after
\Cref{thm:ch35-semiample-criterion}, says the same thing from the other
side.
\end{remark}

\subsection{The two gradings}
\label{sec:ch38-two-gradings}

A weighted section system \eqref{eq:gf-fs-presentation} carries two
indices of different natures: the integer section degree and
the weights (real in a general presentation, integral after clearing
denominators in rational Rees data), and the Rees and base-ideal
arguments of
\Crefrange{chap:bj-test-configurations}{chap:bj-smooth-characteristic}
manipulate both.  Confusing them is the commonest arithmetic error in this
part of the subject, and the homogeneous category is precisely the one in
which the second becomes trivial.  We record the facts.

\begin{proposition}
\label{prop:ch38-two-gradings}
Let $L\in\Pic(X)_{\mathbf Q}$ and let $\phi\in\cH^{\mathrm{gf}}_{\mathbf R}(L)$
be presented by a weighted section system
$(m,(s_j)_{1\leq j\leq r},(\lambda_j)_{1\leq j\leq r})$, so that
$\phi=\frac1m\max_j\{\log|s_j|+\lambda_j\}$ with
$s_j\in H^0(X,mL)$.
\begin{enumerate}
\item\label{itm:ch38-grad:weight}
 $\phi$ is homogeneous if and only if the weights may be taken all equal
 to $0$.  More precisely, if $\phi$ is homogeneous then, for
 \emph{every} presentation, $\max_j\lambda_j=0$ and
 $\phi=m^{-1}\max_j\log|s_j|$; so the weights may be deleted from any
 presentation without changing $\phi$, and
 $\phi=m^{-1}\log|\mathfrak b|$, where
 $\mathfrak b\subseteq\cO_X$ is the ideal locally generated by the
 $s_j$.
\item\label{itm:ch38-grad:degree}
 The degree $m$ is a Veronese ambiguity: for every $r\geq1$ one has
 $\phi=(rm)^{-1}\log|\mathfrak b^r|$, and $rmL\otimes\mathfrak b^r$ is
 globally generated whenever $mL\otimes\mathfrak b$ is.
\item\label{itm:ch38-grad:integral}
 Conversely, if $m^{-1}\log|\mathfrak b|=m'^{-1}\log|\mathfrak b'|$ with
 $\mathfrak b,\mathfrak b'\subseteq\cO_X$ nonzero, then
 $\overline{\mathfrak b^{\,m'}}=\overline{\mathfrak b'^{\,m}}$, the bars
 denoting integral closure.  Thus the pair $(m,\mathfrak b)$ is
 determined by $\phi$ only up to Veronese powers and integral closure.
\item\label{itm:ch38-grad:operations}
 On homogeneous data the two gradings behave differently under the two
 basic operations: with
 $\psi=m^{-1}\log|\mathfrak b|\in\cH_{\mathrm{hom}}(L)$ and
 $\psi'=m'^{-1}\log|\mathfrak b'|\in\cH_{\mathrm{hom}}(L')$,
\[
 \begin{aligned}
 \psi+\psi'&=(mm')^{-1}\log\bigl|\mathfrak b^{m'}\mathfrak b'^{\,m}\bigr|\\
 &\in\cH_{\mathrm{hom}}(L+L'),\\
 \max\{\psi,\psi'\}
 &=(mm')^{-1}\log\bigl|\mathfrak b^{m'}+\mathfrak b'^{\,m}\bigr|,
 \end{aligned}
\tag{38.16}\label{eq:ch38-two-gradings-operations}
\]
 the second belonging to $\cH_{\mathrm{hom}}(L)$ when $L=L'$.
 Degrees multiply; weights, which are all zero, do not enter.  For
 inhomogeneous data the weights add in the first formula and are compared
 in the second, and no formula expresses one grading in terms of the other.
\end{enumerate}
\end{proposition}

\begin{proof}
\emph{\cref{itm:ch38-grad:weight}.}  Suppose $\phi$ is homogeneous.  By
\Cref{prop:ch34-generic-finiteness}\,\cref{itm:ch34-gf:decreasing},
$\sup_{X^{\An}}\phi=\phi(v_{\triv})=m^{-1}\max_j\lambda_j$, and
$\phi(v_{\triv})=0$ by
\Cref{lem:ch35-homogeneity-reformulations}, $v_{\triv}$ being a fixed
point of the scaling action at which $\phi$ is finite; hence
$\max_j\lambda_j=0$ and $\lambda_j\leq0$ for every $j$.

Fix $v\in X^{\val}$ and set $c_j:=\log|s_j|(v)\in[-\infty,0]$, so that
$\phi(v)$ is finite by
\Cref{prop:ch34-generic-finiteness}\,\cref{itm:ch34-gf:equiv} and, by
\Cref{lem:ch34-log-s-properties}\,\cref{itm:ch34-logs:homogeneous},
\[
 m\,\phi(tv)=\max_j\{t\,c_j+\lambda_j\}=:f(t)
 \qquad(t>0).
\]
Homogeneity says $f(t)=t\,f(1)$ for every $t>0$.  For each $j$ with
$c_j>-\infty$ we have $t\,f(1)=f(t)\geq t\,c_j+\lambda_j$; dividing by
$t$ and letting $t\to+\infty$ gives $f(1)\geq c_j$, and this also
holds trivially when $c_j=-\infty$.  Hence
$f(1)\geq\max_jc_j$.  On the other hand $\lambda_j\leq0$ gives
$f(1)=\max_j\{c_j+\lambda_j\}\leq\max_jc_j$.  Therefore
$f(1)=\max_jc_j$, that is,
\[
 \phi(v)=\frac1m\max_j\log|s_j|(v)
 \qquad\text{for every }v\in X^{\val}.
\]
Both sides are continuous on $X^{\An}$ with values in
$[-\infty,+\infty)$, by
\Cref{prop:ch34-generic-finiteness}\,\cref{itm:ch34-gf:cont}, and
$X^{\divis}\subseteq X^{\val}$ is dense in $X^{\An}$ by
\Cref{thm:triv-divisorial-density}; so the identity holds on all of
$X^{\An}$.  Thus $(m,(s_j),(0))$ is a presentation of $\phi$ with all
weights zero.  The converse and the identification of $\mathfrak b$ are
\Cref{prop:hom-fs-ideal}.

\emph{\cref{itm:ch38-grad:degree}.}  By
\eqref{eq:ch34-log-ideal-semiring},
$\log|\mathfrak b^r|=r\log|\mathfrak b|$, so
$(rm)^{-1}\log|\mathfrak b^r|=m^{-1}\log|\mathfrak b|=\phi$.  If
$mL\otimes\mathfrak b$ is globally generated then so is its $r$-th
tensor power, and the natural surjection
$(mL\otimes\mathfrak b)^{\otimes r}\twoheadrightarrow
rmL\otimes\mathfrak b^{r}$ shows that the target is globally generated as
well.

\emph{\cref{itm:ch38-grad:integral}.}  The hypothesis gives
$m'\log|\mathfrak b|=m\log|\mathfrak b'|$, that is,
$v(\mathfrak b^{m'})=v(\mathfrak b'^{\,m})$ for every
$v\in X^{\val}$, and in particular for every divisorial valuation.  By
the valuative criterion of integral dependence
\Cref{thm:appA-valuative-integral-closure}, two ideals with the same
divisorial valuations have the same integral closure.

\emph{\cref{itm:ch38-grad:operations}.}  The two displayed identities are
\eqref{eq:ch34-log-ideal-semiring} after passing to the common degree
$mm'$, exactly as in the proof of
\Cref{prop:ch35-hom-pl-semiring}\,(1); the membership statements are
\Cref{prop:hom-fs-ideal}, using that a product and a sum of globally
generated sheaves twisted by the corresponding bundles are globally
generated.  The last sentence is the content of
\Cref{lem:gf-fs-operations}, where the weights are carried explicitly.
\end{proof}

\begin{warning}{Degree and weight are not interchangeable}
\label{warn:ch38-two-gradings}
In the Rees-algebra picture \eqref{eq:filtration-rees} the first grading is
the degree $m$ of the section ring $R(X,L)=\bigoplus_mR_m$, and the
second is the filtration index $\lambda$, which is the
$\mathbf G_m$-weight of the test configuration and the exponent of
$\varpi$.  A base ideal carries both indices,
$\mathfrak b_{m,\lambda}$, and the flag ideal of
\eqref{eq:ch33-degree-m-flag} assembles the second index at a fixed value
of the first.  Passing to a Veronese multiple replaces the degree $m$ and an
integral filtration index $\lambda$ by $rm$ and $r\lambda$;
the normalized quantity $\lambda/m$ is unchanged. In the
constructions used here, twisting the degree-$m$ series by a fixed
bundle $A$ changes the numerical class by $c_1(A)/m$.
A common bounded shift of all weights changes the function by
that shift divided by $m$; separate bounded weight changes give
an error bound of the same order, rather than necessarily a single
constant. Neither error is caused by the degree index alone.
The class error is treated by \Cref{prop:ch38-drift-ledger}, while
common constant shifts are tracked by normalization.
In the homogeneous category the second
grading is identically zero, by
\Cref{prop:ch38-two-gradings}\,\cref{itm:ch38-grad:weight}; this is exactly
why homogeneous data can be transported along birational morphisms without
an error term, as \Cref{chap:bj-reflection} exploits.
\end{warning}

\subsection{A homogeneous Fubini--Study function, realized and read back}
\label{sec:ch38-fs-example}

\begin{example}[The dictionary computed in both directions]
\label{ex:ch38-fs-from-bdivisor-and-back}
\index{Fubini--Study function!homogeneous, computed}
We continue \Cref{ex:ch38-blowup-surface-evaluation} with $d=3$, so
$X=\mathbf P^2$, $L=\cO(3)$, $\theta=c_1(L)$, and $B$ is the
Cartier b-divisor determined on $Y_2$ by
$B_{Y_2}=-2E_1'-3E_2$.  We check
\Cref{prop:ch38-model-fs-dictionary} in both directions with explicit
sections and an explicit ideal.

\emph{From the b-divisor to the ideal.}  Take $m=1$.  The extension of
the ideal $\mathfrak b=(x^3,xy,y^2)$ of
\eqref{eq:ch38-surface-ideal} to $Y_2$ is computed chart by chart.  In
the chart $(z_1,z_2)$ with $x=z_1z_2$, $y=z_1z_2^2$,
\[
 \mathfrak b\cdot\cO_{Y_2}
 =(z_1^3z_2^3,\;z_1^2z_2^3,\;z_1^2z_2^4)
 =z_1^2z_2^3\,(z_1,\;1,\;z_2)
 =(z_1^2z_2^3),
\]
which is $\cO(-2E_1'-3E_2)$ there.  In the chart of $Y_1$ with
$x=uy$, on which $\pi_2$ is an isomorphism and
$E_1'=\{y=0\}$,
\[
 \mathfrak b\cdot\cO
 =(u^3y^3,\;uy^2,\;y^2)=y^2\,(u^3y,\;u,\;1)=(y^2)
 =\cO(-2E_1').
\]
In the remaining chart, with $x=w_1$ and $y=w_1^2w_2$, on which only
$E_2=\{w_1=0\}$ is visible,
\[
 \mathfrak b\cdot\cO
 =(w_1^3,\;w_1^3w_2,\;w_1^4w_2^2)=(w_1^3)=\cO(-3E_2).
\]
Away from $p$ the ideal is the unit ideal.  Hence
$\mathfrak b\cdot\cO_{Y_2}=\cO_{Y_2}(B_{Y_2})$ is invertible, and
$Y_2$ dominates the normalized blowup of $\mathfrak b$.

\emph{Global generation and membership in $\cH_{\mathrm{hom}}$.}  The
space $H^0(\mathbf P^2,\cO(3)\otimes\mathfrak b)$ consists of the cubic
forms whose dehomogenization at $z=1$ lies in $\mathfrak b$.  Among the
ten monomials of degree three, six qualify:
\[
 x^3,\quad x^2y,\quad xy^2,\quad y^3,\quad xyz,\quad y^2z,
\tag{38.17}\label{eq:ch38-surface-sections}
\]
and the four that do not are $x^2z$, $xz^2$, $yz^2$, $z^3$, whose
dehomogenizations $x^2,x,y,1$ are precisely the four monomials outside
$\mathfrak b$.  So $\mathfrak b$ imposes
$4=\operatorname{colength}_p\mathfrak b$ independent conditions and
$h^0(\cO(3)\otimes\mathfrak b)=6$.  These six sections generate: away
from $p$, if $x(\xi)\neq0$ then $x^3$ does not vanish at $\xi$ and
if $y(\xi)\neq0$ then $y^3$ does not, and $x=y=0$ defines
$p$ alone; at $p$, the dehomogenizations generate
$(x^3,x^2y,xy^2,y^3,xy,y^2)=(x^3,xy,y^2)=\mathfrak b$.  Hence
$\cO(3)\otimes\mathfrak b$ is globally generated, and
\Cref{prop:hom-fs-ideal} gives
\[
 \begin{aligned}
 \psi_B&=\log|\mathfrak b|\\
 &=\max\!\left\{\begin{gathered}
   \log|x^3|,\ \log|x^2y|,\ \log|xy^2|,\\
   \log|y^3|,\ \log|xyz|,\ \log|y^2z|
   \end{gathered}\right\}\\
 &\in\cH_{\mathrm{hom}}(\cO(3)).
 \end{aligned}
\tag{38.18}\label{eq:ch38-surface-fs-presentation}
\]
All six weights are zero, in accordance with
\Cref{prop:ch38-two-gradings}\,\cref{itm:ch38-grad:weight}.

\emph{From the ideal back to the b-divisor.}  Running
\Cref{prop:ch38-model-fs-dictionary}\,\cref{itm:ch38-dict:backward}
backwards: the normalized blowup $\rho:Z\to X$ of $\mathfrak b$ is
dominated by $Y_2$, and on $Y_2$ the exceptional divisor of
$\mathfrak b$ is $2E_1'+3E_2$; so the Cartier b-divisor produced by
the backward direction is determined on $Y_2$ by
$-2E_1'-3E_2=B_{Y_2}$.  The round trip closes, as
\Cref{prop:ch38-model-fs-dictionary}\,\cref{itm:ch38-dict:unique}
guarantees.

\emph{Consistency with the numerical criterion.}  By
\eqref{eq:ch38-surface-nef-test} with $d=3$, the total class
$P=3H-2e_1-e_2$ is nef with $P\cdot\Lambda=0$, so it is not ample.  It
is nevertheless semiample: $\pi^*(\cO(3))\otimes(\mathfrak b\cdot\cO_{Y_2})
=\pi^*(\cO(3)\otimes\mathfrak b)$ is globally generated, being the
pullback of a globally generated sheaf.  This is exactly the right-hand
side of \eqref{eq:ch35-semiample-criterion}, and
\Cref{thm:ch35-semiample-criterion} therefore certifies
\eqref{eq:ch38-surface-fs-presentation} a second time, without computing
sections.

\emph{Values.}  Collecting
\eqref{eq:ch38-surface-divisorial-values},
\eqref{eq:ch38-surface-quasimonomial-value} and
\eqref{eq:ch38-surface-ideal}:
\[
 \begin{aligned}
 \psi_B(v_{\triv})&=0,
 &\psi_B(s\ord_{E_1})&=-2s,\\
 \psi_B(s\ord_{E_2})&=-3s,
 &\psi_B(v_\alpha)&=-2\alpha_1-3\alpha_2,
 \end{aligned}
\]
$\psi_B(\ord_C)=0$ for every curve $C\subset\mathbf P^2$, and
$\psi_B(v_{p,\triv})=-v_{p,\triv}(\mathfrak b)=-\infty$, because
$\mathfrak b\subseteq\mathfrak m_p$.  So $\psi_B$ is a homogeneous
$\theta$-psh function which is finite on $X^{\val}$, unbounded along
two of the exceptional rays, and $-\infty$ at exactly one point of
proper support, namely $v_{p,\triv}$; the polar locus is
$V(\mathfrak b)^{\An}=\{p\}^{\An}$, by
\Cref{prop:ch34-generic-finiteness}\,\cref{itm:ch34-gf:polar}.

\emph{What changes for $d=4$.}  Then $P=4H-2e_1-e_2$ has
$P\cdot\Lambda=1>0$, so $P$ is ample and the whole construction of
\Cref{sec:ch38-buffer} is unnecessary: the b-divisor is already realized by
a Fubini--Study function in the class $c_1(\cO(4))$.  The case $d=3$ is
the boundary case, and the case $d=2$ is excluded, $P$ not being nef.
\end{example}

\begin{remark}[What the example does not show]
\label{rem:ch38-example-limits}
In \Cref{ex:ch38-fs-from-bdivisor-and-back} the b-divisor is rational and
its total class is semiample, so the dictionary
\Cref{prop:ch38-model-fs-dictionary} applies directly and no limit is
needed. The direct dictionary for the specified data does not apply
when one of these features fails: when the coefficients of $B$ are
irrational, so that no single ideal computes $\psi_B$; or when the actual total
line bundle is nef and not semiample, so that no globally
generated multiple is available.
\Cref{ex:ch38-drift-irrational} exhibits the first failure on the same
surface, and a nef non-semiample rational line bundle occurs, for example, as
a non-torsion degree-zero line bundle on a smooth projective curve
over $\mathbf C$. Its numerical class is zero, but no positive
tensor power has a nonzero section. This illustrates why the
dictionary concerns an actual line bundle, not just its numerical class.
\end{remark}

\section{Passage between numerical classes on birational models}
\label{sec:ch38-classes}

Three numerical spaces occur in the construction: $N^1(X)$, where the
target class $\theta$ lives; $N^1(Y)$ for the determination
$\pi:Y\to X$, where the total class $P$ lives; and $N^1(Z)$ for the
higher models on which individual divisorial valuations are evaluated.
This section records how positivity moves between them, and it isolates the
one place where the passage is not free of charge: the numerical class of
the approximating stages differs from the target by an error that must be
tracked and, above all, \emph{oriented}.

\subsection{Pullback, pushforward, and the total class}
\label{sec:ch38-pullback-pushforward}

\begin{lemma}
\label{lem:ch38-numerical-birational}
Let $\pi:Y\to X$ be a projective birational morphism of normal
projective varieties of dimension $n$.
\begin{enumerate}
\item\label{itm:ch38-num:inj}
 $\pi^*:N^1(X)\to N^1(Y)$ is injective, and
 $\pi^*\Nef(X)\subseteq\Nef(Y)$. No pushforward
 $N^1(Y)\to N^1(X)$ is asserted for a general normal target.
\item\label{itm:ch38-num:notample}
 If $\pi$ is not an isomorphism, then $\pi^*\alpha$ is \emph{never}
 ample on $Y$, for any $\alpha\in N^1(X)$.
\item\label{itm:ch38-num:relative}
 Let $G$ be an effective Cartier divisor on $Y$ with $-G$
 $\pi$-ample, and let $\omega\in N^1(X)$ be ample.  Then there is
 $r_1\geq1$ such that
\[
 r\pi^*\omega-G\ \text{ is ample on }Y
 \qquad\text{for every real }r\geq r_1 .
\tag{38.19}\label{eq:ch38-relative-to-absolute}
\]
\item\label{itm:ch38-num:total}
 For $B\leq0$ determined on $Y$ and $\theta\in N^1(X)$, the total
 class $P=\pi^*\theta+B_Y$ of \eqref{eq:bj-nef-total}
 is $\pi$-nef if and only if $B_Y$ is. If an
 $\mathbf R$-Cartier representative $\Theta$ of $\theta$ is chosen,
 then $\pi_*(\pi^*\Theta+B_Y)=\Theta+B_X$ as Weil divisors;
 $B_X$ need not itself be $\mathbf R$-Cartier.
\end{enumerate}
\end{lemma}

\begin{proof}
\emph{\cref{itm:ch38-num:inj}.} For every integral curve
$\Gamma\subset X$, properness and surjectivity provide an integral
curve $C\subset Y$ with $\pi_*C=d\Gamma$, $d>0$, by cutting a
component of $\pi^{-1}(\Gamma)$ dominating $\Gamma$ with general
very ample divisors. Thus
$d(\alpha\cdot\Gamma)=\pi^*\alpha\cdot C$.
If $\pi^*\alpha=0$, all these degrees vanish, so $\alpha=0$.
For nef pullback, the projection formula gives
$\pi^*\alpha\cdot C=\alpha\cdot\pi_*C\geq0$ for a nef $\alpha$;
the degree is zero on contracted curves.

\emph{\cref{itm:ch38-num:notample}.}  Suppose $\pi$ is not an
isomorphism.  If every fibre of $\pi$ were finite, then $\pi$, being
proper and quasi-finite, would be finite; a finite birational morphism onto
a normal variety is an isomorphism, by Zariski's main theorem
\cite[Corollary~III.11.4]{Har}, a contradiction.  Hence some fibre has an
irreducible component $T$ of positive dimension, and cutting $T$ with
$\dim T-1$ general very ample divisors produces an integral curve
$C\subset Y$ contracted by $\pi$, so that
$\pi_*C=0$ and $\pi^*\alpha\cdot C=\alpha\cdot\pi_*C=0$ by
\eqref{eq:appF-projection}.  An ample class has strictly positive degree on
every integral curve.

\emph{\cref{itm:ch38-num:relative}.}  Choose an ample line bundle $A$ on
$X$ and a rational $c>0$ with $\omega-c\,c_1(A)$ ample; this is
possible because $\Amp(X)$ is open and $\omega\in\Amp(X)$.  Since
$\cO_Y(-G)$ is $\pi$-ample and $A$ is ample on $X$, the twist
theorem \cite[4.6.13]{EGAII} provides $r_0\in\mathbf Z_{>0}$ with
$\pi^*A^{\otimes r_0}\otimes\cO_Y(-G)$ ample on $Y$, that is,
$r_0\pi^*c_1(A)-G\in\Amp(Y)$.  Put $r_1:=r_0/c$.  For $r\geq r_1$,
\[
 r\pi^*\omega-G
 =\bigl(r_0\pi^*c_1(A)-G\bigr)
  +\pi^*\bigl(r\omega-r_0c_1(A)\bigr),
\]
and $r\omega-r_0c_1(A)=r\bigl(\omega-(r_0/r)c_1(A)\bigr)$ is ample
because $r_0/r\leq c$ and $\omega-c\,c_1(A)$ is ample, so that
$\omega-(r_0/r)c_1(A)
=(\omega-c\,c_1(A))+(c-r_0/r)c_1(A)$ is a sum of an ample and a nef class.
Its pullback is nef by \cref{itm:ch38-num:inj}, and the sum of an ample and
a nef class is ample.

\emph{\cref{itm:ch38-num:total}.} The Weil-divisor equality uses
$\pi_*\pi^*\Theta=\Theta$ and $\pi_*B_Y=B_X$. It does not require
a numerical pushforward on Cartier classes. For every contracted
curve $C$, $\pi^*\theta\cdot C=0$, hence
$P\cdot C=B_Y\cdot C$.
\end{proof}

\begin{warning}{Ampleness cannot be imported from the base}
\label{warn:ch38-no-ample-from-base}
\Cref{lem:ch38-numerical-birational}\,\cref{itm:ch38-num:notample} is the
structural reason for the shape of the proof of
\Cref{thm:bj-real-cartier-realization}.  The dictionary
\Cref{prop:ch38-model-fs-dictionary} needs an \emph{ample}, or at least
semiample, total class on $Y$; and no amount of positivity on $X$
produces one, because every pullback class is trivial on the contracted
curves.  The missing positivity must therefore be supplied on $Y$, by a
$\pi$-ample correction, and \eqref{eq:ch38-relative-to-absolute} is the
statement that a $\pi$-ample correction plus enough pullback positivity
is ample.  This is the whole content of the ample buffer of
\Cref{sec:ch38-buffer}: the divisor $G$ supplies relative positivity, the
class $r\pi^*\omega$ supplies absolute positivity, and $t$ interpolates
between the nef class $P$ and the ample class $r\pi^*\omega-G$.
\end{warning}

\subsection{The drift ledger}
\label{sec:ch38-drift}

The approximating stages of the realization theorem are homogeneous
Fubini--Study functions for \emph{rational} line bundles, whereas the
target class $\theta$, and hence the activated class
$\omega=\theta+[F]$, is real.  Consequently the class of the stage
$n$ differs from $\omega$.  The next proposition is the exact
bookkeeping for that difference.  It says two things: that the error may be
made $O(n^{-1})$, and --- much more importantly --- that it may be made
\emph{negative} in the nef order, which is what allows the limit to be
taken inside a single class.

\begin{proposition}
\label{prop:ch38-drift-ledger}
Fix a norm $\|\cdot\|$ on the finite-dimensional real vector space
$N^1(X)$.
\begin{enumerate}
\item\label{itm:ch38-drift:dominated}
 \emph{Dominated drift is harmless.}  Let $(\phi_i)_{i\in I}$ be a
 decreasing net with $\phi_i\in\PSH_X(\theta_i)$, and suppose that
 $\theta-\theta_i$ is nef for every $i$.  If the pointwise limit
 $\phi:=\lim_i\phi_i$ is not identically $-\infty$, then
 $\phi\in\PSH_X(\theta)$.
\item\label{itm:ch38-drift:orient}
 \emph{Drift can be oriented.}  Let $\omega\in N^1(X)$ be ample and
 $\delta>0$.  Then
\[
 \bigl\{\omega'\in N^1(X)_{\mathbf Q}\ :\
 \omega-\omega'\ \text{ample},\ \|\omega-\omega'\|<\delta\bigr\}
 \neq\varnothing,
\tag{38.20}\label{eq:ch38-drift-orientable}
\]
 and the set is the intersection of a nonempty open subset of
 $N^1(X)$ with the dense subgroup $N^1(X)_{\mathbf Q}$.
\item\label{itm:ch38-drift:joint}
 \emph{Orientation is compatible with an ampleness constraint upstairs.}
 Let $\pi:Y\to X$ be a model, let $\omega\in N^1(X)$ be ample, and let
 $E$ be an $\mathbf R$-Cartier divisor on $Y$ with
 $\pi^*\omega-E$ ample.  Then for every $\delta>0$ there are
 $\omega'\in N^1(X)_{\mathbf Q}$ and a rational Cartier divisor $R$ on
 $Y$ supported on $\Supp E$ with
\[
 \|\omega-\omega'\|<\delta,\qquad
 \|E-R\|<\delta,\qquad
 \omega-\omega'\ \text{ample},\qquad
 \pi^*\omega'-R\ \text{ample}.
\tag{38.21}\label{eq:ch38-drift-joint}
\]
\item\label{itm:ch38-drift:rate}
 \emph{The rate.}  If in \cref{itm:ch38-drift:joint} one applies the
 statement with $\delta=1/n$ for $n\geq1$, the resulting classes
 satisfy $\|\omega-\omega_n\|<1/n$; the drift is $O(n^{-1})$, and the
 induced error in the total class on $Y$ is
 $\|\pi^*(\omega-\omega_n)\|=O(n^{-1})$ as well.
\end{enumerate}
\end{proposition}

\begin{proof}
\emph{\cref{itm:ch38-drift:dominated}.}  Since $\theta-\theta_i$ is nef,
\Cref{lem:ch34-psh-nonempty} gives $0\in\PSH_X(\theta-\theta_i)$, and
\Cref{thm:bj-basic-closure}\,(3) with $s=t=1$ gives
\[
 \phi_i=\phi_i+0\in\PSH_X\bigl(\theta_i+(\theta-\theta_i)\bigr)
 =\PSH_X(\theta).
\]
The net $(\phi_i)$ is thus a decreasing net in the single class
$\PSH_X(\theta)$, and \Cref{thm:bj-basic-closure}\,(4) applies, its
hypothesis being exactly that the limit is not identically $-\infty$ on
an irreducible component; here $X$ is irreducible.

\emph{\cref{itm:ch38-drift:orient}.}  Fix an ample class $\alpha$ on
$X$; such a class exists since $X$ is projective.  For
$0<\varepsilon<\delta/\|\alpha\|$ the class
$\omega-\varepsilon\alpha$ satisfies
$\|\omega-(\omega-\varepsilon\alpha)\|=\varepsilon\|\alpha\|<\delta$ and
$\omega-(\omega-\varepsilon\alpha)=\varepsilon\alpha$ is ample.  The two
conditions ``$\omega-\omega'$ ample'' and
``$\|\omega-\omega'\|<\delta$'' are open in $\omega'$, the first
because $\Amp(X)$ is open in $N^1(X)$, so the displayed set is the
intersection of a nonempty open set with
$N^1(X)_{\mathbf Q}$.  Rational classes are dense in $N^1(X)$ because
$N^1(X)=N^1(X)_{\mathbf Q}\otimes_{\mathbf Q}\mathbf R$ by construction.

\emph{\cref{itm:ch38-drift:joint}.}  Let
$V\subseteq\Div(Y)_{\mathbf R}$ be the finite-dimensional space of
$\mathbf R$-Cartier divisors supported on the finite set $S$ of
components of $\Supp E$; it is defined over $\mathbf Q$ by
\Cref{lem:ch38-fixed-support-rational}.  Consider the
map
\[
 \Phi:N^1(X)\times V\longrightarrow N^1(Y),
 \qquad
 \Phi(\omega',R):=\pi^*\omega'-[R],
\]
which is $\mathbf R$-linear, hence continuous.  By hypothesis
$\Phi(\omega,E)\in\Amp(Y)$, an open cone, so there is
$\eta>0$ such that $\Phi(\omega',R)$ is ample whenever
$\|\omega-\omega'\|+\|E-R\|<\eta$.  Shrinking $\eta$ below $\delta$,
and intersecting with the nonempty open set produced in
\cref{itm:ch38-drift:orient}, we obtain a nonempty open subset of
$N^1(X)\times V$ on which all four conditions of
\eqref{eq:ch38-drift-joint} hold; it contains a rational point because
$N^1(X)_{\mathbf Q}\times V_{\mathbf Q}$ is dense.

\emph{\cref{itm:ch38-drift:rate}} is immediate, using that
$\pi^*$ is a fixed linear map between finite-dimensional spaces and
hence bounded.
\end{proof}

\begin{warning}{The drift is not absorbed by its size}
\label{warn:ch38-drift-not-absorbed}
It is tempting to argue that since $\omega_n\to\omega$ the classes are
``eventually $\omega$ up to $O(n^{-1})$'' and the limit is therefore
$\omega$-psh.  This is not an argument, and the reason is worth stating
precisely.  The closure calculus available before the envelope theorem
does contain a clause allowing the class to move:
\Cref{thm:bj-basic-closure}\,(5) passes a decreasing net
$\psi_j\in\PSH_X(\theta_j)$ with $\theta_j\to\theta$ to the limit,
provided the limit is not identically $-\infty$.  Class convergence alone does not prove that the limit is
nontrivial. Here that condition is checked separately at
$v_{\triv}$, where every stage is zero. Together with this check,
the moving-class clause would be sufficient.
We use the stronger dominated construction of
\Cref{prop:ch38-drift-ledger}\,\cref{itm:ch38-drift:dominated}:
the condition $\omega-\omega_n\in\Nef(X)$ puts every stage in
one fixed psh class before taking the decreasing limit.
Smallness of the drift identifies the limiting numerical class;
its orientation provides this convenient fixed-class proof.

The same discipline is applied in
\Cref{chap:bj-smooth-characteristic}, where the stages carry classes
$L+A/m$ and the numerical convergence and nontriviality checks
are stated separately, and in \Cref{chap:bj-reflection}, where the exact
drift-cancellation identity \eqref{eq:bj-reflection-total} is used instead
of an estimate.  A statement of the form ``the class error tends to zero,
so it may be ignored'' does not occur anywhere in
\Crefrange{chap:bj-cartier-realization}{chap:bj-envelope-consequences}.
\end{warning}

\begin{remark}[Where the two kinds of $O(m^{-1})$ error occur]
\label{rem:ch38-two-errors}
Two $O(m^{-1})$ errors are audited in this part of the book, and they are
different.  The first is the error in \emph{numerical class} produced when
a real class is replaced by a rational one, or when a degree-$m$
linear series is twisted by a fixed bundle $A$ to make it globally
generated; this is the error of \Cref{prop:ch38-drift-ledger} and
the moving-class regularization in \Cref{thm:ch34-fs-regularization}.
Real weights are allowed independently of this numerical twist.  The second is the error in
the \emph{value} of a function produced when a weight is divided by $m$;
a common bounded weight shift gives an additive constant
$O(m^{-1})$, tracked by \eqref{eq:normalized-psh-space}.  In the homogeneous category
the second error vanishes identically, by
\Cref{prop:ch38-two-gradings}\,\cref{itm:ch38-grad:weight}, which is one of
the reasons the envelope proof passes to homogeneous data before doing
anything delicate.  \Cref{warn:ch38-two-gradings} records the same
distinction from the algebraic side.
\end{remark}

\begin{example}[A drift that cannot be removed]
\label{ex:ch38-drift-irrational}
Return to \Cref{ex:ch38-blowup-surface-evaluation} and take
$\theta:=\sqrt2\,c_1(\cO(3))$, an irrational ample class on
$X=\mathbf P^2$, and $B:=-\sqrt2\,(2\overline{E_1'}+3\overline{E_2})$
in the obvious notation, that is, the real Cartier b-divisor determined on
$Y_2$ by $\sqrt2\,B_{Y_2}^{(1)}$ where
$B_{Y_2}^{(1)}=-2E_1'-3E_2$.  The total class is
$\sqrt2\,P$ with $P$ as in \eqref{eq:ch38-surface-total} for
$d=3$, hence nef.  No rational class equals $\theta$, and no rational
Cartier b-divisor equals $B$; every approximating stage therefore lives
in a class $\omega_n\neq\theta$.  There is no way to arrange
$\omega_n=\theta$ for a single $n$, so the drift is genuinely present
and must be handled, not avoided.

What one \emph{can} arrange, by
\Cref{prop:ch38-drift-ledger}\,\cref{itm:ch38-drift:orient}, is
$\theta-\omega_n$ ample: for instance
$\omega_n:=q_n\,c_1(\cO(3))$ with $q_n\in\mathbf Q$ and
$q_n<\sqrt2$, $q_n\to\sqrt2$, so that
$\theta-\omega_n=(\sqrt2-q_n)c_1(\cO(3))$ is ample.  Then every stage in
the class $\omega_n$ is automatically a stage in the class $\theta$,
by \Cref{thm:bj-basic-closure}\,(3), and
\Cref{prop:ch38-drift-ledger}\,\cref{itm:ch38-drift:dominated} applies.
Numerically, with $q_n=\lfloor n\sqrt2\rfloor/n$ one has
$0<\sqrt2-q_n<1/n$, so the drift is $O(n^{-1})$ and is oriented. Convergence identifies
the target class; orientation puts the stages in its psh space.
\end{example}

\section{Ordered rational approximation}
\label{sec:ch38-approximation}

The real realization theorem requires a decreasing psh system.  For
coefficient functions, decreasing order corresponds to an increasing
sequence of effective divisors.  The following finite-dimensional lemma
provides that order.

Before stating it we record the linear-algebraic fact on which it rests,
since it is the only reason the approximation can be performed with
rational data at all.

\begin{lemma}
\label{lem:ch38-fixed-support-rational}
Let $W$ be a normal projective variety and let
$S=\{E_1,\ldots,E_s\}$ be a finite set of prime divisors on $W$.  Let
\[
 V:=\Bigl\{\textstyle\sum_{j=1}^sa_jE_j\ :\ a_j\in\mathbf R\Bigr\}
 \cap\operatorname{CDiv}(W)_{\mathbf R}
\]
be the space of $\mathbf R$-Cartier divisors supported on $S$.  Then
$V$ is a finite-dimensional real vector space defined over
$\mathbf Q$: setting $V_{\mathbf Q}:=V\cap\bigl(\sum_j\mathbf QE_j\bigr)$,
one has $V=V_{\mathbf Q}\otimes_{\mathbf Q}\mathbf R$, and
$V_{\mathbf Q}$ is dense in $V$.  Moreover the coefficientwise partial
order on $V$ is defined by the $\mathbf Q$-rational linear forms
$D\mapsto\operatorname{coeff}_{E_j}(D)$, so that every strict
coefficientwise inequality defines an open subset of $V$.
\end{lemma}

\begin{proof}
Inside the free abelian group $\Div(W)_{\mathbf Z}$ on the prime divisors
of $W$, let $M_1$ be the subgroup generated by $S$, a direct summand
of rank $s$, and let $M_2:=\operatorname{CDiv}(W)_{\mathbf Z}$ be the
subgroup of Cartier divisors.  Both $M_1\otimes\mathbf Q$ and
$M_2\otimes\mathbf Q$ are $\mathbf Q$-subspaces of
$\Div(W)_{\mathbf Q}$, and $V=(M_1\otimes\mathbf R)\cap
(M_2\otimes\mathbf R)$.  Since $\mathbf R$ is flat over $\mathbf Q$,
the intersection of two $\mathbf Q$-subspaces commutes with the base
change $\mathbf Q\to\mathbf R$:
\[
 (M_1\otimes\mathbf R)\cap(M_2\otimes\mathbf R)
 =\bigl((M_1\otimes\mathbf Q)\cap(M_2\otimes\mathbf Q)\bigr)
   \otimes_{\mathbf Q}\mathbf R .
\]
Thus $V=V_{\mathbf Q}\otimes\mathbf R$ with
$V_{\mathbf Q}=(M_1\cap M_2)\otimes\mathbf Q$, which is dense in
$V$.  Finiteness of the dimension is $\dim V\leq s$.  The coefficient
forms are the restrictions to $V$ of the coordinate forms on
$M_1\otimes\mathbf R$, which are rational; and a strict inequality
between continuous linear forms defines an open set.
\end{proof}

\begin{lemma}\label{lem:bj-ordered-rational-approximation}
Let $\rho:W\to X$ be a projective morphism of normal projective
varieties.  Let $E\ge0$ be an $\mathbf R$-Cartier divisor on $W$.
Suppose that, for all sufficiently small $t>0$, there are
$\mathbf R$-Cartier divisors $E_t$ and classes
$\Omega_t\in N^1(X)_{\mathbf R}$ such that
\[
 0<E_t<E,\qquad
 E_t\uparrow E\quad(t\downarrow0),\qquad
 \Omega_t\longrightarrow\omega,
\tag{38.22}\label{eq:bj-real-buffer-hypotheses}
\]
and $\rho^*\Omega_t-E_t$ is ample.  Assume that all $E_t$ and $E$
have a fixed finite support $S$.  Then there are rational Cartier
divisors $R_n$ and rational classes
$\omega_n\in N^1(X)_{\mathbf Q}$ satisfying
\[
 \begin{gathered}
 0\le R_n<R_{n+1}<E,\qquad
 R_n\longrightarrow E,\qquad
 \omega_n\longrightarrow\omega,\\
                       \rho^*\omega_n-R_n\quad\text{ample}.
 \end{gathered}
\tag{38.23}\label{eq:bj-ordered-rational-approximation}
\]
\end{lemma}

\begin{proof}
Let $V\subset\mathbf R^S$ be the vector space of
$\mathbf R$-Cartier divisors supported on $S$.
By \Cref{lem:ch38-fixed-support-rational}, $V$ is defined
over $\mathbf Q$.  Rational points are therefore dense in
\[
                       V\times N^1(X)_{\mathbf R}.
\]

Set $R_0=0$.  Inductively choose $t_n>0$ so small that
\[
 R_{n-1}<E_{t_n}<E,\qquad
 \|E_{t_n}-E\|+\|\Omega_{t_n}-\omega\|<\frac1n.
\tag{38.24}\label{eq:bj-choice-tn}
\]
The strict coefficient inequalities and ampleness are open.  We may hence
choose a rational pair $(R_n,\omega_n)$ sufficiently close to
$(E_{t_n},\Omega_{t_n})$ that
\[
 R_{n-1}<R_n<E,\qquad
 \rho^*\omega_n-R_n\ \text{is ample},
\]
and its distance from $(E_{t_n},\Omega_{t_n})$ is less than $1/n$.
This proves all assertions.

We supply the details of the two inductive steps, since the ordering is
what the whole construction is for.

\emph{Rationality of the ambient space.}  This is
\Cref{lem:ch38-fixed-support-rational}: the space $V$ of
$\mathbf R$-Cartier divisors supported on $S$ is finite dimensional and
defined over $\mathbf Q$, and $N^1(X)_{\mathbf R}$ is defined over
$\mathbf Q$ by construction, so $V_{\mathbf Q}\times N^1(X)_{\mathbf Q}$
is dense in $V\times N^1(X)_{\mathbf R}$.  Fix norms on both factors.

\emph{Choice of $t_n$.}  Suppose $R_{n-1}\in V_{\mathbf Q}$ has been
constructed with $R_{n-1}<E$ coefficientwise and strictly on
$\Supp E$; for $n=1$ take $R_0=0$. The strict hypotheses
$0<E_t<E$ already require $E\neq0$; the zero divisor is not a
case of the strict conclusion.  Since $E_t\uparrow E$ as $t\downarrow0$ inside the
finite-dimensional space $V$, the coefficients of $E_t$ converge to
those of $E$, and $R_{n-1}<E$ strictly, so
$R_{n-1}<E_{t}$ for all sufficiently small $t$.  Together with
$\Omega_t\to\omega$ this gives \eqref{eq:bj-choice-tn} for a suitable
$t_n$, which we may also take $\leq t_{n-1}$.

\emph{Choice of the rational pair.}  The conditions
\[
 R_{n-1}<R<E,
 \qquad
  \rho^*\Omega-R\in\Amp(W),
 \qquad
 \|(R,\Omega)-(E_{t_n},\Omega_{t_n})\|<1/n
\]
on $(R,\Omega)\in V\times N^1(X)_{\mathbf R}$ define an open subset,
the first two by \Cref{lem:ch38-fixed-support-rational} and by openness of
the ample cone, and the third trivially.  The point
$(E_{t_n},\Omega_{t_n})$ satisfies the first two by
\eqref{eq:bj-choice-tn} and by hypothesis, and the third with strict
inequality $0<1/n$; so the open set is nonempty and contains a rational
point $(R_n,\omega_n)$.  Then
$\|E-R_n\|\leq\|E-E_{t_n}\|+\|E_{t_n}-R_n\|<2/n$ and likewise
$\|\omega-\omega_n\|<2/n$, which gives $R_n\to E$ and
$\omega_n\to\omega$.  The chain $R_{n-1}<R_n$ is part of the choice,
and $R_n>0$ holds because $R_n>R_{n-1}\geq0$ for $n\geq1$.
\end{proof}

\begin{remark}[Why the support is fixed]
\label{rem:ch38-fixed-support}
The proof works in one finite-dimensional rational vector space of Cartier
divisors.  It does not assert an ordered rational approximation for an
arbitrary moving-support family of Weil b-divisors.
\end{remark}

\begin{example}[A moving-support Weil b-divisor]
\label{ex:ch38-moving-support}
Let $X$ be a smooth projective surface and let
\[
 \cdots\longrightarrow Y_3\longrightarrow Y_2\longrightarrow Y_1
 \longrightarrow X
\]
be an infinite chain of point blowups, with $E_j$ the exceptional prime
introduced on $Y_j$. Choose the centres to be pairwise distinct
closed points $p_j$ of $X$, lifted through the preceding blowups.
Then the exceptional divisors $E_j$ are pairwise disjoint. Let
$B$ be the Weil b-divisor with $\operatorname{coeff}_{E_j}(B)=-2^{-j}$
for every $j$ and all other coefficients zero.  Its coefficient function
$\psi_B$ is homogeneous and locally finite, so
\Cref{prop:ch35-weil-coefficient-bijection} applies; but $B$ is not
Cartier. Indeed, its trace on $X$ is zero. A determination
$Y\to X$ would therefore have only exceptional support, whose image
is a finite set of points on this surface. Choose $p_j$ outside
that set and the non-isomorphism image of $Y\to X$.
The determination is zero near $p_j$, so its pullback coefficient
at $E_j$ would be zero, contradicting $-2^{-j}$.
Thus no fixed determination and finite support $S$ supplies the
ambient space of \Cref{lem:bj-ordered-rational-approximation}
for this $B$. The lemma does not apply to it.  Nor does
\Cref{thm:bj-real-cartier-realization}: its hypothesis is that $B$ be
determined on a single model.  This is the precise sense in which the
theorem is a statement about Cartier and not about Weil b-divisors, and
\Cref{rem:ch38-cartier-not-weil} explains why the Cartier statement is
nonetheless the one the envelope proof needs.
\end{example}

\subsection{The dominated form}
\label{sec:ch38-approximation-dominated}

The version of \Cref{lem:bj-ordered-rational-approximation} that is
actually applied in \Cref{sec:ch38-realization} carries the extra
conclusion isolated in
\Cref{prop:ch38-drift-ledger}\,\cref{itm:ch38-drift:orient}: the
approximating classes are dominated by the target class in the nef order.
By \Cref{warn:ch38-drift-not-absorbed}, this permits a fixed-class
limit proof, while convergence and nontriviality would also suffice
for the moving-class clause.

\begin{lemma}
\label{lem:ch38-ordered-rational-dominated}
Let $\rho:W\to X$ be a projective morphism of normal projective
varieties, let $\omega\in N^1(X)$ be \emph{ample}, and let
$E\geq0$ be an $\mathbf R$-Cartier divisor on $W$ with finite
support $S$.  Suppose that, for all sufficiently small $t>0$, there is
an $\mathbf R$-Cartier divisor $E_t$ supported on $S$ with
\[
 0<E_t<E,\qquad
 E_t\uparrow E\quad(t\downarrow0),\qquad
 \rho^*\omega-E_t\ \text{ample}.
\tag{38.25}\label{eq:ch38-dominated-hypotheses}
\]
Then there are rational Cartier divisors $R_n$ supported on $S$ and
rational classes $\omega_n\in N^1(X)_{\mathbf Q}$ with
\[
 \begin{gathered}
 0\leq R_n<R_{n+1}<E,\qquad
 \|E-R_n\|+\|\omega-\omega_n\|<\tfrac2n,\\
 \omega-\omega_n\ \text{ample},\qquad
 \rho^*\omega_n-R_n\ \text{ample}.
 \end{gathered}
\tag{38.26}\label{eq:ch38-dominated-conclusion}
\]
In particular $R_n\uparrow E$, $\omega_n\to\omega$, and every class
$\omega_n$ is dominated by $\omega$ in the nef order.
\end{lemma}

\begin{proof}
Apply \Cref{lem:bj-ordered-rational-approximation} with the constant path
$\Omega_t:=\omega$, whose hypotheses \eqref{eq:bj-real-buffer-hypotheses}
are exactly \eqref{eq:ch38-dominated-hypotheses}, and modify the choice of
the rational pair at each stage as follows.  In the notation of that proof,
the point $(E_{t_n},\omega)$ lies in the open subset of
$V\times N^1(X)$ defined by
\[
 R_{n-1}<R<E,\qquad
 \rho^*\Omega-R\in\Amp(W),\qquad
 \|(R,\Omega)-(E_{t_n},\omega)\|<1/n .
\]
By \Cref{prop:ch38-drift-ledger}\,\cref{itm:ch38-drift:orient} the set
$\{\Omega:\omega-\Omega\in\Amp(X)\}$ is open and has $\omega$ in its
closure, so its intersection with the displayed open set is a nonempty
open subset of $V\times N^1(X)$; it therefore contains a rational point
$(R_n,\omega_n)$, by \Cref{lem:ch38-fixed-support-rational} and density
of $N^1(X)_{\mathbf Q}$.  This is
\eqref{eq:ch38-dominated-conclusion}, with the estimate obtained as before
by the triangle inequality through $(E_{t_n},\omega)$, using
$\|E-E_{t_n}\|<1/n$ from \eqref{eq:bj-choice-tn}.
\end{proof}

\begin{remark}[Two forms, two uses]
\label{rem:ch38-two-approximation-forms}
\Cref{lem:bj-ordered-rational-approximation} allows the class to move along
a prescribed path $\Omega_t$; it is the form to use when the ample buffer
must be built from an auxiliary ample divisor on $X$, and it is the form
quoted in \Cref{chap:bj-reflection}, where the classes are corrected by an
exact identity rather than by domination.
\Cref{lem:ch38-ordered-rational-dominated} holds the class fixed along the
real path and moves only at the rationalization step; it is the form used
in \Cref{sec:ch38-realization}, because its output
satisfies the hypothesis of
\Cref{prop:ch38-drift-ledger}\,\cref{itm:ch38-drift:dominated}.  The two
are not interchangeable: the first requires
$\rho^*\Omega_t-E_t$ ample for a moving $\Omega_t$, which is easy to
arrange, while the second requires $\rho^*\omega-E_t$ ample for the fixed
$\omega$, which is exactly what the sharpened buffer of
\Cref{lem:ch38-buffer} delivers.
\end{remark}

\section{Activation: the fixed-pole construction}
\label{sec:ch38-activation}

The hypotheses of \Cref{thm:bj-real-cartier-realization} give a nef class
on $Y$ and nothing else.  Two deficiencies must be repaired before any
linear series is available: the class $\theta$ downstairs need not be
positive at all (\Cref{warn:ch38-theta-not-nef}), and the support of the
divisor $-B_Y$ need not meet the exceptional locus, so that no relative
positivity is visible along the contracted curves.  Both are repaired by
one device: a single effective divisor $F$ on $X$, chosen large enough
to make $\theta+[F]$ ample and to contain every locus that matters, is
added to the picture.  We call this \emph{activating} the whole support.
The price is a fixed pole, removed at the very end by homogeneous Siu
division.

The construction is quoted verbatim in \Cref{chap:bj-reflection}, where it
is called the fixed-pole construction, so we state it as a free-standing
lemma with all the properties that are used there.

\begin{lemma}
\label{lem:ch38-activation}
Let $X$ be a normal integral projective variety with
$\dim X>0$, let $\theta\in N^1(X)$, and let $B\leq0$ be a real
Cartier b-divisor determined on a model $\pi:Y\to X$, with
\[
                              D:=-B_Y\ge0.
\tag{38.27}\label{eq:bj-realization-D}
\]
Then, after replacing $Y$ by a model dominating it --- which replaces
$D$ by its pullback and does not change $B$, by
\Cref{lem:ch38-determination-equivalence} --- there exist an effective
$\mathbf Q$-Cartier divisor $G$ on $Y$ and an effective Cartier
divisor $F$ on $X$ with the following properties, where
$\omega:=\theta+[F]$, $B':=B-\overline F$ and $D':=D+\pi^*F$.  Items
\cref{itm:ch38-act:ample}--\cref{itm:ch38-act:dominated} require no further
hypothesis; items \cref{itm:ch38-act:total} and \cref{itm:ch38-act:pinef}
assume in addition that the total class $P:=\pi^*\theta+B_Y$ is nef, which
is not automatic --- on $X=\mathbf P^2$ with $\pi=\operatorname{id}$,
$B=0$ and $\theta=-c_1(\cO(1))$ every hypothesis above holds and $P$ is
not nef.  (\Cref{chap:bj-reflection} quotes only the unconditional items.)
\begin{enumerate}
\item\label{itm:ch38-act:ample} $[F]$ is ample and $\omega$ is ample;
\item\label{itm:ch38-act:relample} $-G$ is $\pi$-ample;
\item\label{itm:ch38-act:supports}
 $\Supp F$ contains the image of the non-isomorphism locus of $\pi$,
 the set $\pi(\Supp D)$ and the set $\pi(\Supp G)$; consequently
\[
 \Supp D'=\pi^{-1}(\Supp F)\supseteq\Supp G,
\]
 and $\Supp D'$ is a divisor with finitely many components;
\item\label{itm:ch38-act:dominated} $2G\leq\pi^*F$; in particular
 $R:=\pi^*F-G$ is effective and $R\geq\tfrac12\pi^*F$, so that
 $\pi^*F=G+R$ with $R\geq0$;
\item\label{itm:ch38-act:Qample} $\pi^*\omega-G$ is ample on $Y$;
\item\label{itm:ch38-act:total}
 $\pi^*\omega+B'_Y=\pi^*\theta+B_Y=P$, so that
 $\pi^*\omega-D'=P$ is nef and $B'$ is nef-total in the class
 $\omega$;
\item\label{itm:ch38-act:pinef} $-D'$ is $\pi$-nef, and
 $D'\geq\pi^*F>0$ on every component of $\Supp D'$.
\end{enumerate}
\end{lemma}

\begin{proof}
\emph{Step 1: a relatively ample exceptional divisor.}
By \Cref{thm:appB-admissible-domination}, applied to the projective
birational morphism $\pi:Y\to X$, which is an isomorphism over the
complement of a closed subset $Z_0\subsetneq X$, there is a blowup
$X'\to X$ with centre contained in $Z_0$ which factors through $Y$.
Replace $Y$ by the normalization of $X'$, which is again a model
dominating the original $Y$, and which is projective over $X$ since
normalization is finite; by
\Cref{lem:ch38-determination-equivalence}\,\cref{itm:ch38-de:family} this
does not change $B$, and it replaces $D$ by its pullback, which is
again effective.  On the new $Y$, the ideal $\mathcal I\subseteq\cO_X$
blown up satisfies $\mathcal I\cdot\cO_Y=\cO_Y(-G_0)$ with $G_0$ an
effective Cartier divisor and $-G_0$ relatively ample, by
\eqref{eq:appB-exceptional-ideal} and the construction of the normalized
blowup \eqref{eq:appB-normalized-blowup}.  Moreover
$\Supp G_0\subseteq\pi^{-1}(V(\mathcal I))$ with
$V(\mathcal I)\subseteq Z_0\subsetneq X$.

\emph{Step 2: the activation divisor.}
Set $Z:=Z_0\cup V(\mathcal I)\cup\pi(\Supp D)$, a closed subset of
$X$.  It is a proper closed subset: $Z_0$ and $V(\mathcal I)$ are,
by construction, and $\pi(\Supp D)$ is, because a prime divisor on $Y$
maps either onto a prime divisor of $X$ or into a subvariety of
codimension $\geq2$, $\pi$ being birational, so its image is nowhere
dense or a divisor, and in either case not all of $X$.  Let $A$ be a
very ample divisor on $X$.  Choose $m\in\mathbf Z_{>0}$ so large that
$\mathcal I_Z\otimes\cO_X(mA)$ is globally generated --- possible by
Serre's theorem --- and that $\theta+m\,c_1(A)$ is ample --- possible
because $\Amp(X)$ is open and contains $m\,c_1(A)+\theta$ for
$m\gg0$.  Choose a nonzero global section of
$\mathcal I_Z\otimes\cO_X(mA)$ and let $F\in|mA|$ be its divisor of
zeros; then $F$ is an effective Cartier divisor with
$\Supp F\supseteq Z$, $[F]=m\,c_1(A)$ is ample, and
$\omega=\theta+[F]$ is ample.  This proves
\cref{itm:ch38-act:ample} and the first half of
\cref{itm:ch38-act:supports}.

\emph{Step 3: normalizing $G$.}
By \Cref{lem:ch38-numerical-birational}\,\cref{itm:ch38-num:relative} there
is $r_1\geq1$ with $r\pi^*\omega-G_0$ ample for $r\geq r_1$, that is
$\pi^*\omega-r^{-1}G_0$ ample.  Since
$\Supp G_0\subseteq\pi^{-1}(V(\mathcal I))\subseteq\pi^{-1}(\Supp F)
=\Supp\pi^*F$ and $\pi^*F$ has strictly positive coefficient along every
component of its support, there is $r_2\geq1$ with
$2r_2^{-1}G_0\leq\pi^*F$.  Choose an integer
$r\geq\max\{r_1,r_2\}$ and put $G:=r^{-1}G_0$.  Then $G$ is an effective
$\mathbf Q$-Cartier divisor, $-G$ is $\pi$-ample --- relative
ampleness of an $\mathbf R$-divisor class is preserved by positive
scaling --- and \cref{itm:ch38-act:dominated} and
\cref{itm:ch38-act:Qample} hold.  This is \cref{itm:ch38-act:relample} as
well, and $\Supp G=\Supp G_0\subseteq\pi^{-1}(\Supp F)$ completes
\cref{itm:ch38-act:supports}, the equality
$\Supp D'=\pi^{-1}(\Supp F)$ holding because
$\Supp\pi^*F=\pi^{-1}(\Supp F)$ contains $\Supp D$.

\emph{Step 4: the total class.}
Since $\overline F$ is determined on $X$ by $F$, its trace on $Y$
is $\pi^*F$, so $B'_Y=B_Y-\pi^*F=-D'$ and
\[
 \pi^*\omega+B'_Y=\pi^*\theta+\pi^*F-D'=\pi^*\theta-D=\pi^*\theta+B_Y=P,
\]
which is \cref{itm:ch38-act:total}; $B'\leq B\leq0$ because
$\overline F\geq0$.

\emph{Step 5: relative nefness of $-D'$.}
Let $C\subset Y$ be an integral curve contracted by $\pi$.  Then
$\pi^*\theta\cdot C=0$ and $\pi^*F\cdot C=0$, so by
\cref{itm:ch38-act:total}
\[
 (-D')\cdot C=(P-\pi^*\omega)\cdot C=P\cdot C\geq0 ,
\]
the last inequality because $P$ is nef.  Hence $-D'$ is $\pi$-nef.
The inequality $D'\geq\pi^*F$ is $D\geq0$, and $\pi^*F$ has strictly
positive coefficients on $\Supp\pi^*F=\Supp D'$.
\end{proof}

\begin{remark}[What each hypothesis of \Cref{lem:ch38-activation} is for]
\label{rem:ch38-activation-uses}
Property \cref{itm:ch38-act:ample} is what makes
\Cref{prop:ch38-drift-ledger}\,\cref{itm:ch38-drift:orient} applicable, and
hence what allows the drift to be oriented; without an ample $\omega$
there is no room below the target class.  Property
\cref{itm:ch38-act:relample} is the only source of positivity along the
contracted curves, by \Cref{warn:ch38-no-ample-from-base}.  Property
\cref{itm:ch38-act:supports} is what makes the support of the whole family
$\{E_t\}$ constant, which is the hypothesis of
\Cref{lem:bj-ordered-rational-approximation} and, by
\Cref{ex:ch38-moving-support}, cannot be dispensed with.  Property
\cref{itm:ch38-act:dominated} is the normalization
$\pi^*F=G+R$ with $R\geq0$ quoted in
\Cref{chap:bj-reflection}; it also makes $r=1$ admissible in the buffer,
by \Cref{lem:ch38-buffer}.  Property \cref{itm:ch38-act:Qample} is the
sharpening that eliminates the class drift from the \emph{real} path, and
therefore the reason \Cref{lem:ch38-ordered-rational-dominated} applies
rather than only \Cref{lem:bj-ordered-rational-approximation}.  Property
\cref{itm:ch38-act:pinef} is used only through the buffer.
\end{remark}

\begin{warning}{The fixed pole is not removed by subtraction}
\label{warn:ch38-fixed-pole}
Activation replaces $B$ by $B'=B-\overline F$ and therefore
$\psi_B$ by
$\psi_{B'}=\psi_B+\log|s_F|$ on $X^{\val}$, hence in particular on
$X^{\divis}$.  This is an identity of finite real-valued coefficient
functions; it does not define either coefficient function at a
proper-support point.  At the end of the proof the pole must be removed,
and the
temptation is to define the answer as
$\chi_{B}:=\chi'-\log|s_F|$.  At a semivaluation supported in
$\Supp F$ both terms are $-\infty$ and the difference is undefined by
ledger entry~(4) of \Cref{sec:extended-arithmetic};
\Cref{ex:ch38-naive-realization-fails} shows that no convention repairs
this, because the correct value at two such points can be $0$ at one and
$-\infty$ at the other.  The removal is therefore carried out by
homogeneous Siu division, which produces the identity
$\chi'=\log|s_F|+\chi_B$ as an identity of sums and never rearranges it.
\end{warning}

\begin{example}[Activation on the blown-up plane]
\label{ex:ch38-activation-surface}
Continue \Cref{ex:ch38-blowup-surface-evaluation} with $d=3$, so that
$X=\mathbf P^2$, $Y=Y_2$, $B_{Y_2}=-2E_1'-3E_2$,
$D=2E_1'+3E_2$, $\theta=c_1(\cO(3))$, and
$P=3H-2e_1-e_2$ is nef but not ample.

\emph{The divisor $F$.}  Here $\pi$ is already a composite of blowups,
so Step~1 of \Cref{lem:ch38-activation} does not change $Y$.  The set
$Z$ of Step~2 is the single point $p$, because the non-isomorphism
locus, $\pi(\Supp D)$ and $\pi(\Supp G_0)$ all equal $\{p\}$.  Take
$A$ a line and $m=1$: let $F=L$ be a line through $p$ other than
the one in the direction $q$.  Then $[F]=c_1(\cO(1))$ is ample and
$\omega=\theta+[F]=c_1(\cO(4))$ is ample.  With $L''$ the strict
transform of $L$ on $Y_2$,
\[
 \pi^*L=L''+E_1'+E_2,
 \qquad
 D'=D+\pi^*L=L''+3E_1'+4E_2 .
\tag{38.28}\label{eq:ch38-activation-surface-D}
\]
Indeed $\pi_1^*L=L'+E_1$ and $\pi_2^*L'=L''$, $L'$ not passing
through $q$, while $\pi_2^*E_1=E_1'+E_2$.  The support of $D'$ is
$\{L'',E_1',E_2\}=\pi^{-1}(L)$, as
\Cref{lem:ch38-activation}\,\cref{itm:ch38-act:supports} predicts.

\emph{The divisor $G$.}  A $\pi$-ample $-G_0$ is found by solving two
inequalities.  Writing $G_0=aE_1'+bE_2$, the intersection table
\eqref{eq:ch38-surface-intersection} gives
\[
 (-G_0)\cdot E_2=b-a,\qquad (-G_0)\cdot E_1'=2a-b,
\]
so $-G_0$ is $\pi$-ample exactly when $a<b<2a$; take
$G_0=2E_1'+3E_2$.  Normalizing as in Step~3 with $r=6$ gives
\[
 G:=\tfrac16G_0=\tfrac13E_1'+\tfrac12E_2,
 \qquad
 2G=\tfrac23E_1'+E_2\leq L''+E_1'+E_2=\pi^*L,
\]
so \Cref{lem:ch38-activation}\,\cref{itm:ch38-act:dominated} holds.  For
\cref{itm:ch38-act:Qample}, the class of $G$ is
$\tfrac13e_1+\tfrac16e_2$, and \eqref{eq:ch38-surface-intersection}
gives
\[
 (\pi^*\omega-G)\cdot E_2=\tfrac16,\quad
 (\pi^*\omega-G)\cdot E_1'=\tfrac13-\tfrac16=\tfrac16,\quad
 (\pi^*\omega-G)\cdot\Lambda=4-\tfrac13-\tfrac16=\tfrac72,
\]
with $(\pi^*\omega-G)^2=16-\tfrac19-\tfrac1{36}=\tfrac{571}{36}>0$.  The
Mori cone of $Y_2$ is generated by the negative curves $E_1'$,
$E_2$ and $\Lambda$ --- these are the only irreducible curves of
negative self-intersection on this weak del Pezzo surface of degree
$7$ --- so the Nakai--Moishezon criterion gives that
$\pi^*\omega-G$ is ample.

\emph{Relative nefness.}  Finally
$(-D')\cdot E_2=-(L''+3E_1'+4E_2)\cdot E_2=-(0+3-4)=1$ and
$(-D')\cdot E_1'=-(1-6+4)=1$, both nonnegative, in accordance with
\Cref{lem:ch38-activation}\,\cref{itm:ch38-act:pinef}; and indeed
$P\cdot E_2=P\cdot E_1'=1$ by
\eqref{eq:ch38-surface-nef-test}, which is the same computation, as it
must be.
\end{example}

\section{The ample buffer}
\label{sec:ch38-buffer}

The total class $P=\pi^*\omega-D'$ is nef, and
\Cref{prop:ch38-model-fs-dictionary} needs an ample one.  The buffer is a
one-parameter family of divisors $E_t<D'$, supported on
$\Supp D'$ and increasing to $D'$, for which the total class is ample
for every $t>0$.  Its construction is a single algebraic identity, and
it is worth displaying that identity before the lemma: with
$Q:=rD'-G$,
\[
 \begin{aligned}
  (1+rt)\bigl(\pi^*\omega-E_t\bigr)
   &=(1+rt)\pi^*\omega-D'-tG\\
   &=\bigl(\pi^*\omega-D'\bigr)
     +t\bigl(r\pi^*\omega-G\bigr)\\
   &=P+t\bigl(r\pi^*\omega-G\bigr),
 \end{aligned}
\]
so the total class at parameter $t$ is a convex combination of the nef
class $P$ and the ample class $r\pi^*\omega-G$, rescaled.  Everything
else is bookkeeping.

\begin{lemma}
\label{lem:ch38-buffer}
Keep the notation and conclusions of \Cref{lem:ch38-activation}, and let
$H$ be a nef $\mathbf R$-Cartier divisor class on $X$, possibly zero.
Choose $r\geq1$ such that
\[
                              Q:=rD'-G
\tag{38.29}\label{eq:bj-buffer-Q}
\]
has strictly positive coefficient on every component of $\Supp D'$ and
such that
\[
                         r\pi^*\omega+\pi^*H-G
\tag{38.30}\label{eq:bj-buffer-ample-divisor}
\]
is ample; both hold for $r=1$ and $H=0$, by
\Cref{lem:ch38-activation}\,\cref{itm:ch38-act:dominated} and
\cref{itm:ch38-act:Qample}.  For $t>0$ define
\[
 E_t:=\frac{D'+tG}{1+rt},\qquad
 \Omega_t:=\omega+\frac{t}{1+rt}[H].
\tag{38.31}\label{eq:bj-buffer-path}
\]
Then $E_t$ is an effective $\mathbf R$-Cartier divisor supported on
$\Supp D'$, the coefficient margin is
\[
                         D'-E_t=\frac{tQ}{1+rt}>0,
\tag{38.32}\label{eq:bj-buffer-margin}
\]
so that $0<E_t<D'$ and $E_t\uparrow D'$ as $t\downarrow0$, and the
total class is
\[
 \pi^*\Omega_t-E_t
 =\frac{P+t(r\pi^*\omega+\pi^*H-G)}{1+rt},
\tag{38.33}\label{eq:bj-buffer-ample-total}
\]
which is ample for every $t>0$.  With $H=0$ one has
$\Omega_t=\omega$ for every $t$, and
$\pi^*\omega-E_t$ is ample.
\end{lemma}

\begin{proof}
\emph{Admissibility of $r=1$, $H=0$.}  By
\Cref{lem:ch38-activation}\,\cref{itm:ch38-act:dominated} we have
$2G\leq\pi^*F$, hence $\pi^*F-G\geq\tfrac12\pi^*F$; and
$D'=D+\pi^*F$ with $D\geq0$.  Therefore, for every $r\geq1$,
\[
 Q=rD'-G\geq D'-G=D+(\pi^*F-G)\geq D+\tfrac12\pi^*F
 \geq\tfrac12\pi^*F .
\]
Since $\Supp D'=\Supp\pi^*F$ by
\Cref{lem:ch38-activation}\,\cref{itm:ch38-act:supports}, the divisor
$\tfrac12\pi^*F$ has strictly positive coefficient on every component of
$\Supp D'$; hence so does $Q$, which is the first requirement in
\eqref{eq:bj-buffer-Q}.  For \eqref{eq:bj-buffer-ample-divisor} with
$r=1$ and $H=0$ the class is $\pi^*\omega-G$, ample by
\Cref{lem:ch38-activation}\,\cref{itm:ch38-act:Qample}.  For general
$r\geq1$ and nef $H$,
\[
 r\pi^*\omega+\pi^*H-G=(\pi^*\omega-G)+(r-1)\pi^*\omega+\pi^*H
\]
is the sum of an ample class and two nef classes, by
\Cref{lem:ch38-numerical-birational}\,\cref{itm:ch38-num:inj}, hence ample.

\emph{The margin.}  Directly from \eqref{eq:bj-buffer-path},
\[
 D'-E_t=\frac{(1+rt)D'-D'-tG}{1+rt}=\frac{t(rD'-G)}{1+rt}
 =\frac{tQ}{1+rt},
\]
which is \eqref{eq:bj-buffer-margin} and is strictly positive on every
component of $\Supp D'$ for $t>0$.  Hence $E_t<D'$.  Effectivity of
$E_t$ is clear from $D'\geq0$ and $G\geq0$, and
$\Supp E_t\subseteq\Supp D'\cup\Supp G=\Supp D'$.  Since
$t\mapsto tQ/(1+rt)$ is strictly increasing on $(0,\infty)$ for each
coefficient of $Q$, the family $D'-E_t$ increases with $t$; so
$E_t$ decreases with $t$ and increases to $D'$ as $t\downarrow0$.
Finally $E_t>0$ on every component of $\Supp D'$ for every $t>0$,
because $E_t\geq D'/(1+rt)$ there and $D'>0$ on
$\Supp D'$.

\emph{The total class.}  Substituting \eqref{eq:bj-buffer-path},
\[
 (1+rt)\bigl(\pi^*\Omega_t-E_t\bigr)
 =(1+rt)\pi^*\omega+t\pi^*H-D'-tG
 =\bigl(\pi^*\omega-D'\bigr)+t\bigl(r\pi^*\omega+\pi^*H-G\bigr),
\]
and $\pi^*\omega-D'=P$ by
\Cref{lem:ch38-activation}\,\cref{itm:ch38-act:total}.  This is
\eqref{eq:bj-buffer-ample-total}.  The right-hand side is a sum of the nef
class $P$ and $t>0$ times the ample class
\eqref{eq:bj-buffer-ample-divisor}, hence ample; and a positive multiple of
an ample class is ample.
\end{proof}

\begin{remark}[The role of $H$, and why it may be dispensed with]
\label{rem:ch38-buffer-H}
The auxiliary class $H$ records a more flexible variant of the path; it
is not needed under the activated hypotheses.  Indeed, once $\omega$ is
ample and $-G$ is $\pi$-ample, the relative-to-absolute ampleness
criterion
\Cref{lem:ch38-numerical-birational}\,\cref{itm:ch38-num:relative} already
implies that $r\pi^*\omega-G$ is ample for all sufficiently large
$r$.  The fact that $\pi^*\omega$ itself is not ample merely explains
why the relative correction $-G$ is indispensable; it does not obstruct
the sum from becoming ample.  Activation strengthens this general fact by
scaling $G$ so that $r=1$ works, and hence permits $H=0$ and the
constant path $\Omega_t=\omega$.  The version with $H\neq0$ is retained
because
\Cref{chap:bj-reflection} uses the first, where the class correction is
cancelled by the exact identity \eqref{eq:bj-reflection-total} rather than
by domination, and \Cref{sec:ch38-realization} uses the second.
\end{remark}

\begin{example}[The buffer on the blown-up plane]
\label{ex:ch38-buffer-surface}
Continue \Cref{ex:ch38-activation-surface}: $X=\mathbf P^2$,
$Y=Y_2$, $\omega=c_1(\cO(4))$,
$D'=L''+3E_1'+4E_2$, $G=\tfrac13E_1'+\tfrac12E_2$, and
$P=3H-2e_1-e_2$.  Take $H=0$ and $r=1$.  Then
\[
 Q=D'-G=L''+\tfrac83E_1'+\tfrac72E_2>0
\]
on every component of $\Supp D'$, and
\[
 E_t=\frac{D'+tG}{1+t}
 =\frac{1}{1+t}\Bigl(L''+\bigl(3+\tfrac13t\bigr)E_1'
 +\bigl(4+\tfrac12t\bigr)E_2\Bigr),
 \qquad
 D'-E_t=\frac{t}{1+t}\,Q .
\]
The total class is
\[
 \pi^*\omega-E_t=\frac{P+t(\pi^*\omega-G)}{1+t}
 =\frac{1}{1+t}\Bigl(
   (3+4t)H-\bigl(2+\tfrac13t\bigr)e_1-\bigl(1+\tfrac16t\bigr)e_2\Bigr).
\]
Testing on the three negative curves with
\eqref{eq:ch38-surface-intersection}, and writing
$a=3+4t$, $b=2+\tfrac13t$, $c=1+\tfrac16t$, so that a class
$aH-be_1-ce_2$ has intersection numbers $c$, $b-c$ and $a-b-c$
with $E_2$, $E_1'$ and $\Lambda$ respectively:
\[
 \begin{aligned}
  (1+t)(\pi^*\omega-E_t)\cdot E_2&=1+\tfrac16t,\\
  (1+t)(\pi^*\omega-E_t)\cdot E_1'&=1+\tfrac16t,\\
  (1+t)(\pi^*\omega-E_t)\cdot\Lambda&=\tfrac72t,
 \end{aligned}
\]
all strictly positive for $t>0$, whereas at $t=0$ the third vanishes:
this is exactly the failure of ampleness of $P$ recorded in
\eqref{eq:ch38-surface-nef-test}.  The square is
$a^2-b^2-c^2=4+\tfrac{67}3t+\tfrac{571}{36}t^2>0$, so the
Nakai--Moishezon criterion confirms ampleness for every $t>0$; at
$t=0$ it gives $P^2=9-4-1=4>0$ with $P\cdot\Lambda=0$, the
signature of a nef class on the boundary of the ample cone.  The buffer
therefore does exactly one thing here: it turns the single vanishing
intersection number $P\cdot\Lambda=0$ into the positive number
$\tfrac72t/(1+t)$, at the cost of shrinking the divisor $D'$ by
$tQ/(1+t)$.
\end{example}

\section{The realization theorem}
\label{sec:ch38-realization}

\begin{theorem}\label{thm:bj-real-cartier-realization}
Let $X$ be normal, let $\theta\in N^1(X)_{\mathbf R}$, and let
$B\le0$ be a real Cartier b-divisor determined on a normal projective
birational model $\pi:Y\to X$.  If
\[
                         \pi^*\theta+B_Y
\tag{38.34}\label{eq:bj-realization-nef-hypothesis}
\]
is nef, then there is a unique
\[
                     \chi_B\in\PSH_{\mathrm{hom},X}(\theta)
\]
whose restriction to $X^{\divis}$ is $\psi_B$.  Moreover,
$\chi_B$ admits a decreasing homogeneous rational GF regularization in
rational classes tending to $\theta$.  The result does not use the
envelope property.
\end{theorem}

\begin{proof}
We first record the proof architecture.  The obstruction is that the
hypothesis is numerical and the conclusion is analytic, while the only
bridge between the two --- the dictionary
\Cref{prop:ch38-model-fs-dictionary} --- requires a semiample total class
and rational coefficients.  The proof therefore manufactures a family of
rational data with ample total class, arranges that family to be
\emph{increasing in the divisor} and \emph{dominated in the class}, and
passes to the limit. The two arrangements serve different purposes:
increasing divisors give decreasing functions, the limit
operation used here before the envelope theorem
(\Cref{warn:ch38-firewall}); and dominated classes keep every stage inside
one fixed space $\PSH_X(\omega)$, so its fixed-class
decreasing-limit clause applies
(\Cref{warn:ch38-drift-not-absorbed}).
If $\dim X=0$, then $Y=X$, $N^1(X)=0$, and $B=0$.  The zero
homogeneous function is the unique realization.  Assume $\dim X>0$, and
put
\[
                              D:=-B_Y\ge0.
\]

\medskip
\noindent\emph{Step 1: activate the whole support.}
Apply \Cref{lem:ch38-activation}.  After replacing $Y$ by a dominating
model, which changes neither $B$ nor the nefness of the total class, by
\Cref{lem:ch38-determination-equivalence} and
\Cref{lem:ch38-nef-total-intrinsic}\,\cref{itm:ch38-nt:indep}, we obtain an
effective $\mathbf Q$-Cartier divisor $G$ on $Y$ with $-G$
$\pi$-ample, and an effective Cartier divisor $F$ on $X$ whose
support contains the non-isomorphism locus of $\pi$, $\pi(\Supp D)$,
and $\pi(\Supp G)$, such that
\[
                              \omega:=\theta+[F]
\tag{38.35}\label{eq:bj-activated-class}
\]
is ample.  Put
\[
                         B':=B-\overline F,\qquad
                         D':=D+\pi^*F.
\tag{38.36}\label{eq:bj-activated-data}
\]
Then
\[
 \pi^*\omega+B'_Y=\pi^*\theta+B_Y=P,\qquad
 \Supp D'=\pi^{-1}(\Supp F).
\tag{38.37}\label{eq:bj-activated-total}
\]
On every $\pi$-contracted curve, $-D'$ has the same intersection
number as $P$; hence $-D'$ is $\pi$-nef.  In addition
$2G\leq\pi^*F$ and $\pi^*\omega-G$ is ample, by
\Cref{lem:ch38-activation}\,\cref{itm:ch38-act:dominated} and
\cref{itm:ch38-act:Qample}.

\medskip
\noindent\emph{Step 2: insert an ample buffer.}
Apply \Cref{lem:ch38-buffer} with $H=0$ and $r=1$.  Thus, with
$Q:=D'-G$, which has strictly positive coefficient on every component of
$\Supp D'$, and with
\[
 E_t:=\frac{D'+tG}{1+t}\qquad(t>0),
\]
one has $0<E_t<D'$, $E_t\uparrow D'$ as $t\downarrow0$,
$\Supp E_t\subseteq\Supp D'$, and
\[
 \pi^*\omega-E_t=\frac{P+t(\pi^*\omega-G)}{1+t}
\]
is ample for every $t>0$, being a positive combination of the nef class
$P$ and the ample class $\pi^*\omega-G$.  The class $\omega$ is the
same at every value of $t$: the real path carries no drift.

\medskip
\noindent\emph{Step 3: choose ordered rational stages.}
Apply \Cref{lem:ch38-ordered-rational-dominated} with
$\rho=\pi$, $E=D'$ and the ample class $\omega$; its hypothesis
\eqref{eq:ch38-dominated-hypotheses} is exactly what Step~2 provides, and
the support $S$ of the whole family is the finite set of components of
$\pi^{-1}(\Supp F)$.  We obtain rational Cartier divisors $R_n$
supported on $S$ and rational classes $\omega_n$ such that
\[
 \begin{aligned}
 0&\le R_n<R_{n+1}<D',& R_n&\longrightarrow D',\\
 \omega_n&\longrightarrow\omega,
 &\pi^*\omega_n-R_n&\quad\text{is ample},
 \end{aligned}
\tag{38.38}\label{eq:bj-realization-rational-stages}
\]
and, in addition, $\omega-\omega_n$ is ample for every $n$.  This last
property is the drift orientation of
\Cref{prop:ch38-drift-ledger}\,\cref{itm:ch38-drift:orient}; it is not part
of \Cref{lem:bj-ordered-rational-approximation} and it is what the rest of
the proof uses.

Choose $L_n\in\Pic(X)_{\mathbf Q}$ with $c_1(L_n)=\omega_n$, possible
because $\omega_n$ is a rational numerical class and
$N^1(X)_{\mathbf Q}$ is by definition the image of
$\Pic(X)_{\mathbf Q}$.  Ampleness is numerical, so the actual rational
line bundle $\pi^*L_n-R_n$ is ample, hence semiample.  Formula
\eqref{eq:bj-rational-coefficient}, that is
\Cref{prop:ch38-model-fs-dictionary}\,\cref{itm:ch38-dict:forward} applied
to the rational Cartier b-divisor $-\overline{R_n}$ determined on
$Y$, gives
\[
 \chi_n:=\psi_{-\overline{R_n}}
       \in \cH_{\mathrm{hom}}(L_n)
       =\cH^{\mathrm{gf}}_0(L_n)
       \subset\PSH_{\mathrm{hom},X}(\omega_n).
\tag{38.39}\label{eq:bj-realization-gf-stages}
\]
The inequalities $R_n<R_{n+1}$ imply
$\chi_{n+1}\le\chi_n$ on $X^{\divis}$, and psh comparison extends
the order to $X^{\An}$.

In detail: on $X^{\divis}$ the inequality is
\Cref{prop:ch38-evaluation-welldefined}\,\cref{itm:ch38-ev:linear}, since
$-\overline{R_{n+1}}\leq-\overline{R_n}$; and
\Cref{cor:ch35-divisorial-determination}\,(1), applied to the two psh
functions $\chi_{n+1}$ and $\chi_n$, promotes it to $X^{\An}$.

\medskip
\noindent\emph{Step 3$'$: all stages lie in one class.}
Since $\omega-\omega_n$ is ample, hence nef,
\Cref{lem:ch34-psh-nonempty} gives $0\in\PSH_X(\omega-\omega_n)$, and
\Cref{thm:bj-basic-closure}\,(3) with $s=t=1$ yields
\[
 \chi_n\in\PSH_X\bigl(\omega_n+(\omega-\omega_n)\bigr)=\PSH_X(\omega)
\]
for every $n$; homogeneity is unaffected, so
$\chi_n\in\PSH_{\mathrm{hom},X}(\omega)$.  This is the step at which the
class drift is disposed of, and it uses the \emph{orientation} of the
drift and not its size; see \Cref{warn:ch38-drift-not-absorbed}.

\medskip
\noindent\emph{Step 4: take the psh limit.}
Let
\[
                              \chi':=\lim_n\chi_n,
\]
the pointwise limit of a decreasing sequence in
$\mathbf R\cup\{-\infty\}$, which exists in $\Rbar$ by ledger
entry~(1) of \Cref{sec:extended-arithmetic} and
\Cref{lem:monotone-net-converges}, with values in
$[-\infty,0]$.
Each homogeneous stage vanishes at the trivial valuation.  Therefore
\[
                         \chi'(v_{\triv})=0,
\tag{38.40}\label{eq:bj-realization-trivial-value}
\]
which excludes the identically $-\infty$ alternative.  Since the sequence lies in $\PSH_X(\omega)$ and its limit
is not identically $-\infty$, PSH2 yields
\[
                         \chi'\in\PSH_{\mathrm{hom},X}(\omega).
\tag{38.41}\label{eq:bj-realization-limit}
\]
Precisely: by Step~3$'$ the sequence $(\chi_n)$ is a decreasing
sequence in the single space $\PSH_X(\omega)$, and by
\eqref{eq:bj-realization-trivial-value} its limit is not identically
$-\infty$; so \Cref{thm:bj-basic-closure}\,(4) applies.  Homogeneity of
$\chi'$ is inherited from the $\chi_n$, because an infimum of
homogeneous functions is homogeneous: for $t>0$,
$\chi'(tv)=\inf_n\chi_n(tv)=\inf_nt\chi_n(v)=t\inf_n\chi_n(v)=t\chi'(v)$,
the third equality by \Cref{lem:extended-real-safe}\,(2).

Fix $v\in X^{\divis}$, and pull all divisors to a model on which $v$
appears.  Cartier evaluation is linear there, so
\[
 R_n(v)\longrightarrow D'(v),\qquad
 \chi'(v)=B'(v).
\tag{38.42}\label{eq:bj-realization-trace-limit}
\]
Indeed, write $v=c\,\ord_E$ and choose a model $\rho:Z\to Y$ carrying
$E$.  The map $R\mapsto v(R)=c\,\ord_E(\rho^*R)$ is a linear form on
the finite-dimensional space $V$ of
\Cref{lem:ch38-fixed-support-rational}, hence continuous; since
$R_n\to D'$ in $V$, we get $v(R_n)\to v(D')$.  Therefore
\[
 \chi'(v)=\lim_n\chi_n(v)=\lim_n\bigl(-v(R_n)\bigr)=-v(D')=v(B'_Y)
 =\psi_{B'}(v),
\]
all the quantities being finite real numbers.

\medskip
\noindent\emph{Step 5: remove the activation divisor.}
Because $B\le0$,
\[
                    B'=B-\overline F\le-\overline F.
\]
Hence
\[
                         \chi'\le\log|s_F|
\tag{38.43}\label{eq:bj-realization-pole-bound}
\]
on $X^{\divis}$, and then on $X^{\An}$ by comparison.  Homogeneous
Siu division gives the global identity
\[
 \chi'=\log|s_F|+\chi_B,\qquad
 \chi_B\in\PSH_{\mathrm{hom},X}(\theta).
\tag{38.44}\label{eq:bj-realization-siu-division}
\]
At a point supported in $F$, this identity is supplied by the division
theorem; it is not rearranged into a difference of two $-\infty$ values.
At divisorial valuations, $v(F)<\infty$, and
\eqref{eq:bj-realization-trace-limit} shows that the trace of $\chi_B$
is $B$.

We expand the two assertions.  For \eqref{eq:bj-realization-pole-bound}:
on $X^{\divis}$,
$\psi_{B'}\leq\psi_{-\overline F}=\log|s_F|$ by
\Cref{prop:ch38-evaluation-welldefined}\,\cref{itm:ch38-ev:linear} and
\Cref{lem:ch38-sign-tests}\,\cref{itm:ch38-sign:pole}; the function
$\log|s_F|$ is upper semicontinuous on $X^{\An}$ with values in
$[-\infty,0]$, by \eqref{eq:divisor-section-pole}; and $\chi'$ is
$\omega$-psh; so \Cref{thm:ch35-divisorial-comparison} promotes the
inequality from $X^{\divis}$ to $X^{\An}$.  For
\eqref{eq:bj-realization-siu-division}: apply
\Cref{thm:ch35-homogeneous-siu-general} with the class $\omega$, the
nonzero effective $\mathbf Q$-Cartier divisor $E:=F$ and the
homogeneous $\omega$-psh function $\chi'$; its hypothesis is
\eqref{eq:bj-realization-pole-bound}, and its conclusion is the existence
of a unique
$\chi_B\in\PSH_{\mathrm{hom},X}(\omega-[F])=\PSH_{\mathrm{hom},X}(\theta)$
with $\chi'=\log|s_F|+\chi_B$ on all of $X^{\An}$.  The theorem is
quoted in its general form because the strict hypothesis
\eqref{eq:ch35-siu-strict-hypothesis} of
\Cref{thm:ch35-homogeneous-siu} may fail here;
\Cref{ex:ch38-strict-room-fails} exhibits the failure, and
\Cref{rem:ch38-siu-room} discusses when the strict form suffices.

Finally the trace.  For $v\in X^{\divis}$ all three terms of
\eqref{eq:bj-realization-siu-division} are finite, $v(F)$ being a real
number, so the identity may be solved for $\chi_B(v)$:
\[
 \chi_B(v)=\chi'(v)-\log|s_F|(v)
 =\psi_{B'}(v)+v(F)
 =\bigl(\psi_B(v)-v(F)\bigr)+v(F)=\psi_B(v).
\]
This subtraction is legitimate because it is a subtraction of the finite
number $\log|s_F|(v)=-v(F)$; at a point of proper support no such
manipulation is performed, and \Cref{prop:ch38-safeguards}\,(4) records
the fact.

\medskip
\noindent\emph{Step 6: uniqueness and regularization.}
Homogeneous psh functions are determined by their divisorial traces, which
proves uniqueness.  The general homogeneous GF regularization theorem gives
the asserted decreasing rational GF regularization of $\chi_B$ in
classes tending to $\theta$.

In detail: if $\chi,\chi''\in\PSH_{\mathrm{hom},X}(\theta)$ both restrict
to $\psi_B$ on $X^{\divis}$, then $\chi=\chi''$ on $X^{\An}$ by
\Cref{cor:ch35-divisorial-determination}\,(2).  For the regularization,
$\chi_B$ is homogeneous with $\sup\chi_B=\chi_B(v_{\triv})=0$ by
\Cref{prop:ch35-homogeneous-normalization}, so it is its own
homogenization and \Cref{cor:ch35-homogenization-psh} produces a
decreasing net $(\psi_a)$ with
$\psi_a\in\cH_{\mathrm{hom}}(M_a)$, $c_1(M_a)\to\theta$ and
$\psi_a\searrow\chi_B$ pointwise.  Each $\psi_a$ is a homogeneous
rational generically finite Fubini--Study function, which is the assertion.
\end{proof}

\begin{remark}[Where each hypothesis enters]
\label{rem:ch38-hypothesis-map}
Normality of $X$ is used three times: to have the model category of
\Cref{sec:appF-models} at all; through $\pi_*\cO_Y=\cO_X$ in
\Cref{prop:ch38-model-fs-dictionary}\,\cref{itm:ch38-dict:forward}; and in
the regularity criterion inside the proof of
\Cref{lem:ch35-fs-division}, which underlies
\Cref{thm:ch35-homogeneous-siu-general}.  Antieffectivity of $B$ is used
twice: to make $\psi_B\leq0$, which is what
\Cref{thm:ch35-semiample-criterion} requires at each rational stage, and
to give \eqref{eq:bj-realization-pole-bound}.  Nefness of the total class
is the only numerical input; it is used in Step~2, to make
$P+t(\pi^*\omega-G)$ ample for $t>0$, and again in Step~1, to make
$-D'$ $\pi$-nef.  \Cref{ex:ch38-usc-not-psh} shows that it cannot be
\emph{dropped}; whether it can be weakened to pseudoeffectivity of $P$ is
addressed neither there nor here.  That
$B$ is Cartier, that is, determined on a single model, is used in
Step~3, through the fixed-support hypothesis of
\Cref{lem:ch38-ordered-rational-dominated}; \Cref{ex:ch38-moving-support}
shows that this is not a technical convenience.  No hypothesis on
$\operatorname{char}k$ and no smoothness hypothesis occurs anywhere.
\end{remark}

\begin{remark}[Strict room and the borderline Siu case]
\label{rem:ch38-siu-room}
The division in Step~5 is performed with respect to the divisor $F$ and
the constant $t=1$.  By \Cref{lem:ch35-slice-to-cone}, the strict
hypothesis \eqref{eq:ch35-siu-strict-hypothesis} of
\Cref{thm:ch35-homogeneous-siu} amounts to
$\sup_{K_F}\chi'<-1$, that is, to
$\psi_{B'}\leq-t_0\psi_{\overline F}$ on $X^{\divis}$ for some
$t_0>1$, that is, to
\[
 B\leq(1-t_0)\overline F\qquad\text{for some }t_0>1 .
\]
Equivalently, $B$ has coefficient bounded away from $0$ by a fixed
negative multiple of the corresponding coefficient of $\pi^*F$, along
every component of $\Supp\pi^*F$.  When this holds --- for instance when
$D=-B_Y\geq\varepsilon\pi^*F$ for some $\varepsilon>0$ --- the strict
theorem \Cref{thm:ch35-homogeneous-siu} suffices.
The general Siu result is the foundational input
\cite[Theorem~6.21]{BJ22}. When strict room fails, as it does in
\Cref{ex:ch38-strict-room-fails}, the general form
\Cref{thm:ch35-homogeneous-siu-general} is used.  That theorem carries no
envelope hypothesis --- this is asserted in its statement and in the
warning following it --- so the firewall
\Cref{warn:ch38-firewall} is respected either way.  It is the only
analytic input in the main realization proof that is invoked at a
borderline not reproved in this book.
\end{remark}

\begin{example}[The strict hypothesis genuinely fails]
\label{ex:ch38-strict-room-fails}
Take $B=0$, which certainly satisfies the hypotheses of
\Cref{thm:bj-real-cartier-realization} whenever $\theta$ is nef, with
$P=\pi^*\theta$.  Then $D=0$, $D'=\pi^*F$, and
$B'=-\overline F$, so $\chi'=\log|s_F|$ exactly and
$\sup_{K_F}\chi'=-1$: the inequality
\eqref{eq:bj-realization-pole-bound} is an equality on the compact set
$K_F$ of \eqref{eq:ch35-siu-slice}.  No choice of $F$ repairs this:
the activation divisor is required to contain the exceptional locus and
$\pi(\Supp D)$, so it can only be enlarged, and enlarging it makes
$B'$ more negative by the same amount.  The same phenomenon occurs whenever $B$ has
coefficient $0$ along some component of $\Supp\pi^*F$, which happens
as soon as $\Supp D$ does not meet that component --- and $\Supp F$ is
in general strictly larger than $\pi(\Supp D)$, since it must also
contain the non-isomorphism locus of $\pi$.  Thus the borderline case of
homogeneous Siu division is not an artefact of the presentation; it is the
generic situation for this application.  For $B=0$ the conclusion is of
course trivial, $\chi_0=0$; the point of the example is that no
hypothesis available at the start of the proof distinguishes this case
from the general one.
\end{example}

\begin{remark}[What the proof does not do]
\label{rem:ch38-proof-negative}
Three natural-looking shortcuts are unavailable and are not used.  One
cannot define $\chi_B$ as the upper-semicontinuous regularization of
$\psi_B$ on $X^{\An}$: that function exists, but nothing identifies it
as psh, and \Cref{ex:ch38-usc-not-psh} shows that it need not be.  One
cannot obtain $\chi_B$ as an increasing limit of the functions
$\psi_B+\varepsilon\psi_H$ of the proof of
\Cref{thm:ch35-cartier-psh-criterion}, since increasing limits are exactly
what \Cref{warn:ch38-firewall} forbids.  And one cannot avoid the
activation divisor by working directly with $P$: the total class is nef,
but need not supply a globally generated multiple, and $\theta$ itself need not even be nef, by
\Cref{warn:ch38-theta-not-nef}.
\end{remark}

\section{Extended-real safeguards}
\label{sec:ch38-safeguards}

Every function produced in this chapter takes values in
$\mathbf R\cup\{-\infty\}$, and the value $-\infty$ is attained,
already in \Cref{ex:bj-fixed-divisor-coefficient}.  The forbidden
operation of ledger entry~(4) of \Cref{sec:extended-arithmetic} is the
formation of $(+\infty)+(-\infty)$, equivalently of a difference of two
infinite values of the same sign.  We list every place where such a
difference could arise, and show that none of them does.  The list is
complete for this chapter, in the same sense as the corresponding list in
\Cref{sec:ch34-minus-infinity}.

\subsection{The audit}
\label{sec:ch38-audit}

\begin{proposition}
\label{prop:ch38-safeguards}
In the notation of \Cref{thm:bj-real-cartier-realization} and its proof,
the following hold.
\begin{enumerate}
\item\label{itm:ch38-safe:eval}
 \emph{Evaluation.}  The expression $v(B_Y)$ is formed only for
 $v\in X^{\val}$, where it is a difference of two \emph{finite}
 numbers, and for $v\in X^{\divis}\subseteq X^{\val}$.  It is never
 formed at a point of proper support.  The extension of
 $\psi_C$ to $X^{\An}$ for an antieffective rational $C$ with
 semiample total class is $m^{-1}\log|\mathfrak b|$, a single function
 with values in $[-\infty,0]$, and not a difference.
\item\label{itm:ch38-safe:activation}
 \emph{Activation.}  The identity
 $\psi_{B'}=\psi_B+\log|s_F|$ is used only on $X^{\val}$, where all
 three terms are finite real numbers.  No analytic extension of $\psi_B$
 is assumed during activation.  The corresponding global identity on
 $X^{\An}$ is formed only after realization, as the Siu identity in
 \cref{itm:ch38-safe:siu}.
\item\label{itm:ch38-safe:limit}
 \emph{Limits.}  The limit $\chi'=\lim_n\chi_n$ is the infimum of a
 decreasing sequence in $[-\infty,0]$; it exists in $\Rbar$ by ledger
 entry~(1).  Homogeneity of $\chi'$ is obtained from
 $\inf_n(t\chi_n)=t\inf_n\chi_n$ for $t>0$
 (\Cref{lem:extended-real-safe}\,(2)); the value $t=0$ is never used,
 the product $0\cdot(-\infty)$ being undefined by ledger entry~(3).
\item\label{itm:ch38-safe:trace}
 \emph{The trace computation.}  The only subtraction of two possibly
 infinite quantities that the proof appears to perform is the passage from
 \eqref{eq:bj-realization-siu-division} to
 $\chi_B(v)=\chi'(v)-\log|s_F|(v)$.  It is performed exclusively for
 $v\in X^{\divis}$, where $\log|s_F|(v)=-v(F)$ is a finite real
 number, so that the operation is the subtraction of a finite constant,
 legitimate by ledger entry~(2).
\item\label{itm:ch38-safe:siu}
 \emph{The Siu identity.}  Equation
 \eqref{eq:bj-realization-siu-division} is an identity of sums, valid at
 every point of $X^{\An}$, with both summands in $[-\infty,0]$.  It is
 supplied by \Cref{thm:ch35-homogeneous-siu-general} as an identity and is
 never rearranged.  At a semivaluation $v$ with
 $s(v)\subseteq\Supp F$ both sides equal $-\infty$, and the value of
 $\chi_B$ at such a $v$ is \emph{not} determined by
 \eqref{eq:bj-realization-siu-division}; it is determined by
 \Cref{cor:ch35-divisorial-determination}\,(2), that is, by
 plurisubharmonic determination from the divisorial trace.
\item\label{itm:ch38-safe:comparison}
 \emph{Comparisons.}  All inequalities between extended-real functions
 --- $\chi_{n+1}\leq\chi_n$, $\chi'\leq\log|s_F|$,
 $\chi_{B_1}\leq\chi_{B_2}$ --- are formulated as inequalities and never
 as bounds on a difference, following ledger entry~(4) and
 \Cref{rem:forbidden-cancellations}.
\item\label{itm:ch38-safe:finite}
 \emph{Finite-dimensional operations.}  All differences of divisors and of
 numerical classes, namely $D'-E_t$, $R_{n+1}-R_n$,
 $\omega-\omega_n$, $\pi^*\omega-G$, and $Q=rD'-G$, are formed in
 finite-dimensional real vector spaces, where no infinite value occurs.
\item\label{itm:ch38-safe:sup}
 \emph{Suprema.}  The only supremum formed is
 $\sup_{K_F}\chi'$ in \Cref{lem:ch35-slice-to-cone}, which is defined for
 every family by ledger entry~(5) and is finite or $-\infty$.  No
 supremum of a family of psh functions is formed anywhere, by
 \Cref{warn:ch38-firewall}.
\end{enumerate}
\end{proposition}

\begin{proof}
\emph{\cref{itm:ch38-safe:eval}} is
\Cref{prop:ch38-evaluation-domains}: parts
\cref{itm:ch38-dom:divis} and \cref{itm:ch38-dom:val} confine the formula
$v(B_Y)$ to $X^{\val}$, where the local expression of $B_Y$ as a
difference of two effective Cartier divisors evaluates to a difference of
two finite numbers because a nonzero ideal is not contained in the ideal
of an irreducible component; part \cref{itm:ch38-dom:an} produces the
extension to $X^{\An}$ as a single ideal function.  That the formula
genuinely fails at proper-support points is
\Cref{ex:ch38-evaluation-no-lift} and
\Cref{ex:ch38-fractional-undefined}.

\emph{\cref{itm:ch38-safe:activation}.}  On $X^{\divis}$ the identity is
\Cref{prop:ch38-evaluation-welldefined}\,\cref{itm:ch38-ev:linear} applied
to $B'=B+(-\overline F)$, together with
\Cref{lem:ch38-sign-tests}\,\cref{itm:ch38-sign:pole}.  The same
coefficient calculation holds on $X^{\val}$ by
\Cref{prop:ch38-evaluation-domains}\,\cref{itm:ch38-dom:val}; there every
term is finite.  Thus this step neither forms an extended-real difference
nor presupposes the psh extension that the realization theorem is proving.

\emph{\cref{itm:ch38-safe:limit}.}  A decreasing sequence in the complete
lattice $\Rbar$ converges to its infimum,
\Cref{lem:monotone-net-converges}.  The identity
$\inf_n(t\chi_n)=t\inf_n\chi_n$ for $t>0$ is
\Cref{lem:extended-real-safe}\,(2); it is exactly the statement that
multiplication by a positive real is an order isomorphism of $\Rbar$.
The scaling action of \Cref{chap:trivial-analytification} is by
$t>0$ only, so $t=0$ does not arise.

\emph{\cref{itm:ch38-safe:trace}.}  For $v\in X^{\divis}$, write
$v=c\ord_E$; then $v(F)=c\ord_E(F)\in[0,\infty)$, a finite number,
because $F$ is a divisor on $X$ and $\ord_E$ is a divisorial
valuation of $k(X)$, which takes finite values on nonzero rational
functions.  Hence $\log|s_F|(v)=-v(F)\in(-\infty,0]$, and subtracting it
from the finite number $\chi'(v)$ is subtraction of a finite constant.
That $\chi'(v)$ is finite is
\eqref{eq:bj-realization-trace-limit} together with
\Cref{prop:ch38-evaluation-welldefined}\,\cref{itm:ch38-ev:finite}.

\emph{\cref{itm:ch38-safe:siu}.}  The identity and the uniqueness of
$\chi_B$ are the content of
\Cref{thm:ch35-homogeneous-siu-general}.  For the last assertion, let
$v\in X^{\An}$ with $s(v)\subseteq\Supp F$.  Then
$\log|s_F|(v)=-\infty$ by \eqref{eq:divisor-section-pole}, and
$\chi_B(v)\leq0$ by \Cref{prop:ch35-homogeneous-normalization}, so
$\log|s_F|(v)+\chi_B(v)=-\infty=\chi'(v)$ whatever the value of
$\chi_B(v)$ may be.  The identity therefore carries no information at
such a point, and the value there is pinned down instead by
\Cref{cor:ch35-divisorial-determination}\,(2): two homogeneous psh
functions with the same divisorial trace coincide.
\Cref{ex:ch38-naive-realization-fails} shows that this is not a
pedantic distinction.

\emph{\cref{itm:ch38-safe:comparison}} and
\emph{\cref{itm:ch38-safe:finite}} are inspection of the proof.
\emph{\cref{itm:ch38-safe:sup}} is ledger entry~(5) together with
\Cref{prop:ch35-homogeneous-normalization}, which bounds $\chi'$ above
by $0$.
\end{proof}

\subsection{A naive realization that fails}
\label{sec:ch38-naive}

\begin{example}[The naive removal of the fixed pole is not merely
undefined; it is unrepairable]
\label{ex:ch38-naive-realization-fails}
\index{extended real numbers!forbidden subtraction}
Let $X=\mathbf P^2_k$, let $L_1\neq L_2$ be two distinct lines, and
let
\[
 \theta:=c_1(\cO(1)),\qquad
 B:=-\overline{L_1},
\]
a real (indeed integral) antieffective Cartier b-divisor determined on
$X$ itself, so that $\pi=\operatorname{id}$.  The total class is
\[
 \pi^*\theta+B_X=c_1(\cO(1))-[L_1]=0,
\]
which is nef, so \Cref{thm:bj-real-cartier-realization} applies.  The
answer is known in advance:
$\chi_B=\log|s_{L_1}|\in\cH_{\mathrm{hom}}(\cO(1))$, because
$\cO(1)\otimes\cO(-L_1)\simeq\cO_X$ is globally generated and
\Cref{prop:hom-fs-ideal} applies.

Run activation with $F:=L_1+L_2$, whose support contains
$\Supp(-B_X)=L_1$, as \Cref{lem:ch38-activation} requires.
Then $\omega=\theta+[F]=c_1(\cO(3))$ is ample,
$B'=B-\overline F=-\overline{2L_1+L_2}$, and
\[
 \chi'=\log|s_{2L_1+L_2}|=\log|s_{L_1}|+\log|s_F|,
\]
an identity of functions into $[-\infty,0]$, legitimate by
ledger entry~(2) of \Cref{sec:extended-arithmetic}.

\emph{The naive removal.}  Suppose one attempted to define the answer
pointwise by
\[
 \chi_B^{\mathrm{naive}}(v):=\chi'(v)-\log|s_F|(v),
 \qquad v\in X^{\An}.
\tag{38.45}\label{eq:ch38-naive-difference}
\]
Choose two closed points
\[
 \xi\in L_2\setminus L_1,
 \qquad
 \xi':=L_1\cap L_2,
\]
and consider the proper-support points
$v_{\xi,\triv}$ and $v_{\xi',\triv}$ of
\Cref{ex:ch32-proper-support-trivial}.  At both points
$\log|s_F|=-\infty$, because $s_F$ vanishes at every point of
$L_2$; and at both points $\chi'=-\infty$, for the same reason.  So
\eqref{eq:ch38-naive-difference} reads $(-\infty)-(-\infty)$ at both.
The correct values, however, differ:
\[
 \begin{aligned}
 \chi_B(v_{\xi,\triv})
   &=\log|s_{L_1}|(v_{\xi,\triv})
     =-v_{\xi,\triv}(s_{L_1})=0,\\
 \chi_B(v_{\xi',\triv})
   &=\log|s_{L_1}|(v_{\xi',\triv})=-\infty,
 \end{aligned}
\]
the first because $s_{L_1}$ does not vanish at $\xi\notin L_1$, the
second because it does vanish at $\xi'\in L_1$.  Hence no convention
assigning a value to $(-\infty)-(-\infty)$ can make
\eqref{eq:ch38-naive-difference} correct: the same symbol would have to
denote $0$ at one point and $-\infty$ at another.  This is the failure
predicted by ledger entry~(4) of \Cref{sec:extended-arithmetic} and by
\Cref{rem:forbidden-cancellations}, in the sharpest possible form.

\emph{What replaces it.}  The proof of
\Cref{thm:bj-real-cartier-realization} never writes
\eqref{eq:ch38-naive-difference}.  It obtains
$\chi_B$ from \Cref{thm:ch35-homogeneous-siu-general} as a psh function
satisfying the \emph{additive} identity
$\chi'=\log|s_F|+\chi_B$, and then determines it by its divisorial
trace.  In this example its value at $v_{\xi,\triv}$ is read directly
from the already identified psh function
$\chi_B=\log|s_{L_1}|$, and equals zero because $s_{L_1}$ does not
vanish at $\xi$.  In general there is no local formula of this kind:
uniqueness is the psh determination theorem
\Cref{cor:ch35-divisorial-determination}, not upper semicontinuity alone.

\emph{A third point, for contrast.}  At the proper-support point
$v_{L_2,\triv}$, the trivial valuation of the curve $L_2$, the same
analysis gives $\chi'=-\infty$, $\log|s_F|=-\infty$ and the
correct value $\chi_B(v_{L_2,\triv})=0$, because $s_{L_1}$ does not
vanish identically along $L_2$.  The set on which
\eqref{eq:ch38-naive-difference} is undefined is
$\Supp(F)^{\An}$, and on it the correct answer takes both the value
$0$ and the value $-\infty$.
\end{example}

\begin{warning}{Where the value at a proper-support point comes from}
\label{warn:ch38-proper-support-value}
\Cref{ex:ch38-naive-realization-fails} isolates the single most
important structural fact about
\Cref{thm:bj-real-cartier-realization}: the values of $\chi_B$ at
proper-support points are not supplied by a general evaluation
formula for a divisor on an arbitrary determination.
Among $\theta$-psh functions they are determined by the divisorial trace,
through
\Cref{thm:ch35-divisorial-comparison} and
\Cref{cor:ch35-divisorial-determination}.  This is why the theorem is
stated as an existence-and-uniqueness statement for a psh function with a
prescribed trace rather than as the construction of an explicit function,
and it is why \Cref{sec:ch34-proper-support} had to be developed before
this chapter.  The same discipline governs
\Cref{chap:bj-reflection}: the identity
\eqref{eq:bj-reflection-global-identity} there is likewise proved on
divisorial points first and extended by determination, never by
subtraction.
\end{warning}

\begin{remark}[The ledger entries used, by number]
\label{rem:ch38-ledger-entries}
For the record, the entries of \Cref{sec:extended-arithmetic} used in this
chapter are: (1), for the existence of decreasing limits and of the
supremum in \Cref{lem:ch35-slice-to-cone}; (2), for the sums
$\psi_B+\log|s_F|$ and $\log|s_F|+\chi_B$, and for the subtraction of
the finite quantity $v(F)$ at a divisorial point; (3), for the scaling
$t\chi$ with $t>0$ in the definition of homogeneity, together with the
explicit exclusion of $t=0$; (4), for the prohibition that makes
\Cref{ex:ch38-naive-realization-fails} a failure rather than a convention;
and (5), for the supremum over $K_F$.  Entries (6) and (7) are not used:
no $\limsup$ of a net and no integral occurs in this chapter.  The
corresponding audit for \Cref{chap:bj-fs-pl-psh} is
\Cref{sec:ch34-minus-infinity}, and
\Cref{rem:ledger-uses} records the global pattern.
\end{remark}

\section{Consequences and limits}
\label{sec:ch38-consequences}

\subsection{Order, and the unconditional Cartier criterion}
\label{sec:ch38-order}

\begin{corollary}\label{cor:bj-realization-order}
Suppose $B_1\le B_2\le0$ satisfy the hypotheses of
\Cref{thm:bj-real-cartier-realization} in the same class $\theta$.  Then
\[
                             \chi_{B_1}\le\chi_{B_2}
\]
on all of $X^{\An}$.
\end{corollary}

\begin{proof}
The inequality holds for the divisorial traces.  Both realizations are psh
in the same class, so comparison extends it to every analytic point.

Explicitly: $\psi_{B_1}\leq\psi_{B_2}$ on $X^{\divis}$ by
\Cref{prop:ch38-evaluation-welldefined}\,\cref{itm:ch38-ev:linear}; the two
realizations restrict to these traces, by
\Cref{thm:bj-real-cartier-realization}; and
\Cref{cor:ch35-divisorial-determination}\,(1), applied to
$\chi_{B_1}\in\PSH_X(\theta)$ and $\chi_{B_2}\in\PSH_X(\theta)$, gives
the inequality on $X^{\An}$.
\end{proof}

\begin{corollary}
\label{cor:ch38-realization-monotone}
Let $(B_i)_{i\in I}$ be a net of real Cartier b-divisors, each
$\leq0$, each nef-total in the class $\theta_i$, with
$\theta-\theta_i$ nef for every $i$.  Suppose the net is decreasing:
$B_j\leq B_i$ for $j\geq i$.  Then $(\chi_{B_i})$ is a decreasing net
in $\PSH_{\mathrm{hom},X}(\theta)$, and if
$\inf_i\chi_{B_i}$ is not identically $-\infty$, which holds
as soon as $\inf_i\chi_{B_i}(v_{\triv})=0$, the limit lies in
$\PSH_{\mathrm{hom},X}(\theta)$.
\end{corollary}

\begin{proof}
Each $\chi_{B_i}\in\PSH_{\mathrm{hom},X}(\theta_i)$ by
\Cref{thm:bj-real-cartier-realization}, hence in
$\PSH_{\mathrm{hom},X}(\theta)$ by
\Cref{prop:ch38-drift-ledger}\,\cref{itm:ch38-drift:dominated} applied to
the constant net, that is, by
\Cref{lem:ch34-psh-nonempty} and \Cref{thm:bj-basic-closure}\,(3).  The net
is decreasing on $X^{\divis}$ by
\Cref{prop:ch38-evaluation-welldefined}\,\cref{itm:ch38-ev:linear} and
hence on $X^{\An}$ by
\Cref{cor:ch35-divisorial-determination}\,(1).  Each $\chi_{B_i}$
vanishes at $v_{\triv}$ by
\Cref{prop:ch35-homogeneous-normalization}, so the stated condition holds
automatically and \Cref{thm:bj-basic-closure}\,(4) applies.
\end{proof}

\begin{corollary}
\label{cor:ch38-cartier-psh-criterion-unconditional}
Let $X$ be normal integral projective, $\theta\in N^1(X)$, and let
$B$ be a \emph{real} Cartier b-divisor with determination
$\pi:Y\to X$.  Then
\[
 \begin{gathered}
 \psi_B\ \text{is the divisorial trace of an element}\\
 \text{of }\PSH_{\mathrm{hom},X}(\theta)
 \end{gathered}
 \quad\Longleftrightarrow\quad
 \begin{gathered}
 B\leq0,\\
 \pi^*\theta+B_Y\ \text{is nef}.
 \end{gathered}
\tag{38.46}\label{eq:ch38-unconditional-criterion}
\]
and the element is then unique.  No envelope hypothesis, no smoothness
hypothesis and no hypothesis on $\operatorname{char}k$ is required.  In
particular \Cref{thm:ch35-cartier-psh-criterion} holds unconditionally.
\end{corollary}

\begin{proof}
$\Leftarrow$ is \Cref{thm:bj-real-cartier-realization}, together with the
uniqueness proved there.

$\Rightarrow$ is the unconditional half of
\Cref{thm:ch35-cartier-psh-criterion}, proved in
\Cref{chap:bj-homogeneity-bdivisors} without any envelope hypothesis, and
its proof applies verbatim to real coefficients: if
$\chi\in\PSH_{\mathrm{hom},X}(\theta)$ restricts to $\psi_B$, then
$\psi_B\leq0$ by \Cref{prop:ch35-homogeneous-normalization}, hence
$B\leq0$ by
\Cref{prop:ch38-evaluation-welldefined}\,\cref{itm:ch38-ev:linear}; and
\Cref{cor:ch35-homogenization-psh} produces a decreasing net
$\psi_i\in\cH_{\mathrm{hom}}(L_i)$ with $c_1(L_i)\to\theta$ and
$\psi_i\searrow\chi$, whose coefficient b-divisors $B_i$ satisfy
$B_i\leq0$ with $c_1(L_i)+B_i$ nef, by
\Cref{prop:ch38-model-fs-dictionary} and
\Cref{thm:ch35-semiample-criterion}; the net $(B_i)$ is decreasing and
converges to $B$ pointwise on $X^{\divis}$, and
\Cref{lem:ch35-nef-limit} gives that $\theta+B$ is nef, that is, that
$\pi^*\theta+B_Y$ is nef.  Note that
\Cref{lem:ch35-nef-limit} is stated for rational $B_i$ and an arbitrary
Cartier limit $B$, which is what is used.
\end{proof}

\begin{warning}{The corollary does not remove the envelope property from
\Cref{chap:bj-homogeneity-bdivisors}}
\label{warn:ch38-criterion-scope}
\Cref{cor:ch38-cartier-psh-criterion-unconditional} supersedes
\Cref{thm:ch35-cartier-psh-criterion} logically, but it may not be
substituted for it inside
\Cref{chap:bj-homogeneity-bdivisors}: the numerical b-class theorem
\Cref{thm:bj-nef-bdivisor-correspondence} also uses the compactness
statement \cite[Corollary~6.22]{BJ22}, whose hypothesis is the envelope
property for ample classes and which is not addressed here.  What the
corollary does supply is the Cartier-level criterion in the exact form in
which \Cref{chap:bj-reflection} and
\Cref{chap:bj-envelope-consequences} use it, with real coefficients and
without hypotheses.  The order in which the results are proved --- first
\Cref{thm:bj-real-cartier-realization}, then the criterion --- is what
keeps \Cref{warn:ch38-firewall} intact.
\end{warning}

\subsection{The approximation package}
\label{sec:ch38-approximation-package}

The next corollary is the form in which
\Crefrange{chap:bj-reflection}{chap:bj-subintegral-normalization} consume
the homogeneous theory.  It says that every homogeneous psh function is
approximated from above by \emph{model} data --- rational antieffective
Cartier b-divisors with semiample total class --- in a decreasing way, and
that the b-divisors themselves increase in absolute value.

\begin{corollary}
\label{cor:ch38-gf-regularization-bdivisor}
Let $X$ be normal integral projective, $\theta\in N^1(X)$, and
$\psi\in\PSH_{\mathrm{hom},X}(\theta)$.  Then there are a directed set
$A$, rational Cartier b-divisors $C_a$ over $X$, and
$\mathbf Q$-line bundles $L_a$ on $X$, such that for every
$a\in A$:
\begin{enumerate}
\item\label{itm:ch38-pkg:neg} $C_a\leq0$, and on a determination
 $\pi_a:Y_a\to X$ of $C_a$ the line bundle
 $\pi_a^*L_a+C_{a,Y_a}$ is semiample;
\item\label{itm:ch38-pkg:fs} $\psi_{C_a}\in\cH_{\mathrm{hom}}(L_a)$, so
 that $\psi_{C_a}=m_a^{-1}\log|\mathfrak b_a|$ for some
 $m_a\geq1$ and some nonzero ideal
 $\mathfrak b_a\subseteq\cO_X$ with
 $m_aL_a\otimes\mathfrak b_a$ globally generated;
\item\label{itm:ch38-pkg:mono} the net $(C_a)$ is decreasing, that is,
 the effective divisors $D_a:=-C_a$ increase;
\item\label{itm:ch38-pkg:limit} $\psi_{C_a}\searrow\psi$ pointwise on
 $X^{\An}$;
\item\label{itm:ch38-pkg:class} $c_1(L_a)\to\theta$ in $N^1(X)$.
\end{enumerate}
If $\theta=c_1(L)$ for an ample $\mathbf Q$-line bundle $L$, one may
take $L_a=L$ for every $a$, so that \cref{itm:ch38-pkg:class} becomes
an equality and no class drift occurs.
\end{corollary}

\begin{proof}
By \Cref{prop:ch35-homogeneous-normalization},
\[
 \sup_{X^{\An}}\psi=\psi(v_{\triv})=0.
\]
Thus $\psi$ satisfies the hypothesis $\sup\psi\geq0$ of
\Cref{cor:ch35-homogenization-psh} and is its own homogenization, by
\Cref{lem:ch35-homogenization-maximal}.  That corollary provides a directed
set $A$, $\mathbf Q$-line bundles $L_a$ with
$c_1(L_a)\to\theta$, and a decreasing net
$\psi_a\in\cH_{\mathrm{hom}}(L_a)$ with $\psi_a\searrow\psi$
pointwise; and in the ample rational case it provides $L_a=L$.  This is
\cref{itm:ch38-pkg:limit} and \cref{itm:ch38-pkg:class}.

By \Cref{prop:ch38-model-fs-dictionary}\,\cref{itm:ch38-dict:backward} and
\cref{itm:ch38-dict:unique} there is, for each $a$, a unique rational
Cartier b-divisor $C_a$ with $\psi_{C_a}=\psi_a$; by
\Cref{thm:ch35-semiample-criterion} it satisfies $C_a\leq0$ with
semiample total class, which is \cref{itm:ch38-pkg:neg}, and
\cref{itm:ch38-pkg:fs} is
\Cref{prop:hom-fs-ideal}.  Finally
\cref{itm:ch38-pkg:mono} follows from
\cref{itm:ch38-pkg:limit} and
\Cref{prop:ch38-evaluation-welldefined}\,\cref{itm:ch38-ev:linear}: the net
$(\psi_{C_a})$ is decreasing on $X^{\divis}$, and the evaluation
pairing is an order isomorphism onto its image.
\end{proof}

\begin{remark}[The shape in which the later chapters use it]
\label{rem:ch38-package-use}
\Cref{chap:bj-reflection} applies
\Cref{cor:ch38-gf-regularization-bdivisor} on the source of a birational
morphism, writes $-D_i$ for the coefficient b-divisor of the $i$-th
stage and $\alpha_i=c_1(L_i)$ for its class, and then corrects the
classes by the exact identity \eqref{eq:bj-reflection-total} before
applying \Cref{thm:bj-real-cartier-realization} downstairs; the three
displayed properties \eqref{eq:bj-homogeneous-regularization-data} there
are \cref{itm:ch38-pkg:mono}, \cref{itm:ch38-pkg:class} and the
semiampleness in \cref{itm:ch38-pkg:neg}.
\Cref{chap:bj-smooth-characteristic} uses the presentation
\cref{itm:ch38-pkg:fs} with the ideals $\mathfrak b_a$ replaced by
asymptotic multiplier or test ideals, which is why the corollary is stated
with the ideal displayed rather than only with the b-divisor.
\end{remark}

\subsection{Behaviour under a birational base change}
\label{sec:ch38-base-change}

\begin{lemma}
\label{lem:ch38-pullback-stages}
Let $\beta:X'\to X$ be a projective birational morphism of normal
integral projective varieties and let $\beta^{\An}:X'^{\An}\to X^{\An}$
be the induced map, which is a continuous surjection restricting to a
bijection $X'^{\val}\to X^{\val}$ and to the identity on
$X'^{\divis}=X^{\divis}$.  Write $\beta^*u:=u\circ\beta^{\An}$.
\begin{enumerate}
\item\label{itm:ch38-pb:fs} If $L\in\Pic(X)_{\mathbf Q}$ and
 $\psi\in\cH_{\mathrm{hom}}(L)$, then
 $\beta^*\psi\in\cH_{\mathrm{hom}}(\beta^*L)$.  Explicitly, if
 $\psi=m^{-1}\log|\mathfrak b|$ then
 $\beta^*\psi=m^{-1}\log|\mathfrak b\cdot\cO_{X'}|$.
\item\label{itm:ch38-pb:limit} Let $\omega\in N^1(X)$ be nef, let
 $(\psi_n)$ be a decreasing sequence with
 $\psi_n\in\cH_{\mathrm{hom}}(L_n)$ and $\omega-c_1(L_n)$ nef for
 every $n$, and let $\chi:=\lim_n\psi_n$.  If
 $\chi(v_{\triv})=0$, then
 $\beta^*\chi\in\PSH_{\mathrm{hom},X'}(\beta^*\omega)$.
\end{enumerate}
\end{lemma}

\begin{proof}
The stated properties of $\beta^{\An}$ are standard: continuity and
surjectivity because $\beta$ is proper and surjective; the bijection on
full-support loci because both sides are the set of real valuations of the
common function field $k(X)=k(X')$ admitting a centre, as in the proof of
\Cref{prop:ch35-cartier-continuous-extension}; and the identity on
divisorial points because a prime divisor over $X'$ is a prime divisor
over $X$, by \Cref{prop:ch38-evaluation-functorial}.

\emph{\cref{itm:ch38-pb:fs}.}  If $0\neq s\in H^0(X,mL)$ then
$\beta^*s\in H^0(X',m\beta^*L)$ is nonzero, $\beta$ being dominant and
$X'$ integral, and
$\log|\beta^*s|(w)=-w(\beta^*s)=-\beta^{\An}(w)(s)=\log|s|(\beta^{\An}(w))$
by the definition of the pushforward semivaluation.  Taking a presentation
$\psi=m^{-1}\max_i\log|s_i|$ with all weights zero, as in
\eqref{eq:homogeneous-fs-presentation}, gives
$\beta^*\psi=m^{-1}\max_i\log|\beta^*s_i|$, a presentation with all
weights zero.  Generic finiteness is preserved because the $\beta^*s_i$
do not all vanish identically.  For the ideal formulation, the ideal
generated by the $\beta^*s_i$ in local frames is
$\mathfrak b\cdot\cO_{X'}$, and
$w(\mathfrak b\cdot\cO_{X'})=\beta^{\An}(w)(\mathfrak b)$.

\emph{\cref{itm:ch38-pb:limit}.}  By
\Cref{lem:ch34-psh-nonempty} and \Cref{thm:bj-basic-closure}\,(3),
$\psi_n\in\PSH_X(\omega)$ for every $n$.  By
\cref{itm:ch38-pb:fs}, $\beta^*\psi_n\in\cH_{\mathrm{hom}}(\beta^*L_n)
\subseteq\PSH_{X'}(\beta^*c_1(L_n))$.  Moreover
$\beta^*\omega-\beta^*c_1(L_n)=\beta^*(\omega-c_1(L_n))$ is nef, being
the pullback of a nef class
(\Cref{lem:ch38-numerical-birational}\,\cref{itm:ch38-num:inj}), so
\Cref{thm:bj-basic-closure}\,(3) again gives
$\beta^*\psi_n\in\PSH_{X'}(\beta^*\omega)$.  The sequence
$(\beta^*\psi_n)$ is decreasing, being the composition of a decreasing
sequence with a fixed map, and converges pointwise to $\beta^*\chi$.
Since $\beta^{\An}(v_{X',\triv})=v_{X,\triv}$, one has
$\beta^*\chi(v_{X',\triv})=\chi(v_{\triv})=0$, so the limit is not
identically $-\infty$ and \Cref{thm:bj-basic-closure}\,(4) gives
$\beta^*\chi\in\PSH_{X'}(\beta^*\omega)$.  Homogeneity is preserved
because $\beta^{\An}$ is equivariant for the scaling action,
$\beta^{\An}(tw)=t\,\beta^{\An}(w)$, which is immediate from the
definition of the pushforward semivaluation.
\end{proof}

\begin{corollary}
\label{cor:ch38-realization-pullback}
Let $\beta:X'\to X$ be a projective birational morphism of normal
integral projective varieties, let $\theta\in N^1(X)$, and let
$B\leq0$ be a real Cartier b-divisor over $X$ with a determination
$\pi:Y\to X$ factoring through $\beta$, such that
$\pi^*\theta+B_Y$ is nef.  Then $B$, regarded as a Cartier b-divisor
over $X'$ by
\Cref{prop:ch38-evaluation-functorial}\,\cref{itm:ch38-fun:cofinal}, is
antieffective and nef-total in the class $\beta^*\theta$; both
$\chi^X_B\in\PSH_{\mathrm{hom},X}(\theta)$ and
$\chi^{X'}_B\in\PSH_{\mathrm{hom},X'}(\beta^*\theta)$ exist, and
\[
 \chi^{X'}_B\leq\beta^*\chi^X_B\ \text{ on }X'^{\An},
 \qquad
 \chi^{X'}_B=\beta^*\chi^X_B\ \text{ on }X'^{\val}.
\tag{38.47}\label{eq:ch38-realization-pullback}
\]
\end{corollary}

\begin{proof}
Antieffectivity and nef-totality over $X'$ are
\Cref{prop:ch38-evaluation-functorial}\,\cref{itm:ch38-fun:class}, and the
two realizations exist by \Cref{thm:bj-real-cartier-realization}.  Both
restrict to $\psi_B$ on $X'^{\divis}=X^{\divis}$: the first by the
theorem, the second because
$\beta^*\chi^X_B$ restricted to $X'^{\divis}$ is
$\chi^X_B|_{X^{\divis}}=\psi_B$, by
\Cref{lem:ch38-pullback-stages}.  The function $\beta^*\chi^X_B$ is
upper semicontinuous, being the composition of an upper-semicontinuous
function with a continuous map, so
\Cref{thm:ch35-divisorial-comparison} applied to the psh function
$\chi^{X'}_B$ and the majorant $\beta^*\chi^X_B$ gives the inequality
in \eqref{eq:ch38-realization-pullback}.

For the equality on $X'^{\val}$: run the proof of
\Cref{thm:bj-real-cartier-realization} on $X$, producing
$\chi'=\lim_n\chi_n$ and $\chi^X_B$ with
$\chi'=\log|s_F|+\chi^X_B$.  By
\Cref{lem:ch38-pullback-stages}\,\cref{itm:ch38-pb:limit}, applied with
$\omega$ and the stages $\chi_n$, which satisfy
$\omega-\omega_n$ ample and hence nef,
$\beta^*\chi'\in\PSH_{\mathrm{hom},X'}(\beta^*\omega)$.  Pulling back the
Siu identity, which is an identity of sums of functions with values in
$[-\infty,0]$ and therefore commutes with composition,
\[
 \beta^*\chi'=\log|s_{\beta^*F}|+\beta^*\chi^X_B .
\]
On the other hand $\beta^*\chi'\leq\log|s_{\beta^*F}|$, so
\Cref{thm:ch35-homogeneous-siu-general} produces a unique
$u\in\PSH_{\mathrm{hom},X'}(\beta^*\omega-[\beta^*F])
=\PSH_{\mathrm{hom},X'}(\beta^*\theta)$ with
$\beta^*\chi'=\log|s_{\beta^*F}|+u$; and $u$ has divisorial trace
$\psi_B$, so $u=\chi^{X'}_B$ by
\Cref{cor:ch35-divisorial-determination}\,(2).  At a point
$w\in X'^{\val}$ the number $\log|s_{\beta^*F}|(w)$ is finite, because
$\beta^*F$ is a divisor on the integral variety $X'$ and therefore does
not contain the support of $w$; so the two displayed identities may be
compared and the finite quantity cancelled, giving
$\chi^{X'}_B(w)=u(w)=\beta^*\chi^X_B(w)$.
\end{proof}

\begin{remark}[The full pullback statement is proved later, and is not used
here]
\label{rem:ch38-pullback-forward}
The equality in \eqref{eq:ch38-realization-pullback} extends to all of
$X'^{\An}$, because pullback along $\beta$ is a bijection
$\PSH_{\mathrm{hom},X}(\theta)\to\PSH_{\mathrm{hom},X'}(\beta^*\theta)$;
that is \Cref{thm:bj-homogeneous-reflection}, whose proof uses the present
chapter.  Nothing below \Cref{chap:bj-reflection} uses the extension, and
\Cref{cor:ch38-realization-pullback} is stated in the form that is proved
here so that the dependency runs in one direction only.  The gap between
the two statements is exactly the gap discussed in
\Cref{warn:ch38-proper-support-value}: the identity
$\beta^*\chi'=\log|s_{\beta^*F}|+\beta^*\chi^X_B$ carries no information
at a point whose support lies in $\Supp\beta^*F$, and closing it requires
knowing that $\beta^*\chi^X_B$ is psh, which is the reflection theorem.
\end{remark}

\subsection{The limits of the method}
\label{sec:ch38-limits}

\begin{remark}[Cartier, not arbitrary Weil]
\label{rem:ch38-cartier-not-weil}
The theorem assumes that $B$ is Cartier on one model.  It does not realize
an arbitrary Weil b-divisor, and it does not assert that every homogeneous
psh function is determined on one fixed model.  The stated direction is
exactly the one used by birational reflection.

Three further points make the scope precise.  First, the restriction is not
an artefact of the proof: \Cref{ex:ch38-moving-support} shows that the
ordered rational approximation of \Cref{sec:ch38-approximation}
works in one fixed finite-dimensional space and cannot be applied
directly to that moving-support Weil b-divisor. This example
does not rule out other approximation methods.  Second, the Weil statement is
\emph{true} under stronger hypotheses --- it is
\Cref{thm:bj-nef-bdivisor-correspondence}, for $X$ smooth over a field of
characteristic zero --- and the proof there uses the envelope property for
ample classes, which is unavailable here.  Third, the Cartier statement is
what the envelope proof consumes: in
\Cref{chap:bj-reflection} the b-divisors are the coefficient b-divisors of
homogeneous Fubini--Study functions, hence rational and Cartier by
\Cref{prop:ch38-model-fs-dictionary}, and their later limit is treated through psh closure; no
assertion that an arbitrary homogeneous psh limit is Cartier is made.
\end{remark}

\begin{remark}[Nef limit and ample stages]
\label{rem:ch38-nef-limit-ample-stages}
The total class $P$ is merely nef.  The path
\eqref{eq:bj-buffer-path} replaces it by ample totals, which admit rational
semiample coefficient functions.  The class drift
$\omega_n\to\omega=\theta+[F]$ is handled by PSH2, and the fixed divisor
$F$ is removed by Siu division.

With the sharpened buffer of \Cref{lem:ch38-buffer} the real path carries
no drift at all: $\Omega_t=\omega$ for every $t$.  The drift enters only
at the rationalization step, and is then oriented downwards by
\Cref{lem:ch38-ordered-rational-dominated}, so that ``handled by PSH2''
means precisely: handled by \Cref{thm:bj-basic-closure}\,(3) and~(4) inside
the single space $\PSH_X(\omega)$.  See
\Cref{warn:ch38-drift-not-absorbed}.
\end{remark}

\paragraph{Proof checkpoint.}
Before the fixed-class decreasing-limit clause is applied, every
stage is placed in $\PSH_X(\omega)$ and the limit is checked to have
value $0$ at $v_{\triv}$. Class convergence would additionally be
required if the moving-class clause were used.  A
decreasing family alone could converge to the excluded function
$-\infty$.

\paragraph{Proof checkpoint.}
No upper envelope appears anywhere in the construction.  This independence
is what allows the theorem to be used in \Cref{chap:bj-reflection} to prove homogeneous
birational reflection before envelope property is descended.

\paragraph{Proof checkpoint.}
Three inputs to the proof are numerical and one is analytic.  Nefness of
$P$, ampleness of $\omega$, and relative ampleness of $-G$ are
numerical, and they combine in the single identity
\eqref{eq:bj-buffer-ample-total}.  The analytic input is the homogeneous
Siu division \Cref{thm:ch35-homogeneous-siu-general}, which is the only
borderline analytic input in the main proof that is not reproved in this
book, and which carries no envelope hypothesis.

\section{Chapter dictionary}
\label{sec:ch38-dictionary}

The interfaces established or used in this chapter are the following.
The first and third items are equivalences, the second records the canonical
extension attached to a Cartier b-divisor, the fourth separates two
gradings, and the fifth is the main theorem.  The last two items record
statements used only later and explicitly not proved here.

\begin{description}
\item[Weil coefficients.]
  A locally finite homogeneous function on $X^{\divis}$ corresponds to a
  Weil b-divisor $B$ when $X$ is normal
  (\Cref{prop:ch35-weil-coefficient-bijection}).
\item[Cartier evaluation.]
  A Cartier b-divisor $B$ on a normal $X$ has a coefficient function
  $\psi_B$ on $X^{\divis}$, with a unique continuous extension to
  $X^{\val}$
  (\Cref{prop:ch38-evaluation-domains}\,\cref{itm:ch38-dom:val}).
\item[Homogeneous Fubini--Study data.]
  A function $\psi\in\cH_{\mathrm{hom}}(L)$ presented by a single ideal as
  $m^{-1}\log|\mathfrak b|$ corresponds, when $X$ is normal, to a rational
  Cartier b-divisor $B\leq0$ for which $\pi^*L+B_Y$ is semiample
  (\Cref{prop:ch38-model-fs-dictionary}).
\item[The two gradings.]
  Trivial weight grading is homogeneity, whereas changing the degree grading
  only changes the Veronese presentation of the same function
  (\Cref{prop:ch38-two-gradings}).
\item[Real Cartier realization.]
  If $X$ is normal, integral, and projective, $B\leq0$ is a real Cartier
  b-divisor, and $\pi^*\theta+B_Y$ is nef, there is a unique
  $\chi\in\PSH_{\mathrm{hom},X}(\theta)$ with trace $\psi_B$; no envelope
  property is assumed
  (\Cref{thm:bj-real-cartier-realization} and
  \Cref{cor:ch38-cartier-psh-criterion-unconditional}).
\item[Nef b-divisors, later.]
  For $X$ smooth of characteristic zero and with the envelope property for
  ample classes, a nef b-divisor class $\alpha$ with $\alpha_X=\theta$
  corresponds to a function in $\PSH_{\mathrm{hom},X}(\theta)$ vanishing on
  the primes of $X$ (\Cref{thm:bj-nef-bdivisor-correspondence}).  This
  statement is not proved here.
\item[Birational reflection, later.]
  The pullback $\beta^*$ is a bijection on homogeneous psh classes, reflecting
  the identity of Cartier b-divisors over $X$ and over $X'$; this is
  \Cref{thm:bj-homogeneous-reflection}, not proved here.  The Cartier-level
  input supplied by the present chapter is
  \Cref{cor:ch38-realization-pullback}.
\end{description}

Four separations should be carried forward.

\emph{Divisorial trace versus function.}  A Cartier b-divisor gives a
function on $X^{\divis}$ by \eqref{eq:bj-bdivisor-evaluation}; this has the
unique continuous extension to $X^{\val}$ furnished by
\Cref{prop:ch38-evaluation-domains}\,\cref{itm:ch38-dom:val}.  For a rational antieffective b-divisor with relatively semiample
determination, an ideal gives a continuous function on $X^{\An}$.
Absolute semiampleness of the total line bundle additionally puts
this function in the corresponding homogeneous GF cone.  For a general antieffective $B$ with nef total
class the function on $X^{\An}$ exists but is produced by
\Cref{thm:bj-real-cartier-realization} and is uniquely characterized
among $\theta$-psh functions by its trace, never by a divisor formula
(\Cref{warn:ch38-three-extensions},
\Cref{warn:ch38-proper-support-value}).

\emph{Semiample versus nef.}  The dictionary of
\Cref{prop:ch38-model-fs-dictionary} needs semiampleness, a condition on
linear series; the hypothesis of
\Cref{thm:bj-real-cartier-realization} is nefness, a numerical condition.
The ample buffer of \Cref{lem:ch38-buffer} is the bridge, and by
\Cref{lem:ch38-numerical-birational}\,\cref{itm:ch38-num:notample} it must
involve a $\pi$-ample correction on $Y$; no positivity imported from
$X$ can do the job.

\emph{Size versus orientation of the drift.}  The class error at stage
$n$ is $O(n^{-1})$. In this fixed-class proof,
$\omega-\omega_n$ is ample, so that
\Cref{thm:bj-basic-closure}\,(3) places every stage in the single space
$\PSH_X(\omega)$.  \Cref{warn:ch38-drift-not-absorbed} states the
distinction and \Cref{ex:ch38-drift-irrational} shows that rational numerical classes cannot equal an
irrational target, even though their psh functions can be regarded
in a common larger class.

\emph{Degree versus weight.}  A weighted section system carries a degree
$m$ and weights $\lambda_j$.  Homogeneity is the vanishing of the
weights; the degree has a Veronese ambiguity. Twisting a degree-$m$
series by a fixed bundle changes its numerical class by $O(m^{-1})$;
a common bounded weight shift changes its value by $O(m^{-1})$
(\Cref{prop:ch38-two-gradings}, \Cref{warn:ch38-two-gradings},
\Cref{rem:ch38-two-errors}).

\section{Source note}

\Cref{lem:bj-ordered-rational-approximation,thm:bj-real-cartier-realization}
follow the Cartier realization section of
the supplied envelope manuscript.  The coefficient construction,
homogeneous comparison, regularization, and Siu division are
\cite[Theorem~6.32, Lemma~6.34, Theorem~4.22,
Corollaries~4.23 and~6.17, and Theorem~6.21]{BJ22}.  We retained the
fixed-support approximation, the
trivial-valuation test, and the global-identity safeguard at the pole
divisor.

The remaining attributions are as follows.  The b-divisor formalism, the
model category and the projection formula are those of
\Cref{app:intersection-bdivisors}, ultimately \cite{BFJ09}; the
BFJ citation here is only for birational b-divisor formalism and ordinary
intersection theory on projective models.  No absolute
$(n+1)$-fold pairing of discretely valued BFJ model forms is used in this
chapter.  The
identification of Cartier b-divisors with homogeneous PL functions and the
semiample criterion are \cite[\S6.4]{BJ22}, restated here as
\Cref{prop:ch38-model-fs-dictionary} in the form used at every rational
stage.  The domination theorem used to produce a relatively ample
exceptional divisor in \Cref{lem:ch38-activation} is
\cite[\href{https://stacks.math.columbia.edu/tag/081T}{Tag 081T}]{stacks-project},
imported through \Cref{thm:appB-admissible-domination}; the twist theorem
converting relative ampleness into ampleness,
\eqref{eq:ch38-relative-to-absolute}, is \cite[4.6.13]{EGAII}; and the
negativity lemma used in
\Cref{lem:ch38-nef-total-intrinsic}\,\cref{itm:ch38-nt:exceptional} is
\cite[Lemma~3.39]{KM08}.  Zariski's main theorem in the form
$\pi_*\cO_Y=\cO_X$ for a proper birational morphism onto a normal
variety, and the regularity criterion for a rational section with effective
divisor, are \cite[Proposition~II.6.3A and Corollary~III.11.4]{Har}; the
Nakai--Moishezon criterion used in the surface computations of
\Cref{ex:ch38-activation-surface} and
\Cref{ex:ch38-buffer-surface} is \cite[Theorem~1.2.23]{PAG1}.  The
valuative criterion of integral dependence behind
\Cref{prop:ch38-two-gradings}\,\cref{itm:ch38-grad:integral} is
\Cref{thm:appA-valuative-integral-closure}.

Four things in this chapter are not in the sources in this form and are
supplied here.  The activation data is packaged as a free-standing lemma,
\Cref{lem:ch38-activation}, with the two normalizations
$2G\leq\pi^*F$ and $\pi^*\omega-G$ ample, because
\Cref{chap:bj-reflection} quotes the construction and because the second
normalization is what removes the class drift from the real path.  The
buffer is likewise isolated, \Cref{lem:ch38-buffer}, and stated so that the
auxiliary class $H$ of \eqref{eq:bj-buffer-path} may be taken to be zero.
The drift ledger \Cref{prop:ch38-drift-ledger} and the dominated form
\Cref{lem:ch38-ordered-rational-dominated} of the approximation lemma make
explicit that the approximating classes may be chosen \emph{below} the
target class, which is what allows
\Cref{thm:bj-basic-closure}\,(3) and~(4) to be applied inside a single
class; the sources pass to the limit with a moving class.  And the
extended-real audit \Cref{prop:ch38-safeguards}, together with the failure
example \Cref{ex:ch38-naive-realization-fails}, is written out because the
plan's proof audit asks for every use of subtraction involving an infinite
value to be exhibited.

Two points of scope deserve a final word.  \Cref{ex:ch38-strict-room-fails}
shows that the homogeneous Siu division is applied at the boundary of the
strict regime, so the general form
\Cref{thm:ch35-homogeneous-siu-general}, that is
\cite[Theorem~6.21]{BJ22}, is the correct citation; the strict form
\Cref{thm:ch35-homogeneous-siu} suffices when
$B$ is bounded away from zero along every component of
$\Supp\pi^*F$.  And \Cref{warn:ch38-theta-not-nef} records that the
hypothesis of \Cref{thm:bj-real-cartier-realization} constrains
the nef total class on the determination. The surface trace
conclusion requires $B_X$ to be $\mathbf R$-Cartier, and the
warning gives a threefold example where even $\theta+[B_X]$
is not nef. Neither conclusion makes $\theta$ nef in general.
The historical route to these statements ran through nef b-divisor classes
\cite{BFJ09,DF20} and the divisorial Zariski decomposition
\cite{Bou04,Nak04}; the transcendental counterpart over $\mathbf C$ is
\cite{Xiabdiv,Xiabdiv2}.  The point of the present chapter is that the
Cartier statement, unlike the Weil one, needs none of that machinery and
none of the envelope property, and is therefore available as an input to
\Crefrange{chap:bj-reflection}{chap:bj-subintegral-normalization}.
